
% Cross-references from body sections (e.g.\ Appendix~\ref{appendix} in
% Section~\ref{sec:policy}) point here as the appendix entry point.
% Sections number A through I under the hub's \appendix switch.

% Paragraph purpose: Map each appendix to the reader question it answers.
\paragraph{Appendix roadmap.} The appendices follow the order in which a reader would check the paper: the interface, then the data and accounting, then the statistics, then each result in detail. Appendices~\ref{appendix}, \ref{app:invariants}, and~\ref{app:runtime} define the operations, invariants, and runtime, and show that the runtime never chooses or repairs on the policy's behalf. Appendix~\ref{app:data-splits} shows that no evaluation task leaked into training, Appendix~\ref{app:receipts} traces each reported run to its receipt, and Appendix~\ref{app:statistics} states what an interval licenses. Appendix~\ref{app:cf-contract} specifies the measurement that Box~\ref{box:procedure} validates. The results appendices then hold the submission-order audit behind the enumeration-contract identity failure (Appendix~\ref{app:results-counterfactual}) and the ladder (Appendix~\ref{app:results-ladder}). They also hold the adapter and the policies scored beside it (Appendices~\ref{app:results-adapter}, \ref{app:results-panels}, \ref{app:results-cf-dagger}, and~\ref{app:results-cf-zeroshot}), and the interface survey with both external audits (Appendices~\ref{app:results-external} and~\ref{app:results-external-lfs}). Appendices~\ref{app:results-primary} and~\ref{app:results-contrasts}, with the other Results Detail appendices, show what the enumeration-contract learned results measure, and Appendix~\ref{app:design-contracts} compares the two contracts.

\section{Operation Schemas}
\label{appendix}

% Paragraph purpose: Say what this appendix holds and why a reader would look here.
This appendix gives the operation vocabulary, validity checks, and rejection kinds of each declared algorithm, for readers who need the exact interface behind the Typed Search Operation definition.

% Paragraph purpose: Introduce the typed operation vocabulary and the operand ownership rule that governs all schemas.
The schemas enforce one rule: every operand that determines exploration must come from the policy (Section~\ref{sec:policy}). Each declared algorithm defines a vocabulary of operation types $\tau$, and a policy emission $u_t = (\tau_t, \theta_t)$ pairs a type from that vocabulary with an operand vector $\theta_t$ carrying every exploration-determining argument. The runtime rejects any emission whose type falls outside the declared vocabulary or whose operand vector omits a required field, so it never supplies a missing operand itself.

% Paragraph purpose: Describe the BFS operation schema at the level pinned by the sources.
\paragraph{BFS.}
The primary operation is node expansion. Its operands name the source state (the current frontier head), the grounded successors selected for insertion, the action producing each successor, and the enqueue position for each. The runtime checks that the named source matches the FIFO frontier head, that each successor results from an applicable grounded action in the source state, that the successor is unvisited, and that each insertion targets the FIFO tail. A goal-test operation names the state to test. The BFS pilot corpus exposes the facts these checks use, namely exact grounded successor candidates and full visited membership at each step.
% Evidence: BFS pilot corpus is issue #111.
% TODO: The exact BFS operation-type enumeration and per-field schema (field
% names, types, serialization format) are not recorded in any audited
% document; pin them from the codebase before release.

% Paragraph purpose: Describe the BFWS operation schema at the level pinned by the sources.
\paragraph{BFWS (complete unpruned).}
Expansion operands carry the novelty value and goal-count priority of the source state alongside the fields above. Insertion operands carry each candidate's duplicate membership, novelty, and priority. The runtime checks that the selected source has the declared best priority, that duplicate membership is correctly reported, and that insertions follow the declared ordering. The BFWS corpus exposes the facts these checks use, namely candidate duplicate, novelty, priority, insertion, and enqueue facts, at every training observation.
% Evidence: BFWS corpus is issues #57--#58.
% TODO: The exact BFWS operation-type enumeration, per-field schema, and the
% precise novelty computation (tuple arity, width parameter) are not pinned
% in audited documents; record them before release.

% Paragraph purpose: Describe the additive best-first operation schema at the level pinned by the sources.
\paragraph{Additive best-first.}
Operands carry the source state's accumulated path cost $g$, heuristic estimate $h_{\rm add}$, and composite value $f = g + w \cdot h_{\rm add}$ with $w = 3$ for the weighted variant, or $f = h_{\rm add}$ for the greedy variant. Insertion operands carry per-successor $g$, $h_{\rm add}$, and $f$. The runtime checks that the expanded state has the minimum $f$ in the frontier under the declared tie-breaking rule and that all arithmetic is consistent.
% TODO(schema): Per-field additive operation serialization remains to be pinned; the tie-break is pinned below from the executed controller.

% Paragraph purpose: State the validity preconditions as an ordered rejection hierarchy.
\paragraph{Validity preconditions.}
Every emission passes through the same ordered sequence of checks, and the first check it fails rejects it with that kind and ends the episode. Each recorded rejection kind therefore names the earliest check an emission failed. The checks run in the following order.

\begin{enumerate}[leftmargin=*, itemsep=1pt, parsep=1pt]
\item \textbf{Parse.} The raw output must parse into a well-formed Typed Search Operation with a recognized type and complete operand vector. A failure is recorded as \texttt{malformed\_output}.
\item \textbf{Source identity.} The named source state must match the state the algorithm expects. A failure is recorded as \texttt{source\_state\_mismatch}.
\item \textbf{Action applicability.} Every grounded action in the operands must be applicable in the source state. A failure is recorded as \texttt{inapplicable\_grounded\_action}.
\item \textbf{Successor schema.} Any named successor state must satisfy the domain schema. A failure is recorded as \texttt{schema}.
\item \textbf{Successor effects.} The named successor must match the declared effects of the applied action. A failure is recorded as \texttt{effect}.
\item \textbf{Static context.} Referenced predicates, objects, and constants must be consistent with the task instance. A failure is recorded as \texttt{static\_context}.
\item \textbf{Frontier discipline.} The frontier manipulation must respect the declared ordering (FIFO for BFS, priority for BFWS, $f$-order for additive best-first). A failure is recorded as \texttt{invalid\_frontier\_operation}.
\item \textbf{Algorithm invariant.} Any remaining algorithm-specific constraint not covered above. A failure is recorded as \texttt{other\_invariant\_violation}.
\end{enumerate}

% Paragraph purpose: Separate format failures from search failures, and give the random-valid control's single failure kind.
Rejections fall into two categories, format failures and search failures. \emph{Malformed-output} rejections never reach the operation checker, whereas \emph{post-parse rejections} parse but fail a runtime check. The oracle-assisted random-valid control emits only runtime-accepted operations, so its single failure kind is \texttt{expansion\_budget\_exhausted}, recorded when the frozen episode budget runs out before a goal is reached.

% Paragraph purpose: State what the rejection taxonomy shows and point to its table.
\paragraph{Rejection counts.}
The base model fails mostly on format, while learned policies fail mostly after parsing, on algorithm invariants (Table~\ref{tab:appendix-rejection}). The counts come from the failure-mechanism calibration, which classified all 905 failures across 2,910 replay-verified episodes with zero unclassified. The \texttt{effect}, \texttt{schema}, and \texttt{static\_context} kinds arise in successor-prediction episodes, where the model generates a dynamic-state successor.
% Evidence: failure-mechanism calibration is issue #125 (failure-mechanism-analysis.md).
% Evidence (Table tab:appendix-rejection): docs/experiments/expanded-study/failure-mechanism-analysis.md rejection-taxonomy table
% (arm class | kind | count | share rows; base 447, learned 407 of 906, oracle 51; 2,910 episodes, 905 failures, 0 unclassified).

\begin{table}[t]
\caption{\textbf{Base failures are mostly malformed output, and learned failures are mostly post-parse invariant violations.} Arm labels give failures/episodes where both counts exist; shares use failures within each class.}
\label{tab:appendix-rejection}
\centering\footnotesize
\setlength{\tabcolsep}{4pt}
\begin{tabular}{llrr}
\toprule[0.4mm]
\rowcolor{mygray}\textbf{Arm class} & \textbf{Rejection kind} & \textbf{Failures} & \textbf{Share} \\
\midrule
\multirow{3}{*}{Base (447/447)} & \texttt{malformed\_output} & 379 & 84.8\% \\
 & \texttt{invalid\_frontier\_operation} & 67 & 15.0\% \\
 & \texttt{other\_invariant\_violation} & 1 & 0.2\% \\
\midrule
\multirow{8}{*}{Learned (407/906)} & \texttt{other\_invariant\_violation} & 229 & 56.3\% \\
 & \texttt{malformed\_output} & 72 & 17.7\% \\
 & \texttt{inapplicable\_grounded\_action} & 56 & 13.8\% \\
 & \texttt{effect} & 27 & 6.6\% \\
 & \texttt{schema} & 16 & 3.9\% \\
 & \texttt{invalid\_frontier\_operation} & 5 & 1.2\% \\
 & \texttt{source\_state\_mismatch} & 1 & 0.2\% \\
 & \texttt{static\_context} & 1 & 0.2\% \\
\midrule
Oracle (51 failures) & \texttt{expansion\_budget\_exhausted} & 51 & 100.0\% \\
\bottomrule
\end{tabular}
\end{table}


\section{Algorithm Invariant Definitions}
\label{app:invariants}

% Paragraph purpose: Say what this appendix holds and why a reader would look here.
This appendix states the invariant the runtime enforces for each declared algorithm, which Section~\ref{sec:policy} defines only in outline.

% Paragraph purpose: State that invariants are per-operation checks and therefore say nothing about episode outcomes.
The Trusted Search Runtime checks the following invariants on every operation. Because each invariant constrains a single step, stepwise validity and invariant compliance do not establish termination, completeness, optimality, or episode success (Section~\ref{sec:policy}).

% Paragraph purpose: BFS invariant.
\paragraph{BFS.}
BFS keeps its frontier as a FIFO queue, so every expansion must select the queue head as its source and every insertion must target the queue tail. A state already in the visited set must not be re-inserted. This invariant is sufficient for completeness in finite state spaces with unit-cost transitions.

% Paragraph purpose: BFWS invariant.
\paragraph{BFWS (complete unpruned).}
The frontier is a priority queue ordered first by novelty rank (lower is better), then by goal count among states with equal novelty. No generated successor is pruned. A successor whose state already appears in the frontier or visited set is not re-inserted.
% TODO: The exact novelty definition (tuple arity, width parameter, and
% tie-breaking within equal novelty and equal goal count) is not pinned in
% any audited document. Record the precise rule before release.

% Paragraph purpose: Weighted additive best-first invariant.
\paragraph{Additive best-first, weighted ($w = 3$).}
The frontier is a priority queue ordered by $f(s) = g(s) + 3 \cdot h_{\rm add}(s)$, where $g(s)$ is accumulated path cost and $h_{\rm add}(s)$ is the additive heuristic. Among equal-$f$ states the lower generation serial is expanded first. The expanded state is the heap head. A successor with lower $g$ than an existing frontier entry for the same state reopens that entry with improved cost.
% Evidence (executed controller): examples/planning_benchmark_slice/best_first_controller.py BestFirstController.start_expansion and apply_operation, with PDDLStateAuthority.applicable_actions in pddl_state.py sorting by GroundedAction.serialize.

% Paragraph purpose: Greedy additive best-first invariant.
\paragraph{Additive best-first, greedy.}
The frontier is ordered by $h_{\rm add}(s)$ alone, with no path-cost component, then by ascending generation serial among ties. The expanded state is the heap head. Reopen and duplicate rules follow the weighted variant.
% Paragraph purpose: Show that the executed controller's serial assignment produces the enumeration-contract identity, and that the identity is specific to this runtime.
The executed controller's serial assignment is what makes the enumeration contract insensitive to submission order. In both additive settings, \texttt{BestFirstController.start\_expansion} assigns serials before submissions, in the deterministic \texttt{PDDLStateAuthority.applicable\_actions} order (sorted by \texttt{GroundedAction.serialize}). The \texttt{BestFirstController.apply\_operation} heap orders entries by (priority, serial, state identifier, path cost), and \texttt{finish\_expansion} requires the full candidate set. Submission permutations therefore leave frontier order unchanged on the enumeration contract. A runtime assigning serials in submission order would not be submission-order invariant and would not show this identity, so the identity is a property of this runtime (Appendix~\ref{app:results-counterfactual}).
% Evidence: examples/planning_benchmark_slice/{pddl_state.py PDDLStateAuthority.applicable_actions,best_first_controller.py BestFirstController.start_expansion/apply_operation/finish_expansion}; outputs/native-arms/v1/identity-audit.json structural_basis.


\section{Trusted Search Runtime}
\label{app:runtime}

% Paragraph purpose: Say what this appendix holds and why a reader would look here.
This appendix gives the runtime's guarantees and its step and episode pseudocode, for readers checking that the runtime never chooses or repairs on the policy's behalf.

% Paragraph purpose: State the three guarantees of the runtime.
The runtime provides three guarantees that hold across all declared algorithms.

\begin{enumerate}[leftmargin=*, itemsep=1pt, parsep=1pt]
\item \textbf{No-repair property.} The runtime validates, applies, persists to Search Memory, or rejects each emission. It never chooses omitted operands, reorders candidates, silently repairs an invalid emission, or substitutes a default search decision.
\item \textbf{Reject-and-terminate.} On the first rejected emission the episode terminates with the rejection kind recorded. There is no retry path.
\item \textbf{Replay determinism.} Replaying the raw emissions $u_{1:T}$ against the same task instance and the same initial Search Memory reproduces the explored trace. Any step whose outcome depends on runtime repair would violate this property.
\end{enumerate}

% Paragraph purpose: Say how the pseudocode maps to the ordered checks and the budgeted episode loop.
The two procedures below make the guarantees concrete. \textsc{Step} applies the checks of Appendix~\ref{appendix} in their fixed order, and \textsc{Episode} wraps it in a budgeted loop that ends at the first rejection or at the goal.

\begin{quote}
\small
\begin{tabbing}
\hspace{1.5em}\=\hspace{1.5em}\=\hspace{1.5em}\=\kill
\textbf{procedure} \textsc{Step}($u_t$, $M_{t-1}$, $A$, $x$) \\
\> \textit{// Phase 1: Parse the raw emission} \\
\> \textbf{if} $u_t$ does not parse as $(\tau, \theta)$ with $\tau \in \mathrm{vocab}(A)$ \textbf{then} \\
\>\> \textbf{return} \textsc{Reject}(\texttt{malformed\_output}) \\[2pt]
\> \textit{// Phase 2: Validate source identity} \\
\> \textbf{if} $\theta.\mathrm{source} \neq \mathrm{expected}(A, M_{t-1})$ \textbf{then} \\
\>\> \textbf{return} \textsc{Reject}(\texttt{source\_state\_mismatch}) \\[2pt]
\> \textit{// Phase 3: Validate operands against the task and state} \\
\> \textbf{for each} grounded action $a$ named in $\theta$ \textbf{do} \\
\>\> \textbf{if} $a$ is not applicable in $\theta.\mathrm{source}$ under $x.D$ \textbf{then} \\
\>\>\> \textbf{return} \textsc{Reject}(\texttt{inapplicable\_grounded\_action}) \\
\> \textbf{for each} successor state $s'$ named in $\theta$ \textbf{do} \\
\>\> \textbf{if} $s'$ violates the domain schema \textbf{then} \\
\>\>\> \textbf{return} \textsc{Reject}(\texttt{schema}) \\
\>\> \textbf{if} $s' \neq \mathrm{effects}(a, \theta.\mathrm{source})$ \textbf{then} \\
\>\>\> \textbf{return} \textsc{Reject}(\texttt{effect}) \\
\> \textbf{if} any static-context reference in $\theta$ is inconsistent with $x$ \textbf{then} \\
\>\> \textbf{return} \textsc{Reject}(\texttt{static\_context}) \\[2pt]
\> \textit{// Phase 4: Check Algorithm Invariant} \\
\> \textbf{if} frontier discipline of $A$ is violated by $(\tau, \theta, M_{t-1})$ \textbf{then} \\
\>\> \textbf{return} \textsc{Reject}(\texttt{invalid\_frontier\_operation}) \\
\> \textbf{if} any other invariant of $A$ is violated \textbf{then} \\
\>\> \textbf{return} \textsc{Reject}(\texttt{other\_invariant\_violation}) \\[2pt]
\> \textit{// Phase 5: Apply and persist} \\
\> $M_t \leftarrow \mathrm{apply}(\tau, \theta, M_{t-1})$ \\
\> \textbf{return} \textsc{Accept}($M_t$)
\end{tabbing}
\end{quote}

\begin{quote}
\small
\begin{tabbing}
\hspace{1.5em}\=\hspace{1.5em}\=\hspace{1.5em}\=\kill
\textbf{procedure} \textsc{Episode}($x$, $A$, $\phi_m$, $B$) \\
\> $M_0 \leftarrow \mathrm{InitSearchMemory}(A, x)$ \\
\> \textbf{for} $t = 1, 2, \ldots, B$ \textbf{do} \\
\>\> $o_t \leftarrow \phi_m\bigl(c_t(x, M_{t-1})\bigr)$ \hspace{2em} \textit{// build observation} \\
\>\> $u_t \leftarrow \pi(o_t)$ \hspace{6.4em} \textit{// policy emission} \\
\>\> $r \leftarrow$ \textsc{Step}($u_t$, $M_{t-1}$, $A$, $x$) \\
\>\> \textbf{if} $r$ is \textsc{Reject}(\textit{kind}) \textbf{then} \\
\>\>\> \textbf{return} $(t, \textsc{Fail}, \textit{kind})$ \\
\>\> $M_t \leftarrow r.M$ \\
\>\> \textbf{if} $\mathrm{GoalTest}(M_t, x.G) = \textsc{True}$ \textbf{then} \\
\>\>\> \textbf{return} $(t, \textsc{Success})$ \\
\> \textbf{return} $(B, \textsc{Fail}, \texttt{expansion\_budget\_exhausted})$
\end{tabbing}
\end{quote}

% Paragraph purpose: Note budget enforcement: rejected emissions are charged.
Each call to \textsc{Step} consumes one unit of the frozen episode budget $B$ whether the emission is accepted or rejected, so invalid operations are charged to the budget (Section~\ref{sec:design}).


\section{Data and Split Procedures}
\label{app:data-splits}

% Paragraph purpose: Say what this appendix holds and why a reader would look here.
This appendix describes the corpora, splits, and evaluation manifests behind the reported panels, for readers checking that no evaluation task leaked into training.

% Paragraph purpose: State the whole-instance split rule and the audit contract.
Every training and evaluation split is by whole task instance, so no state, decision, or target from an evaluation task appears in any training row. Every released corpus also passed no-overlap, no-conflict, live/training-parity, and token-budget audits before use.

% Paragraph purpose: Show that the BFS evaluation panel draws only from the dev side of a task-level split, then give the corpus receipts.
\paragraph{BFS pilot corpus.}
The BFS pilot corpus is split by semantic task identity, and its v8 evaluation panel draws only from the development side, so no panel task contributes training rows. The v6 text-state process corpus holds 25,109 process rows split into 12,994 train and 12,115 dev rows, drawn from 90 replay-verified trusted traces covering 90 distinct tasks, with zero cross-split task overlap. The v8 evaluation panel selects its 15 tasks from the development set. Corpus observations expose the canonical task context, exact grounded successor candidates, and full visited membership. Release audits recorded zero semantic, input, and input-target overlap, zero conflicting identical-input labels, zero teacher-operation rejections, and zero live/corpus mismatches, all within an 8,192-token input cap and a 384-token output allowance.
% Evidence: data/bfs_pilot_v6/qualification-attempt-004/qualification-report.json
% (90 tasks, zero semantic split overlap), process-release/manifests/bfs-text-corpus.json
% (split_unit=semantic_task_identity), process-release/audits/leakage.json,
% ms-swift-process/manifest.json, and docs/issue-54-operator-runbook.md L161--168.

% Paragraph purpose: Show that the BFWS corpus covers the development panel and passed its audits, then give its sizes.
\paragraph{BFWS corpus.}
The BFWS corpus covers the frozen 105-task development panel, and its audit counters recorded zero overlap, mismatch, rejection, or final-evaluation-access violations. It holds 69,019 replay-verified teacher decisions. Its release contains 69,019 process records, 67,215 operational records, and a process projection of 47,780 train and 21,239 dev rows. Each step's observation exposes candidate duplicate membership, novelty values, goal-count priorities, insertion positions, and enqueue facts.
% Evidence: docs/issue-58-bfws-text-corpus.md L39-46 (69,019 process records, 67,215 operational records, 47,780 train / 21,239 dev);
% docs/issue-58-bfws-text-corpus.md L42 (105 semantic-task split assignments); docs/issue-57-bfws-expert-traces.md L6 and docs/issue-56-bfws-development-freeze.md L19 (69,019 teacher decisions).

% Paragraph purpose: Give the enumeration-contract training recipes and state that they differ in size, so they are not a matched comparison.
\paragraph{Enumeration training recipe.}
The enumeration-contract recipes differ in size across interventions, so they describe different development-stage interventions rather than a matched-size comparison. Each expanded-baseline process-SFT cell trains a fresh LoRA adapter for one epoch on 512 replay-verified records, whereas the separate BFWS gate checkpoint trains for two epochs on 47,780 records. The pinned Qwen3-VL-8B-Instruct backbone revision is \texttt{0c351dd0}, and the program cap is 336 GPU-hours (Table~\ref{tab:appendix-receipts}).
% Evidence: docs/issue-59-bfws-structural-gate.md L10-13; outputs/expanded-study/v1/seed-replication/analysis.json and synthesis-v1/README.md; docs/experiments/expanded-study/goal10-completion-audit.json (backbone revision); docs/experiments/expanded-study/issue-127-closeout.md (cap).

% Paragraph purpose: State that every evaluation here is development-stage and that the frozen 45-task final evaluation was never accessed, then list the manifests.
\paragraph{Evaluation manifests.}
Every evaluation below is development-stage, and the one manifest reserved for final evaluation has never been accessed. The BFS pilot uses a frozen 15-task outcome-blind development panel. The BFWS structural gate froze a 105-task development panel and, separately, the manifest of the frozen 45-task final evaluation, which no model prediction, training decision, or analysis references. The expanded baseline evaluates 24 tasks from 12 domains under 4 algorithms, 3 modalities, and 4 arms (1,152 logical bindings, 576 model episodes, all independently replayed). The curriculum stage evaluates 27 unseen tasks. The generalization-v2 protocol admits 25 derived tasks from 93 eligible semantics-preserving variants (uniform $k{=}5$ cheapest-prefix per family, membership frozen before execution).
% TODO(release): Release URLs for published corpora and evidence bundles are
% pending handoff (issue #122). Insert URLs on release.
% TODO(final-eval): Held-out final primary evaluation unexecuted; no final
% efficacy finding exists. The pre-registration prerequisites in
% Section~\ref{sec:design} are undecided. No open ticket, so prose only.
% Panel-overlap search (2026-09-23): outputs/ and docs/experiments/ hold no pinned post-#126
% retention manifest; the BFS-v8 / BFWS-fallback overlap count stays a TODO in Section~\ref{sec:design}.


\section{Receipt Tables}
\label{app:receipts}

% Paragraph purpose: Say what this appendix holds and why a reader would look here.
This appendix holds the gate receipts, the per-window ledgers, and the artifact index, so that a reader can trace each reported run to its receipt.

% Paragraph purpose: State that every branch reconciles with zero missing evidence, and how the branches ended.
Every branch of the program reconciles with its ledger and replay, with zero missing evidence (Table~\ref{tab:appendix-receipts}). Five branches completed at v1 scope, while two stopped at frozen v1 cost-admission gates with terminal \textsc{valid\_stop} receipts and later completed at reduced v2 scope. The table combines each branch's frozen scope, gate status, coverage, and compute from the synthesis and follow-up ledgers, and its Tickets column indexes internal ticket numbers.
% Evidence: synthesis ledger is issue #120 (synthesis-v1/README.md).

\begin{table}[t]
\caption{\textbf{Every branch reconciles with zero missing evidence, and each follow-up window keeps its own GPU-hour ledger.} Scope gives bindings or episodes and independently replayed counts. Follow-up windows are not summed into the program total. Gate codes are explained below.}
\label{tab:appendix-receipts}
\centering\footnotesize
\setlength{\tabcolsep}{2pt}
\begin{tabular}{@{}p{0.18\linewidth}p{0.13\linewidth}p{0.12\linewidth}p{0.37\linewidth}r@{}}
\toprule[0.4mm]
\rowcolor{mygray}\textbf{Branch} & \textbf{Tickets} & \textbf{Gate} & \textbf{Scope / replay} & \textbf{GPU-h / cap} \\
\midrule
Expanded baseline & \#116--\#118 & $C$ & 1,152 bindings; 576 model; 1,152/1,152 & 6.53 / 56 \\
DAgger & \#78--\#84 & $C$ & 405 episodes; 81 paired; 405/405 & 10.91 / 48 \\
Successors & \#85--\#89, \#99 & $C$ & 90 episodes; 15 cells; 90/90 & 10.85 / 64 \\
Curriculum & \#119 & $C$ & 243 model; 324 controls; 567/567 & 11.78 / 48 \\
Transfer & \#104--\#107 & $C$ & 39 cells (12 v1, 27 v2); 39/39 & 3.60 / 24 \\
Backbone v2 & \#101--\#103, \#123 & $V\!\to\!C$ & 72 model; 72 controls; 144/144 & 6.03 / 52.11 \\
Generalization & \#96, \#98, \#124 & $V\!\to\!C$ & 1,200 (300 model, 900 controls); 1,200/1,200 & 6.04 / 32 \\
Corruption & \#126 & $C$ & 360/360 new & shared \\
Seed replication & \#127 & $C$ & 534/534 new & 12.79 / --- \\
Recovery reserve & --- & $R$ & --- & 0.00 / 11.89 \\
\midrule
\textbf{Program total} & & & & \textbf{68.52 / 336} \\
\midrule
\multicolumn{5}{@{}l}{\textit{Follow-up windows (independent ledgers)}} \\
Native arms & \#128, \#129 & $H\!\to\!C$ & 216/216 new; 72/72 controls & 3.2086 / 12 \\
Stress R1/R2/R4 & \#130 & $C$ & 72 menu; 80 variants; 18 cross-cell; 170/170 & 7.7998 / 36 \\
Backbone R3 & \#130 & $C$ & 72/72 model; 72/72 controls & shared \\
Choice frontier & \#132 & $Q_s$ & 144/144 panel; 58 preparations & 1.069 / 12 \\
Comparator zoo & \#133 & $C$ & 432/432 bindings; 342 new & CPU \\
Counterfactual / M4 & \#134 & \xmark & 48 pairs; 288 counts; 96 sequences & CPU \\
Validation & \#135 & $P$ & 696/696 CPU; 48/48 model & 1.19 / 4 \\
Scaled adapter & \#136 & \cmark & 24/24 learned & 4.84 / 30 \\
External audit & \#137 & $Q_r$ & 1,000 pairs; 200 tasks & CPU \\
Replication / panels & \#138, \#139 & \cmark & 48/48 new; zoo 638/638 P2, 696/696 P2u; learned 88/88 P2, 96/96 P2u & 15.22 / 24 \\
DAgger ablation & \#140 & $=$ & 24/24 validation; 22/22 P2 & 9.06 / 40 \\
Zero-shot & \#141 & $\equiv$ / $?$ & 46/46 model & 2.25 / 30 \\
LFS audit & \#142 & $\blacksquare$ & 80/80 reference; 400/400 random-valid & 2.89 / 12 \\
\bottomrule
\end{tabular}
% Evidence: all scope and ledger figures are unchanged from the pinned synthesis and follow-up sources in the comments below.
\end{table}
% Paragraph purpose: Decode the gate codes, which differ because branches ended in different ways, and note the GPU-hour special cases.
The Gate column uses one code per kind of outcome, because the branches ended in different ways. $C$ is completion at declared scope, $V\!\to\!C$ a frozen v1 \textsc{valid\_stop} followed by reduced-scope v2 completion, $H\!\to\!C$ the native arms' v1 gate history followed by v2 completion, and $R$ the reduced recovery reserve. $Q_s$ denotes choice sensitivity under the choice-frontier contract, $Q_r$ registered choice in the external audit, $P$ the instrument-validation pass, \xmark\ a failed counterfactual prediction, and $\blacksquare$ saturation under the LLM-First Search paper budget. \cmark\ marks a positive gate, $=$ non-separation, $\equiv$ equivalence on the validation panel, and $?$ an inconclusive result on the held-out panel. In the GPU-hour column, the second-backbone cap reflects the prospective 12.11 GPU-h reserve transfer, and corruption and seed replication share the existing program caps. Stress-test R1--R4 spend covers both stress rows, and all follow-up windows have independent ledgers.
% Evidence: receipt statuses and accounting from the unchanged source comments immediately below.
% Evidence: synthesis-v1/README.md and docs/experiments/expanded-study/issue-127-closeout.md (program/addenda);
% docs/experiments/native-arms/issue-129-closeout.md (native), outputs/native-arms/v1/stress-evaluation.json (170),
% issue-130-closeout.md (R3 72 + 72), outputs/choice-frontier/v1/evaluation/final-audit.json (144),
% preparation/visual-choice-frontier/audit.json (58), outputs/choice-frontier/v1/budget.json (1.069 / 12),
% outputs/choice-frontier/o4/metrics/analysis.json episode_accounting (432, 342 new), #133 closing comment.
% Source conflict: outputs/expanded-study/v1/second-backbone-v3/evaluation/evidence.json lists comparator_episodes 144
% against expected 72; issue-130-closeout.md states 72/72 comparators replayed, used for the R3 row.
% Evidence: v4 window "Shared v4 ledger total: 15.22 / 24 GPU-h (#138 10.11 + #139 5.12)." issue-139-closeout.md L144;
% cap outputs/choice-frontier/v4/budget.json schedule.window_gpu_hours_cap 24.0; A138 episode_accounting.v4_replayed 48, .v4_mismatches [];
% zoo 638/638 and 696/696 per issue-139-closeout.md L113-115; A138 primary.verdict POSITIVE. v5 window "Total 9.06 / 40 GPU-h."
% issue-140-closeout.md L72; cap outputs/choice-frontier/v5/budget.json schedule.window_gpu_hours_cap 40.0; A140 primary.verdict NOT_SEPARATED;
% episode_accounting.{v2,p2}.episodes_replayed 24 / 22, replay_mismatches [] (issue-140-closeout.md L53). Never summed (B4 schedule.accounting.transfers).
% Evidence: v6 window "2.245 / 30" issue-141-closeout.md L77 (stored-row sum of outputs/choice-frontier/v6/budget.json attempts[0..4].gpu_hours,
% cap schedule.authorization.window_gpu_hours_cap 30.0); outputs/choice-frontier/v6/metrics/analysis.json primary.v2.verdict EQUIVALENT, primary.p2.verdict INCONCLUSIVE,
% episode_accounting.{v2,p2}.episodes_replayed 24 / 22, replay_mismatches [] (46/46, issue-141-closeout.md L66). External v2: outputs/external-audit/v2/budget.json
% total_gpu_hours 2.8910, .schema.gpu_hours_cap 12.0; outputs/external-audit/v2/lfs/identity-audit.json paper_budget.verdict SATURATED; replay.reference {80, match 80},
% replay.random_valid {400, match 400}, mismatch 0 (lfs/replay.json). Never summed (.schema.independent).

% Paragraph purpose: Show that the ledger charges failed attempts, with the synthesis-snapshot attempt count.
The program ledger charges failed attempts as well as successful ones, so its total does not hide unsuccessful runs. At the synthesis snapshot, the program recorded 135 attempts (96 succeeded, 36 failed, 3 cutoff), and failed and cutoff attempts retain their hours in the ledger. Every attempt ended before the absolute scheduler cutoff of 2026-09-21T11:55:19Z.

% Paragraph purpose: Credit the corruption (#126) and seed-replication (#127) additions to the program ledger.
Two later addenda brought the program total to 68.52/336 GPU-h. The observation-corruption suite added 2.22 GPU-h to the generalization branch, and the training-seed replication addendum (closed 2026-09-21) added 12.79 GPU-h and 534 independently replayed episodes covering two cells. Appendix~\ref{app:statistics} reports the seed-variance results.
% Evidence: docs/experiments/expanded-study/issue-126-closeout.md records
% 2.2215 GPU-h added and 6.0404/32 branch spend. issue-127-closeout.md records
% 12.7863 GPU-h added and 68.5244/336 total.
% Status 2026-09-21: the v1 full scopes for generalization/robustness (#96/#98) and second backbone (#102/#103)
% were withdrawn (issues CLOSED, supervisor decision) and will not be executed. Only the reduced v2 scopes have executed evidence.

% Paragraph purpose: Show that each follow-up window stayed under its own cap on a separate ledger, then give each window's spend and replay coverage.
\paragraph{Follow-up windows.}
Each follow-up window has its own cap and ledger, none is added to the program total, and every window finished under its cap. Failed attempts retain their hours in every window. The native-arms window (v1 probes, then v2 training and evaluation) spent 3.2086 of 12 GPU-h and independently replayed its 216 new bindings (144 model episodes and 72 CPU random-valid episodes) and 72 reused comparators. The R1-R4 follow-up window spent 7.7998 of 36 GPU-h across two ledgers. Its cap was resized from 20 GPU-h before any outcome, because its frozen worst-case estimate of 27.04 GPU-h including spend exceeded the initial cap, and realized cost landed at 29\% of that estimate. The choice-frontier window spent 1.069 of 12 GPU-h (training 0.81, smoke 0.052, model evaluation 0.208), with corpus derivation and controls on CPU. It replayed its 144 panel episodes and re-derived its 58 preparation episodes bit-exact. The comparator zoo launched no GPU job and replayed all 432 zoo bindings, of which 342 are new CPU episodes and 90 are reused 2$\times$ random-valid controls. The three-seed replication and the held-out and unscreened panels shared one window, which spent 15.22 of 24 GPU-h. The DAgger ablation window spent 9.06 of 40 GPU-h, including 1.79 GPU-h charged for a cancelled collection. The zero-shot window spent 2.25 of 30 GPU-h, and the LLM-First Search audit spent 2.89 of 12 GPU-h on its reference server with no paid API use.
% (moved out of prose; kept in the preceding addenda paragraph and Table tab:appendix-receipts): program total 68.52/336 GPU-h, evidence as in the addenda comment above.
% Evidence: 15.22 / 24 issue-139-closeout.md L144; 9.06 / 40 issue-140-closeout.md L72; 1.79 GPU-h issue-140-closeout.md L80
% (B5 attempts[2].gpu_hours 0.8974 and attempts[3].gpu_hours 0.8973, status failed, cancelled under A1).
% Evidence: 2.25 / 30 issue-141-closeout.md L77 and outputs/choice-frontier/v6/budget.json attempts[0..4].gpu_hours; 2.89 / 12
% outputs/external-audit/v2/budget.json total_gpu_hours 2.8910, .schema.gpu_hours_cap 12.0, .schema.paid_api_usd 0.0 (issue-142-closeout.md L108).
% Window ids: native arms #128 (v1 probes) / #129 (v2); R1-R4 window #130; choice frontier #132; zoo #133.
% Evidence: issue-129-closeout.md (Delivered: 144 GPU model episodes, 72 CPU random_valid, 72/72 comparators, 3.2086 / 12.0);
% issue-130-closeout.md Verification & budget (27.04 GPU-h estimand, 29%) and the #131 status comment (7.7998 / 36, resized from 20);
% outputs/choice-frontier/v1/budget.json and the #131 #132-status comment (train 0.81 + smoke 0.052 + model eval 0.208);
% outputs/choice-frontier/v1/preparation/visual-choice-frontier/audit.json (episodes_rederived 58);
% #133 closing comment (432/432, 342 new, 90 reused, no GPU).
% Evidence: #134 CPU only (issue-134-closeout.md "no GPU was used", LED-3); identity-audit-submission-order.json validation (E134-08).
% Evidence: #135 outputs/choice-frontier/v2/budget.json attempts[].gpu_hours, closing comment "1.19/4 GPU-h" (LED-1); episode_accounting (E135-30, E135-31, E135-40).
% Evidence: #136 outputs/choice-frontier/v3/budget.json attempts[].gpu_hours, closing comment "4.84/30 GPU-h" (LED-2); episode_accounting (E136-31).
% Evidence: #137 CPU only (issue-137-closeout.md "No GPU", LED-4); scienceworld/identity-audit.json pairs_checked 1000, tasks 200 (E137-02, E137-03).

% Paragraph purpose: Map each new pinned evidence path to the claim class it supports.
\paragraph{Artifact index.}
Each follow-up claim class in the paper traces to a pinned evidence file, and Table~\ref{tab:appendix-artifacts} gives the mapping. Paths are relative to the evidence repository.

\begin{table}[t]
\caption{\textbf{Each follow-up claim class traces to pinned evidence files.} Filenames are relative to each directory prefix and the evidence repository, and claim classes identify their use.}
\label{tab:appendix-artifacts}
\centering\footnotesize
\setlength{\tabcolsep}{4pt}
\begin{tabular}{@{}p{0.56\linewidth}p{0.37\linewidth}@{}}
\toprule[0.4mm]
\rowcolor{mygray}\textbf{Relative file / directory} & \textbf{Claim class} \\
\midrule
\multicolumn{2}{@{}l}{\path|outputs/native-arms/v1/|} \\
\path|evaluation/|: \path|evaluation.json|, \path|analysis.json| & Clean / corruption contrasts \\
\path|identity-audit.json| & Enumeration identity \\
\path|stress-evaluation.json| & Menu / R2 / replay \\
\path|r2-decomposition.json| & Termination modes \\
\path|budget.json|; R1--R4 v1/v2 ledgers & Window accounting \\
\midrule
\multicolumn{2}{@{}l}{\path|outputs/expanded-study/v1/|} \\
\path|second-backbone-v2/evaluation/| & Backbone BFS \\
\path|second-backbone-v3/evaluation/| & Backbone additive \\
\midrule
\multicolumn{2}{@{}l}{\path|outputs/choice-frontier/|} \\
\path|v1/evaluation/| (evaluation, identity-audit, analysis, final-audit) & Gate / panel / replay \\
\path|v1/preparation/visual-choice-frontier/| (report, audit) & Corpus derivation \\
\path|v1/budget.json| & Window accounting \\
\path|o4/zoo/manifest.json| & Zoo coverage / replay \\
\path|o4/metrics/| (analysis, stored-posthoc-analysis, optimal-costs, native-four-algorithm-reference-costs) & Zoo M1--M5 / Holm / costs \\
\path|v4/seeds/metrics/analysis.json| & Three-seed validation \\
\path|v4/panels/metrics/analysis.json| & P2 / P2u / concentration \\
\path|v5/metrics/analysis.json| & DAgger ablation \\
\path|v4/budget.json|; \path|v5/budget.json| & Window accounting \\
\path|v6/metrics/analysis.json| & Zero-shot policy \\
\path|v6/budget.json| & Window accounting \\
\midrule
\multicolumn{2}{@{}l}{\path|outputs/external-audit/v2/|} \\
\path|lfs/identity-audit.json| & LFS identity \\
\path|budget.json| & External v2 ledger \\
\midrule
\multicolumn{2}{@{}l}{\path|docs/experiments/| and ticket \#133} \\
\path|native-arms/issue-129-closeout.md|; \path|issue-130-closeout.md| & Native / stress receipts \\
\path|choice-frontier/issue-132-closeout.md| & Binding erratum \\
\path|choice-frontier/|\newline\path|metric-adjudication-2026-09-23.md| & Metric adjudication \\
Ticket \#133 closing comment & Zoo numbers / boundary \\
\bottomrule
\end{tabular}
% Evidence: original path inventory and claim classes retained from the source comments below.
\end{table}
% Evidence: paths confirmed present in the evidence repository on 2026-09-23; the #133 closing comment is
% https://github.com/Sino-Huang/multimodality_on_planning/issues/133#issuecomment-5788941573.


\section{Statistics}
\label{app:statistics}

% Paragraph purpose: Say what this appendix holds and why a reader would look here.
This appendix states the interval, test, margin, and tiny-strata conventions behind every reported contrast, for readers judging what an interval licenses.

% Paragraph purpose: State that main-program intervals are descriptive bounds, and the default contrast procedure behind them.
\paragraph{Default contrast procedure.}
Main-program intervals are descriptive bounds, not evidence of broad superiority, because they come from one training seed and panels of 24 or fewer instances. Contrasts are paired by whole task instance, and paired differences use a bootstrap with 10,000 resamples at seed 1729 and 95\% percentile intervals.

% Paragraph purpose: Pin the BFWS gate bound to its two-sided percentile interval.
The BFWS gate's reported lower bound is the 0.025 quantile, i.e., the lower endpoint of a two-sided 95\% percentile interval. The gate follows the default procedure, with confidence level 0.95 and 10,000 paired whole-instance resamples (bootstrap seed 1729).
% Evidence: BFWS gate is issue #59. configs/experiments/bfws_phase_threshold_v1.json bootstrap {confidence_level 0.95, resamples 10000, seed 1729,
% interval paired_percentile_bootstrap}; docs/issue-59-bfws-structural-gate.md L250-253 (lower bound 0.2).

% Paragraph purpose: State the tiny-strata rule: no intervals below 8 episodes, and no silent pooling or dropping.
\paragraph{Tiny-strata rule.}
No intervals are computed over strata with fewer than 8 episodes, so such strata are reported as descriptive only. They are listed explicitly and never silently pooled or dropped.

% Paragraph purpose: State that no transfer comparison survives correction, then the test procedure.
\paragraph{Transfer tests.}
No transfer comparison survives multiple-testing correction, and this non-significance is not evidence of equivalence. Transfer comparisons use McNemar exact tests against the per-example base predictions, with Holm correction across all 36 adapter-by-benchmark comparisons, and the minimum adjusted $p$ in the executed transfer branch is 0.8156.
% Evidence: outputs/expanded-study/v1/transfer-paired-analysis.json min_holm_p 0.81561 over 36 comparisons.
% Evidence: 36 Holm-corrected McNemar comparisons per synthesis-v1/transfer-paired.csv (Table tab:results-transfer).

% Paragraph purpose: State the follow-up bootstrap seeds separately from the main-program default.
\paragraph{Follow-up bootstrap conventions.}
The follow-up windows fix their own bootstrap seeds, separate from the main program's seed 1729 default. The native-arms contrasts and the choice-frontier contrasts use a paired whole-problem bootstrap at seed 61813 with 10,000 resamples and 95\% percentile intervals, and the second-backbone v3 contrasts (stress test R3) use seed 1729 with 10,000 resamples. The comparator-zoo metrics use a task-cluster bootstrap that draws nine tasks with replacement at seed 133 for 10,000 draws, keeping both additive cells and all random seeds inside their task cluster.
% Window ids: native arms #129, choice frontier #132, R3 #130, zoo #133.
% Evidence: configs/experiments/native-arms/native-arms-v2-protocol.json analysis.bootstrap (seed 61813, 10000, percentile);
% docs/experiments/choice-frontier/issue-132-design.md section 6 (seed 61813); outputs/expanded-study/v1/second-backbone-v3/evaluation/analysis.json
% (bootstrap seed 1729, 10000); docs/experiments/choice-frontier/issue-133-protocol.md (nine tasks, seed 133, 10,000 draws).

% Paragraph purpose: State the materiality rule and the non-equivalence sentence requirement.
\paragraph{Materiality and non-equivalence.}
A follow-up contrast is material when its 95\% interval excludes zero, and equivalence can be read only against a pre-registered margin. The follow-up protocols through the comparator zoo pre-registered no equivalence margin, so a non-material contrast there supports neither superiority nor equivalence. The later choice-frontier protocols pre-register an equivalence margin of $\pm 0.05$ for $D$ (adapter minus random-valid M1) and for the DAgger change $\Delta$, and none for the separation $S$ from an exact-$\epsilon$ rung, so a non-separated $S$ is not equivalence to the rung. Wherever the text reports a non-material contrast or a comparison that fails Holm correction, an adjacent sentence states that non-significance is not equivalence.
The exact sign-flip $p$ value is likewise descriptive. It enumerates all $2^n$ sign vectors of the $n$ per-task differences, is two-sided, and tests a zero difference, so it never supports equivalence.
% Evidence: outputs/choice-frontier/v6/metrics/analysis.json per_panel.*.*.sign_flip.{exact,sign_vectors,statistic}; pooled.D3.sign_flip (8388608 sign vectors).
% Evidence: A138 primary.equivalence_margin 0.05; A139 heldout_p2.primary.equivalence_margin 0.05, unscreened_p2u.primary.equivalence_margin 0.05;
% A140 co_primary.equivalence_margin 0.05; separation verdict rules SEPARATED_ABOVE / SEPARATED_BELOW / NOT_SEPARATED only (A140 primary.rule).
% Evidence: native-arms-v2-protocol.json analysis.bootstrap.materiality ("interval excludes 0"); issue-132-design.md section 6;
% issue-132-closeout.md Read of the evidence item 3; issue-133-protocol.md ("no equivalence or generalized negative capability claim").

% Paragraph purpose: Apply the tiny-strata rule to the single-task-driven contrast and the one-task M5 subset.
\paragraph{Tiny strata in the follow-ups.}
The fewer-than-8 rule also governs discriminative subsets. The choice-frontier learned-minus-random-valid contrast is computed over 18 cells, but it rests entirely on storage, the only panel task where random-valid's seed frequency lies strictly between 0 and 1 (0.2). A one-task discriminative subset falls under the rule, so that contrast is descriptive. The zoo's M5 subset, defined from controls only, holds the same single task and is reported descriptively.
% Evidence: issue-132-closeout.md Erratum item 3 (storage 0.2, one discriminative task, below the <8 rule);
% outputs/choice-frontier/o4/metrics/analysis.json m5_informative_tasks_D (storage-compact-919000) and the #133 closing comment.

% Paragraph purpose: State the #133 Holm and dominance test procedure.
\paragraph{Zoo dominance tests.}
A positive learned-versus-comparator dominance claim in the comparator zoo must survive Holm step-down correction over fifteen tests (five comparators by three metrics), applied to paired task-cluster bootstrap one-sided $p$ values. Undefined or uninformative co-solved contrasts prohibit a positive dominance claim and never count as wins. All zoo metrics are exploratory, and those computed on the stored choice-frontier episodes are post-hoc.
% Evidence: docs/experiments/choice-frontier/issue-133-protocol.md (claim unit, 15 tests, undefined contrasts, post-hoc labelling).

% Paragraph purpose: State the seed-variance addendum's scope, its frozen seeds, and its descriptive status.
\paragraph{Seed-variance addendum.}
A prospective training-seed replication measured variance for exactly two cells, and its cross-seed aggregation is descriptive ($n = 3$ seeds). Seeds 17, 29, and 71 were frozen before any seed-29 or seed-71 artifact existed, and all 534 new episodes were independently replayed.
% Evidence: outputs/expanded-study/v1/seed-replication/analysis.json (seeds.original 17, additional_prospective 29/71;
% cross_seed_aggregation "descriptive only (n = 3 seeds)"); docs/experiments/expanded-study/synthesis-v1/seed-variance-addendum.md (534 episodes).

% Paragraph purpose: State that Cell A's process-SFT gain over the base and deficit against random-valid repeat at every seed, and define the cell.
In Cell~A, process SFT beats the base and trails random-valid at every seed (Table~\ref{tab:appendix-seeds}). Because the cell is additive, the identity audit reads that success as validity rather than choice. Cell~A is the greedy $\times$ multimodal replication cell (\texttt{best\_first\_add\_greedy} $\times$ multimodal-state) evaluated on the 24-problem panel, and each seed trained a fresh LoRA adapter with new initialization and dropout randomness on the identical frozen 512-record membership.

\begin{table}[t]
\caption{\textbf{Process SFT beats the base and trails random-valid at all three seeds (descriptive).} Cell~A, greedy $\times$ multimodal. Successes and paired differences compare process SFT with the base and random-valid.}
\label{tab:appendix-seeds}
\centering\footnotesize
\setlength{\tabcolsep}{4pt}
\begin{tabular}{rccc}
\toprule[0.4mm]
\rowcolor{mygray}\textbf{Seed} & \textbf{Successes / 24} & \textbf{SFT $-$ base} & \textbf{SFT $-$ random-valid} \\
\midrule
17 & 21/24 & $+0.875$ & $-0.125$ \\
29 & 21/24 & $+0.875$ & $-0.125$ \\
71 & 22/24 & $+0.917$ & $-0.083$ \\
\bottomrule
\end{tabular}
% Evidence: outputs/expanded-study/v1/seed-replication/analysis.json baseline_process_sft.per_seed.{17,29,71}.successes (21, 21, 22 of 24)
% and .contrasts.process_sft_minus_{pretrained_base,random_valid}.difference; bootstrap seed 1729, 10000 resamples.
\end{table}

% Paragraph purpose: State that the Cell B null replicates at every seed.
Cell set~B replicates a null result at every seed. It compares DAgger iteration-1 with the exposure-matched continued-SFT control under BFS across all three modalities, and every arm at every seed has 0/72 unseen successes. The paired contrast (DAgger minus continued SFT) is therefore $+0.000$ $[0.000, 0.000]$ at each seed, which is not evidence of equivalence.
% Evidence: outputs/expanded-study/v1/seed-replication/analysis.json dagger_iteration_1.per_seed.*.arms.*.unseen (successes 0, episodes 72)
% and .contrasts.dagger_minus_continued_sft_unseen (difference 0.0); aggregated.contrast_point_estimates_per_seed.

% Paragraph purpose: Scope the seed-replication boundary to the two replicated cells.
This addendum supersedes the blanket single-training-seed boundary for these two cells only. Every other enumeration-contract cell in the program remains single-seed (seed 17) evidence.
% TODO(seed-coverage): Seed replication beyond these two cells is not
% authorized under the current program. All remaining cells are single-seed.


\section{Modality Adapters and Observation Contract}

% Paragraph purpose: Say what this appendix holds and why a reader would look here.
This appendix defines the observation adapters and what each one exposes to the policy, which bounds any reading of the modality results.

% Paragraph purpose: State the decision-sufficiency requirement every adapter meets, then define the adapter map.
Every adapter must be decision-sufficient: its observation must expose every fact the exact reference uses to resolve the current step under the declared algorithm. Three fixed Modality Observation adapters are declared. Each maps a decision context $c_t$ (the canonical task context plus the decision-relevant view of Search Memory $M_t$) to an observation $o_t = \phi_m(c_t)$ presented to the policy, and the adapter stays fixed for an entire training or evaluation cell.

% Paragraph purpose: text-state adapter.
\paragraph{text-state.}
The state, goal, and search context are serialized as a compact relational representation. Object identities, predicate names, and all Search Memory contents (frontier, visited set, algorithm-specific state) appear as structured text.

% Paragraph purpose: visual-state adapter.
\paragraph{visual-state.}
The state and goal are rendered as a pair of 128$\times$128 pixel images, a rendered state image and a partial-goal constraint image. Object labels are not embedded in the rendered image. Search Memory contents and the grounded candidate list remain in the textual serialization, even in this adapter. This is the published frozen \texttt{visual-state} arm that the native arms below are compared against.

% Paragraph purpose: State that the figures show the policy's own scenes without retouching, and how each figure's images were checked.
The figures show the scenes the policy saw, enlarged but never retouched. Figure~\ref{fig:teaser} displays unchanged cached 128-pixel policy scenes enlarged by nearest-neighbour sampling. The validation filmstrips in Appendix~\ref{app:qualitative} show label-free 256-pixel re-renders, used only after the same scenes at the policy's 128-pixel resolution matched the cached inputs pixel for pixel. The held-out puzzle shows the original cached 128-pixel frames instead, because re-rendering its initial state changes tile glyphs. Figure scripts retain the exact image paths and checks.

% Paragraph purpose: multimodal-state adapter.
\paragraph{multimodal-state.}
The observation combines the text-state serialization and the visual-state rendered images. The policy receives the relational text and the image pair jointly.

% Paragraph purpose: Record the executed native-arms v2 adapters and how they reduce the textual payload.
\paragraph{Native image-only arms.}
Two further adapters executed under the native-arms v2 contract. \texttt{visual-nomem-state} keeps only the static task interface and the grounded, unscored candidate menu as text, and carries state facts in the current rendered scene alone. \texttt{visual-seq-state} adds up to $K{=}8$ rendered history frames, each labelled with its absolute path depth. Neither receives textual Search Memory, accepted deltas, $g$/$h$/priority scalars, candidate scores, or membership flags.
% Evidence: docs/experiments/native-arms/native-arms-design.md (removed and retained keys, history frame labels);
% configs/experiments/native-arms/native-arms-v2-protocol.json history_window.

% Paragraph purpose: State the text-scaffold caveat: search bookkeeping always travels as text, and the corruption suite agrees.
\paragraph{Text-scaffold caveat.}
In all three declared adapters, including visual-state, search bookkeeping travels as text. Rendered images carry the planning state and goal visuals, but the frontier, visited set, novelty counters, and candidate list are always serialized as structured text. The observation-corruption suite is consistent with this (Table~\ref{tab:results-corruption}). Copyability-preserving text shuffling drops success by up to $-0.944$ (text) and $-0.722$ (multimodal), while visual corruption costs at most $-0.111$ with intervals that include zero.
% Evidence: observation-corruption suite is issue #126; shuffled deltas from Results Table tab:results-corruption
% (Text-shuffled rows); visual rows max -0.111.

% Paragraph purpose: State what the executed native arms changed in the text-scaffold caveat, with the menu reading scoped to what the receipts identify.
The native arms narrow this caveat to the candidate menu without removing it. With the textual bookkeeping and scalars removed, both arms solve 18/18 clean tasks with zero invalid operations, and blanking or degrading every image of the seq arm gives a corrupted-minus-clean contrast of $0.000$ $[0.000, 0.000]$ (Appendix~\ref{app:results-native}). As a post-hoc interpretation, the grounded unscored menu is the remaining text channel that could carry the signal. Menu manipulation shows that the policies do not reject schema-valid inapplicable menu entries (pick-rate at chance), and on the additive cells submission-order invariance makes any listed candidate a correct emission (Appendix~\ref{app:results-menu}). The caveat therefore still holds for every adapter that presents the menu, and no executed adapter isolates an image-only decision channel.
% Evidence: docs/experiments/native-arms/issue-129-closeout.md (18/18 per arm, zero invalid, corruption 0.000 [0.000, 0.000],
% menu-leak interpretation); docs/experiments/native-arms/issue-130-closeout.md R1. Same numbers as Section~\ref{sec:results}.

% Paragraph purpose: State that adapters match format, not information, so no intrinsic modality effect is established, with the information-loss evidence.
\paragraph{Information-loss boundaries.}
Matching the adapter between training and evaluation ensures identical observation format, not information parity, so the available evidence does not establish semantic parity across adapters or an intrinsic modality effect. The 128-pixel unlabelled rendered views can lose symbolic identities, and static-context and goal fields of the observation contain text in every adapter. The loss also differs by modality. The P3 name-compression family in the generalization-v2 suite is information-lossy, and family-wide (24 eligible P3 variants), 24/24 variants are text-lossy, 18/24 are multimodal-lossy, and 0/24 are visual-lossy.
% Evidence: docs/experiments/expanded-study/issue-124-closeout.md L111-114 (24 is the family-wide eligible P3 set; the admitted
% subset of 5 reported in Section~\ref{sec:results-secondary} is a different stratum).
% #128 image-only proposal marker retired: visual-nomem-state and visual-seq-state executed under the #129 v2
% amendment (issue-129-closeout.md); the #128 v1 gate is retained as history.


\section{Retained-Infrastructure Feasibility Notes}

% Paragraph purpose: Scope the section: feasibility evidence, not efficacy.
This appendix defines the verified infrastructure components behind corpus generation, evaluation, and evidence production. They are feasibility evidence for the runtime and rendering pipeline and cannot support a Search Process Policy efficacy claim.

% Paragraph purpose: Define the plan pipeline terms.
\paragraph{Plan pipeline.}
\textbf{Plan Submission} transmits a supplied plan from the project adapter to a planning or rendering backend. \textbf{Plan Interpretation} converts the submitted plan into the backend's ordered action representation. An \textbf{Action Sequence} is the canonical identity of a supplied plan as a normalized, ordered sequence of grounded action predicates. Casing and whitespace are not part of its identity.

% Paragraph purpose: Define the rendering pipeline terms.
\paragraph{Rendering pipeline.}
\textbf{Render Production} creates a visualisation output and its rendered image for a planning state. \textbf{Render Validation} confirms that the produced visualisation and image satisfy declared structural and semantic checks. \textbf{Plan Provenance} links a produced visualisation to the exact supplied plan that the backend interpreted.

% Paragraph purpose: Define the certification and evidence terms.
\paragraph{Certification and evidence.}
\textbf{Integration Certification} is the finding that Plan Submission, Plan Interpretation, Plan Provenance, Render Production, Render Validation, and required execution-isolation checks all passed. An \textbf{Attempt} is one authorized, immutable execution with its own identity and retained outcome, whether successful or not. An \textbf{Evidence Bundle} is the self-contained records from an Attempt that allow an independent reproducer to verify its claims offline. The retained infrastructure supplies deterministic compact delta/event evidence with semantic replay while preserving search semantics, serving the replay-determinism requirement of the Trusted Search Runtime.
% Evidence: compact delta/event evidence per issue #110.

% Paragraph purpose: Note historical CGAS terms and their claim boundary.
\paragraph{Historical terms.}
Planning Certificate, Joint Action-and-Certificate SFT, Adaptive Scaffolding, Support Route, Live Memory, Verified Joint Step, and Route Label are historical CGAS terms for retained infrastructure internals. They cannot support any current Search Process Policy efficacy claim and are not referenced elsewhere in this manuscript.


\section{Choice-Frontier Contract and Comparator Selectors}
\label{app:cf-contract}

% Paragraph purpose: Say what this appendix holds and why a reader would look here.
This appendix specifies the choice-frontier contract, its corpus derivation, its metrics, and the comparator selectors, for readers reproducing the measurement that Section~\ref{sec:results-ladder} validates.

% Paragraph purpose: The menu withholds scores, positions, and search bookkeeping, so a choice can rest only on the scenes and task pages.
\paragraph{Menu and observation.}
The choice-frontier observation withholds every score, position cue, and search counter, so a choice can rest only on the rendered scenes and the task pages. At each decision the runtime exposes the live frontier as an unscored menu of opaque labels $c_0, \ldots, c_K$. The labels are assigned after a seeded Fisher--Yates permutation of the generation-serial member order, with a per-decision seed derived from master seed 51131 and a hash of the member state identifiers, so a label carries no positional information. The observation holds static task-context pages, the unlabelled 128px initial-state scene, one unlabelled 128px scene per frontier state in menu order, and the partial-goal pages. Its text payload is limited to the algorithm name, the label list, the representation tag, the schema version, and the frozen legend, and a leak check fails closed on any forbidden key marker.
% Evidence: docs/experiments/choice-frontier/issue-132-design.md section 2 (menu, seed rule, payload, assert_no_text_leak).

% Paragraph purpose: The output contract admits one offered label per decision, and the first invalid emission ends the episode.
\paragraph{Output and validation.}
The output contract admits exactly one offered label per decision, and the first invalid emission ends the episode. The policy emits canonical JSON with the single key \texttt{expand\_choice} and a label value. The runtime rejects malformed JSON, a wrong key set, a non-string label, a label not offered by the pending menu, and a label bound to a non-live member, and it records the rejection as \texttt{deterministic\_invalid\_operation}. An episode otherwise ends as \texttt{goal\_reached}, \texttt{frontier\_exhausted}, \texttt{decision\_budget\_exhausted}, or \texttt{expansion\_budget\_exhausted}.
% Evidence: issue-132-design.md section 2 (output contract, rejections, termination reasons).

% Paragraph purpose: The runtime performs every expansion, so one decision is one choice, and the goal-at-selection rule fixes the exact reference's decision count.
\paragraph{Expansion and budget.}
The policy only chooses and the runtime does everything else, so one submitted choice is one decision. On acceptance the runtime removes the chosen member from the frontier and expands it atomically. Candidate generation, serial assignment, admission, duplicate detection, reopening, and expansion bookkeeping are inherited unchanged from the frozen additive best-first controller, which is not modified, and per-candidate admissions stay runtime-internal. Selecting a member that satisfies the goal solves the task without an expansion and without budget charge. The exact episode therefore takes its reference expansion count plus one decisions. Decision and expansion caps both equal 2$\times$ the reference expansion count.
% Evidence: issue-132-design.md section 2 (expand_member, inherited BestFirstController internals, goal at selection,
% decision_cap = expansion_cap = 2 x reference_expansions); issue-132-closeout.md Erratum items 1 and 5.

% Paragraph purpose: Each training record teaches the exact reference's own choice, derived by replaying the frozen exact traces through the new contract.
\paragraph{Corpus derivation.}
Each training record teaches the exact reference's own choice, the label bound to the heap head. The training corpus was derived on CPU by replaying the frozen exact traces through the new contract and verified expansion-for-expansion against them (58 episodes bit-exact). Records are selected through the frozen matched-modalities membership, which indexes enumeration-contract decisions, and each such decision maps to the expansion whose candidate set it submitted into. Unique (task, expansion) pairs in first-occurrence order form the base set, and a round-robin over the frozen task order tops the set up to 512 training records per algorithm. Diagnostics are the unique expansions of the frozen diagnostic identifiers, with no top-up.

% Paragraph purpose: Training ran only after five fail-closed gates passed, and the pinned preparation report records the corpus counts.
Training ran only after five fail-closed gates passed. The gates check expansion-for-expansion agreement with the stored exact traces, a token budget of measured input plus 384 within 32,768, the text-leak check, image resolution, and identity with the frozen membership. The pinned preparation report records outcome PASS over 29 tasks, 1,046 records (512 training and 11 diagnostic per algorithm), 3,544 scenes, and 58 episodes re-derived bit-exact. The smoke gate accepted 40/40 model calls.
% Evidence: issue-132-design.md section 3 (membership freeze rule, teacher target, five gates);
% outputs/choice-frontier/v1/preparation/visual-choice-frontier/report.json counts (records 1046, tasks 29, scenes 3544,
% episodes_verified 58, outcome PASS) and audit.json (episodes_rederived 58); issue-132-closeout.md (smoke 40/40).
% Source conflict: issue-132-design.md anticipated 26 training tasks plus 1 diagnostic-only task and a 54/54 exact-trace gate;
% the pinned report records 29 tasks and 58 verified episodes. The artifact values are used, matching Section~\ref{sec:results}.
% Provenance sentence identical across Results, Experimental Design, and this appendix; report.json expansion_gate
% "every derived episode matched its stored exact trace expansion for expansion". issue-132-closeout.md L66 "derived independently
% of the stored exact traces" conflicts with the artifact and is not used.

% Paragraph purpose: Both adapters trained to completion under the frozen recipe, their reports were recovered without retraining, and every episode is independently replayed.
\paragraph{Training and replay.}
Both adapters trained to completion under the frozen recipe, and their episodes are independently replayed. Each of the two cells trains a LoRA adapter (rank 64, $\alpha$ 128) for 16 updates on its 512 records at seed 17 with the frozen expanded-program recipe. Both jobs completed all 16 steps and saved final adapters, and only the later post-save summary construction failed. Their reports were therefore recovered from the primary training-state artifacts with provenance marked, no retraining was needed, and the training audit passes on the recovered reports. At finalize, independent replay re-derives every menu binding, payload, and trusted transition from the stored events and checks pixel bindings against the retained scenes.
% Evidence: issue-132-design.md sections 4 and 6; issue-132-closeout.md Execution record and Budget accounting (failed train attempts,
% recovered reports, audit-train PASS); outputs/choice-frontier/v1/training/visual-choice-frontier/audit.json.

% Paragraph purpose: The comparator selectors replace only the choice, so every other search step stays additive, and their replay omits image bindings.
\paragraph{Comparator selectors.}
The comparator selectors replace only the choice. Each selector's choice is mapped back to the menu binding, and every candidate still enters through the additive admission path, so frontier admission, duplicate detection, $h_{\rm add}$ computation, reopening, and serial ordering stay additive for \texttt{bfs-order} and novelty-first selection alike. Novelty-first selection uses positive PDDL state atoms as its novelty features, and its seen-feature sets advance only when a state is expanded, including the initial expansion. Each new episode is re-run against an independently constructed PDDL authority, which checks menus, decisions, admissions, and result dictionaries. This replay covers the model input, menu, actions, admissions, and budget, but selector episodes neither generate nor replay image bindings. The optimal costs $C^*$ come from one offline shortest-plan computation per task with replay-verified plans, and the native four-algorithm reference costs are recorded separately from the zoo.
% Evidence: docs/experiments/choice-frontier/issue-133-protocol.md (label mapping, additive admission path, novelty features,
% independent PDDL authority, one offline C* computation per task); #133 closing comment (no image-binding replay, native references
% checked separately); outputs/choice-frontier/o4/metrics/{optimal-costs,native-four-algorithm-reference-costs}.json.

% Paragraph purpose: The protocol fixes each metric's edge cases so every arm is scored by one rule, and the goal-at-selection rule fixes the exact reference's M1 at 0.875.
\paragraph{Metric computation.}
The protocol fixes each metric's edge cases so that every arm is scored by the same rule. Under M1 an episode solves at multiplier $m$ when its first goal selection occurs within $\lfloor m R_t \rfloor$ decisions with at most $\lfloor m R_t \rfloor$ preceding expansions, where $R_t$ is the additive exact-reference expansion count. Unsolved 2$\times$ runs count as failures at every $m$, the same prefix truncation applies to every arm, and task and algorithm weights are equal. M2 also has a censored form, in which failures are right-censored at the budget rather than assigned a fictitious expansion count. M3 counts transitions along the selected goal state's trusted predecessor chain and is undefined, never imputed, for unsolved episodes. M4 measures agreement with the current additive heap head over decisions with at least two live choices, against a chance rate equal to the decision-weighted mean of $1/K$. M5 restricts to tasks where the mean 2$\times$ random-valid solve indicator lies strictly between 0 and 1.
The goal-at-selection rule fixes the reference's own M1. Goal selection is the $(R_t{+}1)$-th decision, so the exact reference fails at $m = 1$ and scores M1 $= 0.875$ by construction. An arm that reaches the goal in fewer decisions can therefore exceed the reference, as the privileged hadd-greedy selector does on the validation, held-out, and unscreened panels (0.882, 0.928, and 0.906).
% Evidence: goal-at-selection rule in the Expansion and budget paragraph above (issue-132-design.md section 2); A139 arms.{p135,p2,p2u}.exact_reference.m1 0.875;
% arms.{p135,p2,p2u}.hadd-greedy.m1 0.8823 / 0.9284 / 0.90625, rounded half-up to 3 dp from the named keys.
% Evidence: docs/experiments/choice-frontier/issue-133-protocol.md, Pre-registered metrics and claims.


\section{Results Detail: Gate Receipts and Qualification Gates}
\label{app:results-gates}

% Paragraph purpose: Say what this appendix holds and why a reader would look here.
This appendix gives the program's receipts, the follow-up ledgers, and the qualification gates that the accounting paragraph of Section~\ref{sec:results-gates} summarizes.

% Paragraph purpose: The expanded program stayed well inside its compute cap and every executed branch reconciles with zero missing evidence.
The expanded program stayed well inside its compute cap, and every executed branch reconciles with zero missing evidence. It ran on 2 A100 80GB GPUs under a 336 GPU-hour cap with an absolute scheduler cutoff, and spent 68.5244 of those hours in total, 53.52 (15.9\%) at the synthesis snapshot and the remainder on the corruption and seed-replication addenda. Of the seven admitted branches, five completed at v1 scope. The other two, the second backbone and generalization/robustness, stopped at frozen v1 cost-admission gates with terminal \textsc{valid\_stop} receipts and then completed at reduced v2 scope under new versioned protocols. Independent replay covers 1,152/1,152 baseline episodes, 405/405 DAgger episodes, and 90/90 successor-prediction episodes. Table~\ref{tab:appendix-receipts} gives frozen scope, status, coverage, and ledger spend per branch.
% Evidence: 53.52 / 336 and 15.9% from docs/experiments/expanded-study/synthesis-v1/README.md compute-accounting claim (budget.csv;
% verification.json LEDGER); 1,152 / 405 / 90 replay counts from synthesis-v1/branch-reconciliation.csv and verification.json.
% Evidence: synthesis snapshot #120; addenda #126/#127.
% Evidence: 68.5244 / 336 program total after the addenda per docs/experiments/expanded-study/issue-127-closeout.md
% (12.7863 GPU-h added, 68.5244/336 total); same value as the Reproducibility statement and Appendix tab:appendix-receipts (68.52).
% Status 2026-09-21: issues #96/#98/#102/#103 are CLOSED (supervisor decision, v1 scope withdrawn); prose states the withdrawal.

% Paragraph purpose: The follow-up windows ran on their own ledgers, never summed with the program, and all stayed within their caps.
The follow-up windows ran on their own ledgers and are never summed with the main program. Most used separate GPU ledgers, and three ran on CPU only (the comparator zoo, the counterfactual audit, and the first external audit). Each stayed within its cap: the native-arms window spent 3.2086 of 12 GPU-hours, the R1-R4 follow-up window 7.7998 of 36, and the choice-frontier window 1.069 of 12.
% Window ids: native arms #128/#129, R1-R4 window #130, choice frontier #132, zoo #133.
% Evidence: Sigma attempts[].gpu_hours in outputs/native-arms/v1/budget.json (3.20857, cap 12.0), outputs/reviewer-v130/v1+v2/budget.json
% (3.98625 + 3.81354 = 7.79978, cap 36), outputs/choice-frontier/v1/budget.json (1.06939, cap 12); docs/experiments/native-arms/issue-129-closeout.md (3.2086 / 12.0);
% docs/experiments/native-arms/issue-130-closeout.md (states 7.80 GPU-h, i.e. the 7.7998 ledger value rounded; cap 36);
% docs/experiments/choice-frontier/issue-132-closeout.md and outputs/choice-frontier/v1/budget.json (1.069 / 12.0);
% issue Sino-Huang/multimodality_on_planning#133 closing comment (no GPU launched).
The R1-R4 cap was raised from 20 to 36 GPU-hours before any outcome, because its frozen worst-case estimate of 27.04 GPU-hours including spend exceeded the initial cap. Realized cost landed at 29\% of that estimate. Each follow-up closed with complete coverage and independent replay (Table~\ref{tab:appendix-receipts}).
% Evidence: docs/experiments/native-arms/issue-130-closeout.md, Verification & budget (27.04 GPU-h estimand, 29%).


% Paragraph purpose: The earlier BFS pilot showed executable BFS with zero gain over random-valid, which gated it off; the BFWS gate then passed, and the 45-task final evaluation manifest stayed untouched.
Three earlier gates fixed which algorithms the active matrix contains. The first, the bounded BFS pilot, showed executable BFS with zero gain over the strongest control and therefore closed with a governed \textsc{valid\_stop}. On a frozen outcome-blind 15-task BFS development panel, final checkpoints from five process-SFT seeds achieved 1.0 invariant-valid success with a 0.0 invalid-operation rate, while the base achieved 0.0 and random-valid achieved 1.0, so downstream BFS work was gated off. The BFWS structural gate then passed on its own frozen development panel (Section~\ref{sec:results-primary}), although the overlap between the BFWS and BFS development panels is not pinned. The companion manifest of the frozen 45-task final evaluation was fixed before outcomes and has never been accessed.
% TODO(panel-overlap): issue-54-operator-runbook.md and issue-59-bfws-structural-gate.md pin the selection rules but not their
% shared task identities. A 2026-09-23 search found no pinned post-#126 retention manifest. Count stays open.
% Evidence: docs/experiments/deadline-study/prior-evidence.json historical_bfs_stages[stage v8] (15 selected tasks, outcome VALID_STOP,
% metrics process_sft/random_valid invariant-valid success 1.0, base 0.0, per_seed 17/29/43/71/101); BFWS gate docs/issue-59-bfws-structural-gate.md L250-256.
% TODO: the authorized held-out final primary evaluation on the frozen
% 45-task manifest remains unexecuted; no final efficacy number exists.

% Paragraph purpose: The third gate retired optimal A*, and the five-seed BFS pilot predates the one-seed budget rule.
The third gate retired the optimal A* arm over $h_{\max}$ with landmark counting after its own governed \textsc{valid\_stop}. The two additive best-first variants supersede it, and its original heuristic-representation estimand remains unidentified. The five-seed BFS pilot above predates the budget rule, which authorizes one training run at seed 17 per enumeration-contract model and cell (Section~\ref{sec:design}).

\section{Results Detail: Primary-Matrix Tables}
\label{app:results-primary}

% Paragraph purpose: The two primary tables keep what the model learned apart from what the controls achieve, with no measure combined.
The two primary tables keep what the model learned apart from what the controls achieve, and no measure is combined with another. Table~\ref{tab:results-bfws} gives the BFWS development panel, and Table~\ref{tab:results-primary-matrix} gives the expanded 24-task baseline by algorithm and modality. Section~\ref{sec:results-primary} summarizes both.
% Evidence: sources docs/issue-59-bfws-structural-gate.md L250-256
% (Table tab:results-bfws) and synthesis-v1/baseline-summary.csv, baseline-contrasts.csv (Table tab:results-primary-matrix, tab:results-contrasts).

\begin{table}[t]
\caption{\textbf{Process SFT beats oracle-assisted random-valid on the BFWS development panel.} Development-stage success over 15 tasks for one fixed training seed, alongside the pretrained base and the exact reference. No final evaluation has been executed.}
\label{tab:results-bfws}
\centering\footnotesize
\setlength{\tabcolsep}{4pt}
\begin{tabular}{lc}
\toprule[0.4mm]
\rowcolor{mygray}\textbf{Condition} & \textbf{Success (15 tasks)} \\
\midrule
Pretrained base & 0.0 \\
Process SFT & 0.933 \\
Random-valid (oracle-assisted control) & 0.467 \\
Exact execution & 1.0 \\
\bottomrule
\end{tabular}
\end{table}

\begin{table}[t]
\caption{\textbf{Learned success appears only in the additive cells and never exceeds random-valid.} Descriptive enumeration-contract baseline by algorithm and modality. Successes, SFT invalid operations per decision, SFT expansions, and exact-reference decisions are separate columns. Random-valid (RV) and exact are controls, not model abilities.}
\label{tab:results-primary-matrix}
\centering\footnotesize
\setlength{\tabcolsep}{2pt}
\begin{tabular}{@{}ll|cccc|cc|c@{}}
\toprule[0.4mm]
\rowcolor{mygray}\textbf{Algorithm} & \textbf{Modality} & \textbf{Base} & \textbf{SFT} & \textbf{RV} & \textbf{Exact} & \textbf{SFT inv/dec} & \textbf{SFT exp.} & \textbf{Exact dec.} \\
\midrule
\multirow{3}{*}{BFS} & text & 0/24 & 0/24 & 15/24 & 24/24 & 24/66 & 15 & 1969 \\
 & visual & 0/24 & 0/24 & 15/24 & 24/24 & 24/69 & 17 & 1969 \\
 & multimodal & 0/24 & 0/24 & 15/24 & 24/24 & 24/63 & 15 & 1969 \\
\midrule
\multirow{3}{*}{BFWS} & text & 0/24 & 0/24 & 17/24 & 24/24 & 24/24 & 0 & 974 \\
 & visual & 0/24 & 0/24 & 17/24 & 24/24 & 24/24 & 0 & 974 \\
 & multimodal & 0/24 & 0/24 & 17/24 & 24/24 & 24/24 & 0 & 974 \\
\midrule
\multirow{3}{*}{Weighted A* ($w{=}3$)} & text & 0/24 & 20/24 & 24/24 & 24/24 & 4/609 & 179 & 694 \\
 & visual & 0/24 & 18/24 & 24/24 & 24/24 & 6/499 & 143 & 694 \\
 & multimodal & 0/24 & 22/24 & 24/24 & 24/24 & 2/666 & 196 & 694 \\
\midrule
\multirow{3}{*}{Greedy best-first} & text & 0/24 & 22/24 & 24/24 & 24/24 & 2/614 & 179 & 673 \\
 & visual & 0/24 & 22/24 & 24/24 & 24/24 & 2/643 & 189 & 673 \\
 & multimodal & 0/24 & 21/24 & 24/24 & 24/24 & 3/617 & 179 & 673 \\
\bottomrule
\end{tabular}
\end{table}

\section{Results Detail: Primary-Matrix Paired Contrasts}
\label{app:results-contrasts}

% Paragraph purpose: Say what this appendix holds and why a reader would look here.
This appendix gives the paired whole-problem contrasts of the expanded baseline, for readers asking how the learned gain compares with each control.

% Paragraph purpose: Process SFT gains over the base only in the additive cells and never beats random-valid, whose tie with the exact reference there means those cells cannot register choice.
Process SFT gains over the base only in the additive cells, and it never beats random-valid (Table~\ref{tab:results-contrasts}). Against the base it gains $+0.833$ under weighted A* and $+0.903$ under greedy best-first, and nothing under BFS or BFWS. Against random-valid it loses in every algorithm column and wins none of the 72 problem-modality units in any column. The additive gains do not reflect choice, because random-valid itself ties exact execution in both additive columns, so those cells cannot measure any advantage beyond valid bookkeeping. The development-panel BFWS positive does not reappear on this suite. These contrasts are tiny-subgroup descriptive bounds over 24 problems and one training seed, so they support no modality or algorithm superiority claim beyond the executed panel. The complete six-measure matrix for the frozen final protocol, including per-cell Algorithm Invariant compliance and charged-budget columns, belongs to the unexecuted frozen 45-task final evaluation.
% Evidence: docs/experiments/expanded-study/synthesis-v1/baseline-contrasts.csv process_sft_minus_pretrained_base rows (modality means +0.833 w3, +0.903 greedy).
% TODO: same pending evaluation as above; the final protocol's full
% six-measure matrix is part of that unexecuted slot.

\begin{table}[t]
\caption{\textbf{Process SFT gains over the base only in the additive cells and trails random-valid in every column.} Descriptive paired success contrasts. Each cell gives the point difference and wins/losses/ties over problem-modality units.}
\label{tab:results-contrasts}
\centering\footnotesize
\setlength{\tabcolsep}{2pt}
\begin{tabular}{lcccc}
\toprule[0.4mm]
\rowcolor{mygray}\textbf{Contrast} & \textbf{BFS} & \textbf{BFWS} & \textbf{Weighted A* ($w{=}3$)} & \textbf{Greedy} \\
\midrule
SFT $-$ base & $+0.000$ (0/0/72) & $+0.000$ (0/0/72) & $+0.833$ (60/0/12) & $+0.903$ (65/0/7) \\
SFT $-$ random-valid & $-0.625$ (0/45/27) & $-0.708$ (0/51/21) & $-0.167$ (0/12/60) & $-0.097$ (0/7/65) \\
SFT $-$ exact & $-1.000$ (0/72/0) & $-1.000$ (0/72/0) & $-0.167$ (0/12/60) & $-0.097$ (0/7/65) \\
Random-valid $-$ exact & $-0.375$ (0/27/45) & $-0.292$ (0/21/51) & $+0.000$ (0/0/72) & $+0.000$ (0/0/72) \\
\bottomrule
\end{tabular}
\end{table}

\section{Results Detail: Secondary Panels}
\label{app:results-secondary}
\label{sec:results-secondary}

% Paragraph purpose: Say what this appendix holds and why a reader would look here.
This appendix reports the secondary enumeration-contract panels, which Section~\ref{sec:results-primary} points to and which sit outside the choice-quality claim.

% Paragraph purpose: No executed enumeration-contract panel shows a modality effect, and parity limits would bound one if it appeared.
\paragraph{Modality.} No executed enumeration-contract panel shows a modality effect. The per-modality rows of Table~\ref{tab:results-primary-matrix} have the same profile in all three adapters: zero learned success under BFS and BFWS, and learned success below the saturating controls in the additive cells. The earlier matched-modality study, a 3-problem panel on its own ledger that is never pooled with the expanded baseline, agrees. Its process-SFT success is 5/12 in text-state, 6/12 in visual-state, and 6/12 in multimodal-state, against 10/12 for random-valid in each modality, 0/12 for the base, and 12/12 for exact execution. A difference, had one appeared, could not be read as an intrinsic modality effect, because semantic parity across adapters is not established. The unlabelled 128px rendered views can lose symbolic identities, and the static and goal pages still contain text in every modality.
% Evidence: docs/experiments/expanded-study/synthesis-v1/README.md "Historical v5, kept separate" and claims historical-separation
% (v5-final-evaluation.json; analysis-v5/analysis.json): SFT 5/12, 6/12, 6/12; base 0/12; random-valid 10/12; exact 12/12.

% Paragraph purpose: The learned additive adapters succeed on most semantics-preserving perturbations, shifted initial states hurt most, and the saturated controls remain bounds.
\paragraph{Generalization and robustness v2.} The learned additive best-first adapters succeed on most semantics-preserving perturbations, and shifted initial states hurt most. The reduced v2 protocol admitted a uniform cheapest-prefix sample of 5 variants per family over 5 perturbation families (25 of 93 eligible variants, membership frozen before execution). Learned adapters succeeded on 128/150 episodes, against 0/150 for the pretrained base, whose every episode is a one-call invalid termination (150 decisions total). By family, learned success was 29/30 on object-renaming, 27/30 on name-compression, 27/30 on render-restyle, 26/30 on scale-up, and 19/30 on shifted-init. The modality split was text 44/50, multimodal 43/50, and visual 41/50, and the algorithm split was greedy 67/75 and weighted A* 61/75. All 22 learned failures are single-invalid-operation terminations, 11 of them in shifted-init. 

The controls saturate (random-valid 750/750, exact 150/150), so they remain bounds rather than a measure of learned ability. The lossy name-compression stratum is reported separately and never pooled. All 5 admitted variants are text-lossy, 3 multimodal-lossy, and 0 visual-lossy, and learned success on that stratum is text 9/10, visual 8/10, and multimodal 10/10. Family strata hold 5 problems each, so family-level rates are descriptive only. The full-scope v1 generalization branch was withdrawn after its terminal VALID\_STOP receipt (closed 2026-09-21) and will not be executed.
% Evidence: docs/experiments/expanded-study/synthesis-v1/README.md claim generalization-v2-headline (128/150, 0/150, per-family and
% per-modality splits); docs/experiments/expanded-study/issue-124-closeout.md (algorithm split, lossy stratum, 750/750, 150/150);
% outputs/expanded-study/v1/generalization-robustness/evaluation.json.
% Status 2026-09-21: issues #96/#98 CLOSED (supervisor decision, v1 scope withdrawn); the executed evidence is the reduced
% v2 scope only. Prose states the withdrawal.

% Paragraph purpose: DAgger does not improve BFS over exposure-matched continued SFT, and both stay far below random-valid.
\paragraph{DAgger.} DAgger \citep{ross2011dagger} does not improve BFS execution over an exposure-matched continued-SFT control, and both stay far below random-valid (Table~\ref{tab:results-dagger}). The comparison ran two DAgger iterations on 3 development and 24 unseen modality-task pairs. On the unseen panel, DAgger and continued SFT each achieve 1/72 invariant-valid successes, with one DAgger win, one loss, and 70 ties over whole problems. Random-valid reaches 45/72 and exact execution 72/72, while the original process-SFT checkpoint achieves 0/72. DAgger also has the higher unseen invalid-operation rate of the two iterated arms (0.316 versus 0.175). Final cumulative unique training corrections were text 201, visual 203, and multimodal 181, accumulated over 12 training cells of 512 records and 16 optimizer updates per cell-iteration.
% Evidence: docs/experiments/expanded-study/synthesis-v1/dagger-summary.csv and README.md section 3 (DAgger, #78-#84).

\begin{table}[t]
\caption{\textbf{DAgger does not improve BFS over exposure-matched continued SFT.} Development-stage success and unseen invalid-operation rate for the iterated arms, the original process-SFT checkpoint, and both references.}
\label{tab:results-dagger}
\centering\footnotesize
\begin{tabular}{lccc}
\toprule[0.4mm]
\rowcolor{mygray}\textbf{Arm} & \textbf{Dev success /9} & \textbf{Unseen success /72} & \textbf{Unseen invalid rate} \\
\midrule
DAgger iteration 2 & 3/9 & 1/72 & 0.316 \\
Continued SFT iteration 2 & 2/9 & 1/72 & 0.175 \\
Original process SFT & 1/9 & 0/72 & 0.364 \\
\midrule
Random-valid & 6/9 & 45/72 & 0.000 \\
Exact reference & 9/9 & 72/72 & 0.000 \\
\bottomrule
\end{tabular}
% Evidence: docs/experiments/expanded-study/synthesis-v1/dagger-summary.csv (dev /9, unseen /72, unseen invalid rate per arm).
\end{table}

% Paragraph purpose: The policy depends on trusted successors; model-generated successors collapse success, and their failures lie in state identity and effects rather than schema.
\paragraph{Model-generated successors.} The learned policy depends on the runtime's trusted successors. When the model generates dynamic-state successors itself, goal success collapses from 45/45 to 1/45 on the evaluated panel, and the single success is the one multimodal unseen cell. The trusted arm made zero trusted-state substitutions. The generated successors fail mostly on content rather than format. The trained successor model is exactly correct on 65/109 retained raw predictions, with 93 schema-valid but only 66 effect-valid, and downstream failure kinds are dominated by exact state-identity and effect failures (27 effect, 16 schema, 1 static-context). Collection-time verification of the untrained exact-successor accepted 119/1,536 predictions (text 43, visual 52, multimodal 24). Coverage is partial by design. The frozen outcome-blind admission accepted a 12-task fallback panel after the full 24-task unseen panel was priced at 147.66 GPU-hours against the 64 GPU-hour branch cap, so full unseen coverage remains missing.
% Evidence: docs/experiments/expanded-study/synthesis-v1/successor-summary.csv and README.md successor section (45/45, 1/45, 65/109,
% 93 / 66, 119/1,536); failure kinds per failure-mechanism-analysis.md taxonomy (effect 27, schema 16, static_context 1).

% Paragraph purpose: Training order does not interact materially with modality, although ceiling and floor effects bound what the panel can show.
\paragraph{Curriculum by modality.} Training order shows no material interaction with modality, although saturated controls leave little room for one to appear. Nine training cells crossed three modalities with staged, shuffled, and mixed ordering, all under greedy additive best-first and trained fresh from the same base on identical record sets, and were evaluated on 27 unseen whole problems. Successes were 24/24/25 (staged/shuffled/mixed) in text-state, 24/27/26 in visual-state, and 25/27/27 in multimodal-state. A sequential-SFT comparator reached 24/27, 25/27, and 24/27 in text, visual, and multimodal cells respectively. All four paired interaction intervals include zero: $-0.111$ $[-0.259, +0.037]$ and $-0.074$ $[-0.259, +0.074]$ for the visual and multimodal staged-minus-shuffled contrasts against text, and $-0.074$ $[-0.259, +0.111]$ and $-0.037$ $[-0.185, +0.111]$ for the corresponding mixed-minus-shuffled contrasts. The base stayed at 0/27, and random-valid and exact execution saturated at 27/27 in every modality, so ceiling and floor effects bound the observable differences.
% Evidence: docs/experiments/expanded-study/synthesis-v1/curriculum-summary.csv (successes of 27 per cell and control) and
% synthesis-v1/analysis.json curriculum interaction intervals (README.md curriculum section).

% Paragraph purpose: No planning adapter changes zero-shot reasoning, math, or coding accuracy detectably, the null is not equivalence, and the leakage screen finds no content overlap.
\paragraph{Transfer.} No planning adapter changes zero-shot accuracy detectably on reasoning, math, or coding benchmarks. All 12 verified planning adapters and the base were evaluated zero-shot on frozen subsets of FOLIO \citep{han2022folio} (200 items), GSM8K \citep{cobbe2021gsm8k} (200), and HumanEval \citep{chen2021humaneval} (164), with per-adapter accuracies in Table~\ref{tab:results-transfer}. McNemar exact comparisons \citep{mcnemar1947note} of each adapter against the per-example base predictions, Holm-corrected across all 36 comparisons \citep{holm1979simple}, leave none significant. The minimum adjusted $p$ is 0.8156, and the smallest raw $p$ is 0.0227 (FOLIO with the text width-based adapter, at $+12$ examples). This null does not establish equivalence. A leakage screen over the frozen 8-gram window finds 7/564 benchmark items sharing only trivial numeric 8-grams with the planning corpus and zero content overlap.
% Evidence: docs/experiments/expanded-study/synthesis-v1/README.md claim transfer-null (transfer-paired.csv; transfer-paired-analysis.json;
% transfer/leakage.json): 0/36 after Holm, min adjusted p 0.8156, smallest raw p 0.0227, 7/564 items.

\begin{table}[t]
\caption{\textbf{No planning adapter differs detectably from the base on zero-shot transfer.} Descriptive accuracy on frozen FOLIO, GSM8K, and HumanEval subsets. t/v/m denote the adapter training modality. No comparison with the base survives Holm correction.}
\label{tab:results-transfer}
\centering\footnotesize
\setlength{\tabcolsep}{5pt}
\begin{tabular}{@{}llccc@{}}
\toprule[0.4mm]
\rowcolor{mygray}\textbf{Adapter} & \textbf{Modality} & \textbf{FOLIO} & \textbf{GSM8K} & \textbf{HumanEval} \\
\midrule
Base & --- & 0.600 & 0.955 & 0.780 \\
\midrule
\texttt{bfs} & t & 0.600 & 0.950 & 0.793 \\
 & v & 0.605 & 0.945 & 0.799 \\
 & m & 0.605 & 0.965 & 0.780 \\
\midrule
\texttt{iw} & t & 0.660 & 0.960 & 0.768 \\
 & v & 0.645 & 0.965 & 0.762 \\
 & m & 0.645 & 0.960 & 0.756 \\
\midrule
\texttt{astar\_w3} & t & 0.590 & 0.950 & 0.817 \\
 & v & 0.575 & 0.965 & 0.799 \\
 & m & 0.590 & 0.945 & 0.774 \\
\midrule
\texttt{astar\_greedy} & t & 0.590 & 0.950 & 0.799 \\
 & v & 0.590 & 0.955 & 0.811 \\
 & m & 0.590 & 0.950 & 0.793 \\
\bottomrule
\end{tabular}
% Evidence: docs/experiments/expanded-study/synthesis-v1/transfer-accuracy.csv (per-adapter accuracy per benchmark) and transfer-paired.csv.
\end{table}

% Paragraph purpose: A second backbone reproduces the BFS null: neither the base nor process SFT reaches a goal.
\paragraph{Second backbone.} A second backbone, InternVL3.5-8B \citep{wang2025internvl35}, reproduces the enumeration-contract pattern across two reduced-scope cells. Each cell used a frozen 12-task key-cell panel selected by lowest reference decision count before any model outcome, with one seed-17 training run per modality cell. Under BFS (protocol v2) neither arm reached a goal. The pretrained base scored 0/36 at 1 decision per episode (36 total), and process SFT scored 0/36 at 1 to 4 decisions (59 total), while the reused BFS controls scored random-valid 54/72 (18/24 per modality, oracle-assisted) and exact 72/72.
% Evidence: outputs/expanded-study/v1/second-backbone-v2/evaluation/{evidence,analysis}.json (protocol expanded-second-backbone-v2;
% comparator sources are bfs-* baseline episodes; 12 tasks; complete_coverage true).
Under greedy additive best-first (protocol v3, stress test R3), the same backbone, recipe, and panel shape give high learned success that still trails random-valid. Process SFT succeeds on 11/12 text-state, 10/12 visual-state, and 10/12 multimodal-state problems (31/36), while the base scores 0/36 with one invalid call per episode. Process-SFT invalid-operation rates are 0.00546, 0.01220, and 0.01258, and decision usage is 183, 164, and 159 of a 390-call allowance (text, visual, multimodal). Random-valid and exact reference both score 24/24 per modality. These control counts hold two entries per panel task, and both entries of a task come from its one reused, hash-pinned greedy baseline control episode per condition and modality. Paired whole-problem contrasts (seed 1729, 10,000 resamples) give SFT minus base $+0.917$ $[+0.750, +1.000]$ in text-state and $+0.833$ $[+0.583, +1.000]$ in both visual-state and multimodal-state. SFT minus random-valid is $-0.083$ $[-0.250, +0.000]$ in text-state, and $-0.167$ $[-0.417, +0.000]$ in both visual-state and multimodal-state. Coverage is complete, with 72/72 model episodes and 72/72 comparators replayed.
% Evidence: outputs/expanded-study/v1/second-backbone-v3/evaluation/evidence.json (by_condition: process_sft 31/36, base 0/36 with
% 36 invalid over 36 decisions; comparator sources best_first_add_greedy-*) and analysis.json (by_modality successes, invalid_operation_rates
% 0.00546448 / 0.01219512 / 0.01257862, decision_usage 183 / 164 / 159 of 390; contrasts; bootstrap seed 1729, 10000).
% Source conflict: issue-130-closeout.md summarizes rates as "0.5-1.2%" and usage as "164-183 of 390"; neither covers the
% multimodal cell (0.01258, 159). The per-cell artifact values are used. The artifact reduced_scope string reads "bfs decision
% count"; issue-130-closeout.md records this as an inherited template string and states selection used additive reference decision
% counts (task lists of v2 and v3 differ), so the prose says "reference decision count" without naming the algorithm.
% Control denominators: analysis.json by_modality lists episodes 12 and random_valid_successes 24 / exact_reference_successes 24;
% evidence.json comparator_provenance.sources lists 72 files, while by_condition counts 72 episodes per control and comparator_episodes
% 144 against expected 72. Source conflict: the artifacts count each reused control episode twice; the prose states the
% two-entries-per-task basis and does not print a deduplicated count.
% Paragraph purpose: Read with the identity audit, the InternVL cells fit the enumeration-contract pattern: the greedy cell shows contract fluency, not choice, and the BFS null is consistent with an algorithm rather than a backbone difference.
Read with the identity audit, the two InternVL cells fit the enumeration-contract pattern seen on Qwen3-VL. Learned success under the enumeration contract appears only in the additive cells, where the identity audit shows no decision to make (Section~\ref{sec:results-identity}), and the v3 cell replicates that contract fluency (31/36 versus 0/36 base, below random-valid 24/24 per modality). The earlier total null on this backbone is therefore consistent with an algorithm (choice-requirement) difference rather than a backbone null. The cross-backbone contrast against the pinned Qwen3-VL process-SFT rate is descriptive rather than paired, because the InternVL panel is the key-cell subset while the Qwen rate is the pinned full-panel aggregate. It is $+0.000$ $[-0.167, +0.083]$ in text-state, $-0.083$ $[-0.333, +0.083]$ in visual-state, and $-0.042$ $[-0.292, +0.125]$ in multimodal-state. InternVL3.5-8B shares the Qwen3-8B LLM family \citep{yang2025qwen3}, so the comparison differs in the vision tower, connector, image tokenization, and multimodal recipe, and it is reported as an architecture-difference limitation. The full-scope v1 second-backbone branch was withdrawn after its terminal VALID\_STOP receipt (closed 2026-09-21) and will not be executed, and its gate is retained as history.
% Evidence: second-backbone-v3/evaluation/analysis.json cross_backbone (Qwen rates 0.917 / 0.917 / 0.875; method string states
% descriptive, not paired; architecture_note "reported as an architecture-difference limitation"); issue-130-closeout.md R3 (algorithm-confound reading).
% Status 2026-09-21: issues #102/#103 CLOSED (v1 scope withdrawn); prose states the withdrawal.

\section{Results Detail: Failure-Mechanism Calibration}
\label{app:results-failures}
\label{sec:results-failures}

% Paragraph purpose: Say what this appendix holds and why a reader would look here.
This appendix localizes where each arm class fails, for readers asking why learned success is absent under BFS and BFWS.

% Paragraph purpose: The three arm classes fail in distinct ways: the base mostly cannot emit the format, learned arms mostly break the search invariant early, and the oracle control runs out of budget late.
Each arm class fails in its own way. The base mostly cannot produce the output format, learned arms mostly emit well-formed operations that break the search invariant early, and the oracle control runs out of budget late. These profiles come from a CPU-only calibration, frozen before any outcome was computed, that mined 2,910 unique replay-verified episodes across the five executed model branches and classified all 905 failures with zero unclassified. The analysis is descriptive. The base policy fails in 447 of 447 episodes, always on the first decision. Of these, 379 emit a malformed output that never reaches the operation checker, and 67 issue rejected BFS frontier operations. Learned arms fail in 407 of 906 episodes. Their dominant failure kind is \texttt{other\_invariant\_violation} with 229 occurrences (56.3\% of learned failures), namely parsed, source-correct, applicable operations that the runtime rejects because they violate the declared search invariant. The remaining learned failures are 72 malformed outputs, 56 inapplicable grounded actions, 27 effect failures, 16 schema failures, and smaller kinds. First failures cluster early, at decision indices 0, 1, and 2 with counts 137, 88, and 46. The oracle control fails 51 times, all late and all by expansion-budget exhaustion (buckets 5--8, 9--16, 17--32, and 33+ hold 9, 9, 3, and 30). Table~\ref{tab:appendix-rejection} gives the rejection taxonomy.
% Evidence: docs/experiments/expanded-study/failure-mechanism-analysis.md (arm-level profiles, taxonomy table, first-failure index
% counts 137/88/46, oracle budget buckets 9/9/3/30; 144/144, 155/162, 56 of 546 in the family breakdown).

% Paragraph purpose: The failures concentrate where the primary matrix shows no learned success, and process supervision leaves an invariant problem rather than a format problem.
The failures concentrate where the primary matrix shows no learned success (Section~\ref{sec:results-primary}). Every baseline process-SFT episode under BFS or BFWS fails (144/144), and the DAgger-iteration BFS arms fail 155/162, while the learned additive best-first families fail 56 of 546 episodes (10.3\%). Process supervision thus leaves mainly an invariant problem rather than a format problem, since most learned failures are rejected invariant violations and malformed outputs are a minority. Read with the horizon profile, the learned policy fails early on invariant grounds, whereas the oracle control, which always emits valid operations, fails only when its expansion budget runs out late in the episode. Like the rest of the calibration, this reading is descriptive.
% (moved out of prose; kept in the preceding paragraph and Table tab:appendix-rejection: 229/407 invariant violations, 72/407 malformed outputs)
% Evidence: docs/experiments/expanded-study/failure-mechanism-analysis.md, "Learned arms: 407 failures" and taxonomy rows
% learned|other_invariant_violation|229 and learned|malformed_output|72 (Table tab:appendix-rejection source).

\section{Results Detail: Observation-Corruption Channel Tests}
\label{app:results-corruption}
\label{sec:results-corruption}

% Paragraph purpose: Say what this appendix holds and why a reader would look here.
This appendix reports the observation-corruption tests, which show which observation channel the trained enumeration-contract controller relies on.

% Paragraph purpose: Corrupting the text collapses the frozen enumeration-contract controller's success, while corrupting the images barely moves it.
Corrupting the text collapses the trained enumeration-contract controller's success, while corrupting the images barely moves it. The suite evaluated frozen checkpoints, with no retraining, on a prospectively frozen set of $k{=}9$ admitted tasks under four deterministic corruption families applied at observation time. Clean learned success on these tasks is 0.917 (text-state 0.944, visual-state 0.889, multimodal-state 0.889, with 9 tasks $\times$ 2 algorithms per cell). Copyability-preserving text shuffling lowers success by up to $-0.944$, and the multimodal arm retains at most 0.167 success under shuffled text. Masking instead destroys the output contract. Blanking or degrading the rendered state costs at most $-0.111$, and every visual interval includes zero, so visual-channel corruption is not material (Table~\ref{tab:results-corruption}). The multimodal channel-isolation contrast, paired by task and algorithm ($n{=}18$), measures the gap directly. Text-channel corruption costs $-0.806$ against $-0.111$ for visual-channel corruption, giving text minus visual $-0.694$ with 95\% interval $[-0.889, -0.472]$, which is material under the frozen rule. This contrast pools the output-contract-destruction (masked) regime. The controls behave as the contract predicts. The corrupted base stays at floor (0/144, all first-decision invalid), exact execution remains at 18/18, and random-valid succeeds 72/72 on the new frozen seeds, so per the saturation rule no structural-advantage claim is made.
% (moved out of prose; kept in Table tab:results-corruption) intervals $[-1.000, -0.833]$ for the text-shuffled text arm and $[-0.278, +0.000]$ for the visual multimodal arm.
% Evidence: outputs/expanded-study/v1/modality-stress/analysis.json contrasts.degradation_by_family / by family x modality groups
% (learned_clean_success, learned_corrupted_success, mean_degradation, degradation_interval_95, material) and the channel-isolation block
% (text_channel_mean_degradation -0.8056, visual -0.1111, mean_text_minus_visual_degradation -0.6944); docs/experiments/expanded-study/issue-126-closeout.md.

\begin{table}[t]
\caption{\textbf{Text corruption collapses the success of frozen checkpoints and visual corruption does not.} Paired differences compare corrupted with clean runs. \cmark\ marks a material interval-based verdict and $=$ a non-material verdict.}
\label{tab:results-corruption}
\centering\footnotesize
\setlength{\tabcolsep}{4pt}
\begin{tabular}{@{}ll|cc|cc|c@{}}
\toprule[0.4mm]
\rowcolor{mygray}\textbf{Corruption} & \textbf{Arm} & \textbf{Corrupted} & \textbf{Clean} & \textbf{$\Delta$} & \textbf{95\% interval} & \textbf{Material} \\
\midrule
Text-masked & text & 0.000 & 0.944 & $-0.944$ & $[-1.000, -0.833]$ & \cmark \\
Text-masked & multimodal & 0.000 & 0.889 & $-0.889$ & $[-1.000, -0.722]$ & \cmark \\
\midrule
Text-shuffled & text & 0.000 & 0.944 & $-0.944$ & $[-1.000, -0.833]$ & \cmark \\
Text-shuffled & multimodal & 0.167 & 0.889 & $-0.722$ & $[-0.889, -0.500]$ & \cmark \\
\midrule
Visual-blank & visual & 0.833 & 0.889 & $-0.056$ & $[-0.167, +0.000]$ & $=$ \\
Visual-blank & multimodal & 0.778 & 0.889 & $-0.111$ & $[-0.278, +0.000]$ & $=$ \\
\midrule
Visual-degraded & visual & 0.833 & 0.889 & $-0.056$ & $[-0.167, +0.000]$ & $=$ \\
Visual-degraded & multimodal & 0.778 & 0.889 & $-0.111$ & $[-0.278, +0.000]$ & $=$ \\
\bottomrule
\end{tabular}
% Evidence: outputs/expanded-study/v1/modality-stress/analysis.json contrasts (per corruption family x arm: corrupted, clean, delta,
% 95% interval, material); values as tabulated in docs/experiments/expanded-study/issue-126-closeout.md.
\end{table}

% Paragraph purpose: Masking and shuffling the text fail for different reasons: masking destroys the output contract, shuffling removes information while valid operations keep flowing.
Masking and shuffling the text fail for different reasons, which the success deltas pool together. A decomposition of the 72 published text-corruption episodes (stress test R2) separates the two regimes. Under text masking, every episode of the text-state and multimodal-state cells is all-invalid (18/18 in each, 18 invalid operations over 18 decisions). Each emitted operation fails the schema, so the policy never enters the search. This is output-contract destruction, not information loss. Under copyability-preserving shuffling, by contrast, valid operations keep flowing while the search fails. In the text-state cell 16/18 episodes are partial-valid, 2/18 are all-invalid, and none reaches the goal (18 invalid operations over 49 decisions). In the multimodal-state cell 15/18 are partial-valid, none is all-invalid, and 3/18 reach the goal (15 invalid operations over 105 decisions). This partial-valid pattern is the information-loss signature. A new cross cell shows the same format failure on an image-state arm. Text masking applied to the frozen visual-state arm scores 0/18 with every episode all-invalid, so masking the text scaffold destroys the output contract even when images are the nominal state channel.
% Evidence: outputs/native-arms/v1/r2-decomposition.json (per-cell all_invalid / partial_valid / goal_reached / decisions / invalid);
% cross cell: outputs/native-arms/v1/stress-evaluation.json textmask|text-masked|visual-state cells and issue-130-closeout.md ("0/18, all-invalid").
% Source conflict: issue-130-closeout.md reports shuffled "34/36 partial_valid"; the JSON records partial_valid 16 + 15 with 3 goal-reached
% and 2 all-invalid episodes. Per-cell JSON counts are used; no pooled shuffled count is printed.

% Paragraph purpose: The text channel carries decision-relevant information only in the shuffled sense, because masking mostly destroys the output format; this is a channel finding, not a modality claim.
The language-channel claim survives only in its shuffled form. The masked catastrophe in Table~\ref{tab:results-corruption} is mostly a loss of the copyable output format rather than a loss of state information. Shuffling the text leaves operations copyable and still stops the search, so under the frozen observation contract the textual serialization carries decision-relevant information for these adapters. Removing or degrading the visual channel barely moves success, including on the visual-state arm whose only rendered-semantics channel was corrupted. Appendices~\ref{app:results-menu} and~\ref{app:results-native} probe the grounded candidate menu directly. This is a channel finding about the trained controller under the declared adapters, not evidence about intrinsic modality value.
% Evidence: docs/experiments/native-arms/issue-130-closeout.md R2 ("survives only in the shuffled form; the masked catastrophe was mostly copyability").
% The published #126 numbers in the table above are unchanged by the decomposition.

\section{Results Detail: Menu-Manipulation Receipts}
\label{app:results-menu}
\label{sec:results-menu}

% Paragraph purpose: Say what this appendix holds and why a reader would look here.
This appendix reports the menu-manipulation receipts, for readers asking whether the trained enumeration-contract policy reads its menu's content or only its surface.

% Paragraph purpose: Menu manipulation asks whether the trained policy responds to the content of its grounded menu, using two families on the frozen native adapters.
Menu manipulation (stress test R1) asks whether the trained policy responds to the content of its grounded menu. It runs on the frozen image-only adapters of both native arms (Appendix~\ref{app:results-native}) with no retraining. Menu permutation reorders the candidate menu, and distractor injection adds $k{=}3$ schema-valid but inapplicable actions to it. Table~\ref{tab:results-menu} reports the 72 episodes, all independently replayed within the 170/170 stress-evaluation episodes.
% Evidence: docs/experiments/native-arms/issue-130-closeout.md R1; outputs/native-arms/v1/stress-evaluation.json (episodes_replayed 170);
% outputs/native-arms/v1/evaluation/episodes/menu/. k=3 per the #131 brief vocabulary contract.

\begin{table}[t]
\caption{\textbf{Reordering the menu changes nothing, while injected inapplicable actions end every episode.} Menu manipulation on the frozen native arms. Invalid operations count emitted operations, and distractor picks count decisions with injected distractors.}
\label{tab:results-menu}
\centering\footnotesize
\setlength{\tabcolsep}{4pt}
\begin{tabular}{lccc}
\toprule[0.4mm]
\rowcolor{mygray}\textbf{Family} & \textbf{Success} & \textbf{Invalid ops} & \textbf{Distractor picks} \\
\midrule
Menu permutation & 36/36 & 0/512 & n/a \\
Distractor injection & 0/36 & 36/60 & 36/60 (0.6) \\
\bottomrule
\end{tabular}
\end{table}
% Evidence: issue-130-closeout.md R1 table; stress-evaluation.json distractor block (decisions_with_injection 60, distractor_picks 36, pick_rate 0.6).
% stress-evaluation.json: 9 episodes and 15 decisions per distractor-injection cell, 60 decisions over 36 episodes; episodes terminate
% on the first invalid operation. 16/36 episodes open with one or two accepted operations (8 with two decisions, 8 with three).
% Source conflict: docs/experiments/native-arms/issue-130-closeout.md L13 states "the first emitted operation of every episode is a
% distractor" (and L16 "every episode's first operation picks an injected ... action"). stress-evaluation.json and the episode
% records contradict it; e.g. visual-nomem-state/expanded-final__ferry-compact-915000/
% best_first_add_greedy-distractor-injection-learned_adapter.json.gz accepts events 0 and 1 (sail l1 l0, sail l1 l2) and rejects
% event 2 (debark c0 l1). The pinned JSON and episode records are used.
% The closeout's "rules out order-riding" and "rules out state grounding / reads menu CONTENT" rulings are not used.

% Paragraph purpose: The distractor receipt shows that the policy does not reject schema-valid inapplicable menu entries, but it cannot separate content reading from position habits; permutation is uninformative about order.
The distractor receipt shows that the policy does not reject schema-valid inapplicable menu entries, but it cannot tell content reading from position habits. On the frozen native adapters (stress test R1, 72 replayed episodes), injecting $k{=}3$ schema-valid but inapplicable actions collapses success to 0/36, with 36/60 invalid operations and a distractor pick-rate of 0.6 against a uniform-picking expectation of 0.592 to 0.625. Every episode ends on a distractor pick, 20/36 on the first decision. The pick-rate is therefore at chance, and these out-of-distribution distractors never appeared in training menus. Content-reading and position habits both predict a pick-rate near 0.6, so this receipt does not separate them, and the discriminating fixed-versus-random-position test remains unrun. Menu permutation leaves 36/36 successes and 0/512 invalid operations, which is uninformative about order under the submission-order invariance of Section~\ref{sec:results-identity}.
% Evidence: native adapters from issue #129; stress test R1 is issue #130 R1.
% Evidence: docs/experiments/native-arms/issue-130-closeout.md R1; outputs/native-arms/v1/stress-evaluation.json (distractor block:
% decisions_with_injection 60, distractor_picks 36, pick_rate 0.6). Chance rate 0.592 / 0.625 and the 20/36 first-decision count:
% post-hoc recomputation from outputs/native-arms/v1/evaluation/episodes/menu/*/*/*distractor-injection-learned_adapter.json.gz. Source conflict with issue-130-closeout.md L13/L16 recorded above.

\section{Results Detail: Image-Only Native Arms}
\label{app:results-native}
\label{sec:results-native}

% Paragraph purpose: Say what this appendix holds and why a reader would look here.
This appendix reports the image-only native arms, which test whether removing the textual scaffold changes the enumeration-contract results.

% Paragraph purpose: The native arms remove every textual search aid except the grounded candidate menu, and both passed the pre-registered smoke gate with full replay.
The native-arms follow-up (contract v2) removes every textual search aid except the grounded candidate menu. In its two image-only arms, \texttt{visual-nomem-state} (nomem) and \texttt{visual-seq-state} (seq), state facts arrive only as rendered images, and seq adds a window of $K{=}8$ history frames. Neither arm receives textual search memory, $g$/$h$/priority scalars, candidate scores, or membership flags, but the grounded, unscored candidate menu is still presented. Both arms passed a pre-registered post-training smoke gate with 62/62 schema-valid grounded calls against a threshold of 0.5. That gate comes from the v2 amendment, frozen before any trained-arm outcome, which replaced the v1 admission gate on pretrained-base format compliance (retained as history). Four adapters (two arms by two additive best-first algorithms) used the frozen recipe with 512 records, 16 updates, and seed 17. All 216/216 new bindings were recorded and independently replayed, together with 72/72 reused replay-verified comparators. The window spent 3.2086 of 12 GPU-hours.
% Evidence: native arms v1 gate #128, v2 #129.
% Evidence: docs/experiments/native-arms/issue-129-closeout.md (Delivered, Interpretation: K=8 window, smoke gate 62/62, threshold 0.5,
% 216/216 + 72/72, 3.2086 / 12.0); docs/experiments/native-arms/native-arms-v2-amendment.md; docs/experiments/native-arms/issue-128-closeout.md (v1 gate history).

% Paragraph purpose: Without textual bookkeeping both arms still solve every clean task, but the saturated random-valid control leaves no room to call this a structural advantage.
Without textual bookkeeping both arms still solve every clean task, but a saturated control leaves no room to call this an advantage. Both arms nominally succeed on 18/18 with zero invalid operations, against 16/18 for the published frozen \texttt{visual-state} arm. The pretrained base scores 0/36 under both new contracts, and the exact reference scores 18/18. Random-valid is at ceiling on its five frozen seeds, so under the saturation rule no structural-advantage claim is made. Paired whole-problem contrasts ($n{=}18$, bootstrap seed 61813, 10,000 resamples, 95\%) give nomem minus \texttt{visual-state} $+0.111$ $[0.000, +0.278]$ and seq minus \texttt{visual-state} $+0.111$ $[0.000, +0.278]$. Neither is material, and neither establishes equivalence. The history window changes nothing on these tasks (seq minus nomem $0.000$ $[0.000, 0.000]$). Learned minus pretrained base is $+1.000$ $[+1.000, +1.000]$ for each arm, which is material.
% Evidence: issue-129-closeout.md Results table and clean-success paragraph; outputs/native-arms/v1/evaluation/{evaluation,analysis}.json (outcome PASS, 144 paired cells).

% Paragraph purpose: The pre-registered images-only hypothesis fails because corrupting every image of the seq arm changes nothing; the post-hoc menu reading stays bounded by the identity audit.
The pre-registered corruption payoff is a null, so the images-only hypothesis fails on this panel. Under seq the images are the only state and history channel, yet blanking every attached image gives a seq corrupted-minus-clean contrast of $0.000$ $[0.000, 0.000]$, and degrading every image gives $0.000$ $[0.000, 0.000]$. As a post-hoc interpretation, the grounded unscored candidate menu is the remaining channel that could carry the signal, since neither the rendered scenes nor the removed scalar scaffold is necessary once the menu is present. That reading has a limit. On these additive cells the identity audit (Section~\ref{sec:results-identity}) makes any listed candidate a correct emission, so the null is equally consistent with a policy that uses no state information.
% Evidence: issue-129-closeout.md Results (corruption rows) and Interpretation.

% Paragraph purpose: The native arms reach ceiling on the two strains where the full-scaffold adapters degraded, but the comparison is post-hoc and unpaired and cannot show that the removed text carried signal.
The native arms reach ceiling on the two generalization strains where the full-scaffold adapters degraded, although this comparison is post-hoc and unpaired. The degraded-strain test (stress test R4) evaluated the same frozen native adapters on the shifted-init and scale-up families. The learned adapters succeed on 10/10 shifted-init and 10/10 scale-up episodes for both arms, and the pretrained base scores 0/20 on each strain (0/40 in total). The published full-scaffold rates on the same variants are 19/30 for shifted-init and 26/30 for scale-up (Appendix~\ref{app:results-secondary}). The two rates come from different episode sets, with 10 visual-only episodes per arm against 30 full-scaffold episodes pooled over three modalities. Read with that caveat, the result is consistent with the reduced contract removing scalar operands whose emission the full-scaffold adapters failed. It does not measure whether that text carried signal, since any listed candidate succeeds on these cells.
% Evidence: outputs/native-arms/v1/stress-evaluation.json variants|* cells (8 cells of 10 episodes, each pooling learned and base
% episodes: 5 successes, 5 invalid); the episode store holds 40 learned_adapter and 40 pretrained_base variant files (80 total,
% matching Table tab:appendix-receipts "80 variant"); learned/base split per docs/experiments/native-arms/issue-130-closeout.md R4
% ("base 0/20" per strain).

\section{Results Detail: Choice-Frontier and Comparator-Zoo Receipts}
\label{app:results-cf}

% Paragraph purpose: State what the first choice-frontier measurement and the comparator zoo establish on the reused 9-task panel, and what they cannot.
This appendix documents the first measurement under the choice-frontier contract, on the reused 9-task panel, and the comparator zoo scored on the same panel. Both predate the validated ladder. They show that the contract separates random-valid from the exact reference, but they cannot calibrate choice quality.

% Paragraph purpose: Show that the first adapter's corpus reproduces the exact reference and contains real choices, then give the training, replay, and budget receipts.
The first adapter (``Learned adapter'' in Tables~\ref{tab:results-cf} and~\ref{tab:results-zoo}) was trained on a corpus that reproduces the exact reference and contains real choices. The corpus was derived on CPU by replaying the frozen exact traces through the new contract, and all 58 derived episodes matched their stored traces expansion for expansion. It covers 29 training tasks and holds 1,046 records (512 training plus 11 diagnostic per algorithm) with 3,544 scenes. Every corpus episode carries at least one decision whose menu offers two or more states, with up to 315 such decisions in one episode and menus of up to 222 states. Two adapters, one per additive algorithm, were trained with the frozen recipe at seed 17 (16 updates on 512 records each). The smoke gate accepted 40/40 model calls, all 144/144 panel episodes were independently replayed with complete coverage, and the final audit returned CHOICE\_SENSITIVE over those 144 replayed episodes. The window spent 1.069 of 12 GPU-hours (training 0.81, smoke 0.052, model evaluation 0.208), and corpus and controls ran on CPU.
% Evidence (final audit, merged from the former closing protocol paragraph): outputs/choice-frontier/v1/evaluation/final-audit.json (episodes_replayed 144, CHOICE_SENSITIVE, ok true).
% Evidence: issue-132-closeout.md Execution record and Budget accounting; outputs/choice-frontier/v1/preparation/visual-choice-frontier/{report,audit}.json
% (episodes_rederived 58, PASS; report.json expansion_gate "every derived episode matched its stored exact trace expansion for
% expansion", menu_size_range max 222; audit.json headroom_present true, headroom_decisions_per_episode max 315 at
% astar-pair-68102553...:best_first_add_w3); outputs/choice-frontier/v1/budget.json; train/smoke/evaluation split per the #131 brief ledger.
% Source conflict: issue-132-closeout.md L66 says the corpus was "derived independently of the stored exact traces" and gives
% "max menu 315"; report.json records replay with an expansion-for-expansion gate and a maximum menu of 222. The artifact wording is used.

\begin{table}[t]
\caption{\textbf{First choice-frontier measurement on the frozen panel.} Random-valid falls well below the exact reference, and the learned adapter emits no invalid operation. Cap denotes expansion-budget exhaustion and first-call an invalid first emission. These goal rates are not comparable with enumeration-contract performance.}
\label{tab:results-cf}
\centering\footnotesize
\setlength{\tabcolsep}{2.5pt}
\begin{tabular}{lccc}
\toprule[0.4mm]
\rowcolor{mygray}\textbf{Condition} & \textbf{Goal rate} & \textbf{Failures} & \textbf{Invalid ops} \\
\midrule
Exact reference & 18/18 & --- & 0 \\
Learned adapter & 7/18 & 11 cap & 0 \\
Pretrained base & 0/18 & 18 first-call & 18 \\
Random-valid (5 frozen seeds) & 32/90 & 58 cap & 0 \\
\bottomrule
\end{tabular}
\end{table}
% Evidence: outputs/choice-frontier/o4/metrics/all-episode-metrics.json records[multiplier=2].solves["2"] per arm (18, 7, 0, 32 of 18/18/18/90);
% per-cell outcomes in outputs/choice-frontier/v1/evaluation/{cells,analysis}.json cells.*; failure kinds per issue-132-closeout.md.

% Paragraph purpose: Show that the contract registers choice (controls separate) while the first adapter does not separate from random-valid, with the repeated-measures, one-task, and no-equivalence boundaries.
The first measurement shows that the contract registers choice, but not that the first adapter chooses better than chance. Random-valid minus the exact reference is $-0.644$ $[-0.856, -0.422]$, material, so the two controls separate under this contract. The first adapter minus the pretrained base is $+0.389$ $[+0.167, +0.611]$, material, and with zero invalid operations this is a gain in contract fluency. The first adapter minus random-valid is $+0.033$ $[-0.033, +0.133]$, not material. Even that small difference is weaker evidence than its interval suggests. The 18 cells are 9 tasks measured under two algorithms as repeated measures, and the whole contrast rests on one task (storage), the only one where random-valid's seed frequency lies strictly between 0 and 1 (it is 0.2). The frozen rule makes strata under eight cells descriptive only, so the contrast is descriptive. The evidence therefore supports no learning-to-plan claim for the first adapter, and no equivalence with random-valid, since no equivalence margin was pre-registered. All contrasts are paired percentile 95\% bootstrap intervals over the frozen resampling cells (seed 61813, 10,000 resamples).
% Evidence: outputs/choice-frontier/v1/evaluation/{identity-audit,analysis}.json (18 divergent pairs; contrasts, ci95, material flags);
% issue-132-closeout.md Erratum items 2, 3, and 5 (binding; no expansion-parity wording). Task-clustered intervals in Appendix app:results-cf.
% Source conflict: analysis.json descriptive_only false; the Erratum item 3 one-task rule governs the prose.
% Evidence (seed frequency 0.2 and percentile intervals, merged from the former closing protocol paragraph): outputs/choice-frontier/v1/evaluation/analysis.json; issue-132-closeout.md Erratum item 3.

% Paragraph purpose: Show that task-clustered intervals, which respect the repeated measures, leave both readings of the first measurement unchanged.
Clustering by task leaves both readings unchanged. Clustering is needed because the control expansion sequences coincide across the two algorithms on every panel task. With the same seed and resample count, the post-hoc task-clustered interval is $[0.000, +0.100]$ for the first adapter minus random-valid, which still reaches zero, and $[-0.889, -0.333]$ for random-valid minus the exact reference, which still excludes it.
% Evidence: issue-132-closeout.md Erratum items 2 and 3.

\paragraph{Comparator zoo.}
\label{sec:results-calibration}
% Paragraph purpose: Define the three rule selectors and state why the zoo comparison is exploratory and not information-matched.
The comparator zoo is an exploratory CPU-only comparison that places three deterministic rule selectors on the same measurement at the same 2$\times$ cap. \texttt{bfs-order} picks the minimum frontier $g$, then the earliest serial. Novelty-first selection picks the lowest width-1, then width-2, novelty over positive PDDL atoms. \texttt{worst-first} is a negative control. The comparison is exploratory because the metric set was pre-registered before comparator episodes, but the stored choice-frontier episodes enter post-hoc. It is also not information-matched. The selectors act on the additive frontier, not native BFS or BFWS, and read privileged scalars or PDDL atoms while the model sees only scenes, so no learned-versus-rule claim follows.
% Evidence: zoo is issue #133. docs/experiments/choice-frontier/issue-133-protocol.md (selector definitions, frozen 2x cap); issue
% Sino-Huang/multimodality_on_planning#133 closing comment (additive frontier, privileged g/h or PDDL atoms, 432/432 replayed).
% Selector definitions merged here from the former selector paragraph under Zoo receipts.
% Evidence: docs/experiments/choice-frontier/issue-133-protocol.md (selector definitions); #133 closing comment (M1 points, 15-test family);
% outputs/choice-frontier/o4/metrics/analysis.json arms.*.m1_auc.

\begin{table}[t]
\caption{\textbf{Exploratory comparator-zoo metrics at the shared cap.} No rule selector exceeds random-valid on M1, so no comparator lies between random-valid and the exact reference. M1 has task-clustered intervals, and M2 and M3 are solved-only. The rule selectors read privileged information unavailable to the model.}
\label{tab:results-zoo}
\centering\footnotesize
\setlength{\tabcolsep}{4pt}
\begin{tabular}{@{}lcccc@{}}
\toprule[0.4mm]
\rowcolor{mygray}\textbf{Arm} & \textbf{M1 [95\% CI]} & \textbf{2$\times$ solved} & \textbf{3$\times$} & \textbf{4$\times$} \\
\midrule
Learned adapter & .1042 [.0278, .2014] & 7/18 & n/a & n/a \\
Pretrained base & 0 [0, 0] & 0/18 & n/a & n/a \\
\midrule
Random-valid & .1389 [.0250, .2722] & 32/90 & 42/90 & 50/90 \\
\midrule
Exact reference & .8472 [.7917, .8750] & 18/18 & n/a & n/a \\
\midrule
\texttt{bfs-order} & .0694 [0, .1528] & 6/18 & 8/18 & 10/18 \\
Novelty-first selection & .1250 [0, .2778] & 6/18 & 8/18 & 12/18 \\
\texttt{worst-first} & .0556 [0, .1389] & 4/18 & 6/18 & 8/18 \\
\bottomrule
\end{tabular}
\vspace{4pt}
\begin{tabular}{@{}lcc@{}}
\toprule[0.4mm]
\rowcolor{mygray}\textbf{Arm} & \textbf{M2 $\rho$} & \textbf{M3 $\kappa$} \\
\midrule
Learned adapter & 1.432 & 1.025 \\
Pretrained base & undefined & undefined \\
\midrule
Random-valid & 1.407 & 1.000 \\
\midrule
Exact reference & 1.000 & 1.000 \\
\midrule
\texttt{bfs-order} & 1.613 & 1.000 \\
Novelty-first selection & 1.452 & 1.000 \\
\texttt{worst-first} & 1.549 & 1.050 \\
\bottomrule
\end{tabular}
\end{table}
% Evidence: #133 closing comment (M1 points, 2x/3x/4x counts, M2, M3); outputs/choice-frontier/o4/metrics/analysis.json arms.*
% (m1_task_cluster_95pct_ci_10000_seed133 endpoints, shown to four decimals from the stored floats; m2/m3 base null = undefined).

% Paragraph purpose: Show that the zoo cannot calibrate choice quality because no comparator lies between random-valid and the exact reference, with the Holm/Pareto outcome and its non-equivalence boundary.
The zoo bounds the first adapter only against rules weaker than random-valid, so it cannot calibrate choice quality. On M1 the first adapter falls between \texttt{bfs-order} and random-valid and below novelty-first selection (Table~\ref{tab:results-zoo}, Figure~\ref{fig:zoo-m1}), and no comparator lies between random-valid and the exact reference on this panel. The formal tests find no ordering either. Across the fifteen-test Holm family no comparison supports positive Pareto dominance for the first adapter, and three co-solved tests are undefined. Non-significance here is neither equivalence nor a general negative capability claim. The post-hoc label diagnostics and the one-task M5 subset below are descriptive only.
% (moved out of prose; kept in the M4 paragraph below and Table~tab:appendix-m4) The 17.4% / 26.9% / 54.5% M4 sentence was a duplicate of the M4 paragraph below.
% Evidence (R2): outputs/choice-frontier/o4/metrics/m4-task-cluster-intervals.json statistics.observed_minus_chance_agreement
% .task_cluster_95pct_ci_10000_seed133 [-0.1609, 0.0323] (E134-27).
% Evidence: #133 closing comment (15-test Holm, no positive learned Pareto dominance, three co-solved tests lack support; M5 one task);
% outputs/choice-frontier/o4/metrics/analysis.json m1_m2_m3_dominance_holm (inferential_positive_after_holm false), m5_informative_tasks_D.
% Evidence: outputs/choice-frontier/o4/metrics/analysis.json M4 (23/132 heap-head, 72/132 last label), post-hoc per-decision permutation; Appendix~\ref{app:results-cf}.

\paragraph{Zoo receipts.}
% Paragraph purpose: Show that zoo selectors face the same menus and budgets as the model, then give coverage, replay, and optimal-cost receipts.
The zoo reuses the model's episode conditions, so selectors and model face the same menus and budgets. Random-valid serves as the uniform-choice reference and the exact reference as the imitation target. Selector episodes replay the frozen model input, randomized menu, actions, admission, and budget, but they neither generate nor replay image view bindings. Because the selectors operate on the additive frontier, novelty-first selection does not reproduce BFWS pruning. Random-valid and the three selectors also ran at 3$\times$ and 4$\times$ caps, outside the primary metric. All 432/432 zoo bindings were observed and independently replayed, with 342 new CPU episodes, 90 original 2$\times$ random-valid controls reused, and none missing. The nine unit-cost optimal costs $C^*$ come from one offline planner computation per task with replay-verified plans: storage 3; elevators, ferry, and 15puzzle-compact 4; depot, towers-of-hanoi, and visitall 5; blocksworld and 15puzzle-expanded 6.
% Evidence: docs/experiments/choice-frontier/issue-133-protocol.md; #133 closing comment (coverage and C* values);
% outputs/choice-frontier/o4/metrics/analysis.json episode_accounting; outputs/choice-frontier/o4/metrics/optimal-costs.json.

% Paragraph purpose: Show that post-hoc label diagnostics find no heap-head agreement beyond chance and a last-label excess that does not establish a position habit.
Post-hoc diagnostics on the stored episodes find no heap-head agreement beyond chance and an excess of last-label choices that the data cannot attribute to a habit. The first adapter selects the reference's heap-head state on 23/132 multi-option decisions (17.4\%) against an on-policy chance rate of 26.9\%, and it emits the last menu label on 72/132 (54.5\%). These decisions cluster within 9 tasks. Under task clustering the agreement gap does not survive while the last-label excess does (Appendix~\ref{app:results-counterfactual}). The post-hoc choice-frontier diagnostic records a first-label rate of 2.3\% and a teacher-corpus last-label rate of 12.0\% (chance 10.4\%, 978 records). Menus are permuted per decision with a seed, and no last-label comparator was run, so these post-hoc rates do not establish a position habit.
% Evidence: #133 closing comment (M4, M5); outputs/choice-frontier/o4/metrics/stored-posthoc-analysis.json; analysis.json m5_informative_tasks_D
% and m5_D_descriptive (0.3125 / 0.075 / 0.625); issue-132-closeout.md Erratum item 4 (first label 2.3%, teacher 12.0%, chance 10.4%, 978 records).

% Paragraph purpose: Show that the M5 subset is a single task and therefore descriptive only.
M5 is too small to read beyond description. It restricts to tasks where random-valid solves at a 2$\times$ rate strictly between 0 and 1, and only one task qualifies (storage-compact-919000, random-valid rate .2). On it the first adapter's M1 is .3125, against .075 for random-valid and .625 for the exact reference.
% Evidence: #133 closing comment (M5); analysis.json m5_informative_tasks_D and m5_D_descriptive (0.3125 / 0.075 / 0.625) (copied from the M4 comment above).
\begin{figure}[t]
\centering
\includegraphics[width=0.75\columnwidth]{figures/fig_zoo_m1.pdf}
\caption{Only the exact reference stands apart on M1 (solve-versus-budget AUC at multipliers 1 to 2) on the frozen 9-task choice-frontier panel. The intervals of all non-exact arms overlap, and the Holm family licenses no ordering. Exact reference .847, random-valid uniform-choice reference .139, learned adapter .104 [.028, .201], novelty-first .125, bfs-order .069, worst-first .056, pretrained base 0. Under the enumeration contract, by contrast, the BFS/BFWS random-valid rows (15/24 and 17/24 versus exact 24/24) already registered choice.}
\label{fig:zoo-m1}
\end{figure}
% Evidence: caption fixed by local wave2-contract (figure contract item 2); values from outputs/choice-frontier/o4/metrics/analysis.json arms.*.m1_auc.

% Source conflict: analysis.json contrasts.learned_minus_random_valid.descriptive_only false; the Erratum item 3 one-task rule governs the prose.
% Evidence: outputs/choice-frontier/v1/evaluation/analysis.json; final-audit.json (episodes_replayed 144, CHOICE_SENSITIVE, ok true); issue-132-closeout.md Erratum item 3.
% (Former closing protocol paragraph merged into the contrast paragraph and the corpus paragraph above; its source-conflict and evidence comments are kept here and copied there.)

\section{Results Detail: Submission-Order Counterfactual and M4 Intervals}
\label{app:results-counterfactual}

% Paragraph purpose: State that the counterfactual locates the 48/48 identity in the executed serial rule, and that the section also holds the first adapter's task-clustered M4 intervals.
This appendix shows that the enumeration contract's 48/48 identity comes from the executed heap-serial rule. It also gives the task-clustered intervals behind the first adapter's post-hoc label diagnostics.

% Paragraph purpose: Show that the counterfactual changes only the heap-serial rule, then give the protocol timing and re-validation receipts.
The counterfactual changes only the rule that assigns heap serials and re-runs the identity audit on the same 48 additive pairs (24 tasks $\times$ 2 algorithms, seed 17). The executed rule assigns serials at expansion start from the sorted candidate order. The counterfactual assigns them at application time in submission order. The protocol was committed before any output of this CPU-only follow-up, and the audit first re-validated 288 stored episode counts and 96 raw-output sequences with no mismatch.
% Evidence: #134 protocol docs/experiments/choice-frontier/issue-134-protocol.md (commit de8f99a precedes outputs commit c020da4);
% outputs/native-arms/v1/identity-audit-submission-order.json pairs_checked 48 (E134-05), evaluation_seed 17 (E134-10), rule definitions
% (E134-03, E134-04), validation.stored_episode_counts_checked 288, .stored_raw_output_sequences_checked 96, .mismatches [] (E134-08). CPU only (LED-3).

\begin{table}[t]
\caption{\textbf{Heap-serial counterfactual identity audit.} Random-valid matches the exact reference on every additive pair only under the executed sorted-order rule. Each row compares random-valid and exact-reference choices on the same additive pairs.}
\label{tab:appendix-counterfactual}
\centering\footnotesize
\setlength{\tabcolsep}{4pt}
\begin{tabular}{lcc}
\toprule[0.4mm]
\rowcolor{mygray}\textbf{Heap-serial rule} & \textbf{Identical} & \textbf{Divergent} \\
\midrule
Sorted candidate order (executed) & 48 & 0 \\
Submission order (counterfactual) & 29 & 19 \\
\bottomrule
\end{tabular}
\end{table}
% Evidence: outputs/native-arms/v1/identity-audit-submission-order.json rules.sorted_candidate_order.pairs_identical 48 / .pairs_divergent 0;
% rules.submission_order.pairs_identical 29 / .pairs_divergent 19 (E134-03, E134-04); validation.exact_reference_rule_invariance.pairs_identical 48 (E134-09).

% Paragraph purpose: Show that the 48/48 identity is a property of the executed serial rule, with the failed pre-registered prediction, the divergence location, and the goal flips.
The 48/48 identity is a property of the executed serial rule, not of the enumeration contract alone. Under the submission-order rule 19/48 pairs diverge (Table~\ref{tab:appendix-counterfactual}), while the executed rule stays identical as predicted. The 19/48 falls short of the pre-registered threshold of more than 24/48, so the stored verdict reads ``PREDICTION\_FAILED: see preregistered\_prediction''. Where pairs diverge, the first divergence falls between decisions 2 and 50 (median 6), and both algorithms on one ferry task flip the goal outcome.
% Evidence: verdict "PREDICTION_FAILED: see preregistered_prediction" (E134-01); preregistered_prediction.held false,
% .executed_rule_stays_identical true, .variant_diverges_in_most_pairs false (E134-02, E134-06, E134-07);
% first divergence 2-50, median 6 (E134-11); ferry-expanded-915109 both algorithms flip (E134-12); tie-break and scope reading per issue-134-closeout.md R1.

% Paragraph purpose: Show that task clustering keeps the first adapter's last-label excess but not its heap-head agreement gap, and state M4's classification.
Clustering by task keeps the first adapter's last-label excess but not its gap in heap-head agreement. The post-hoc analysis recomputes the M4 diagnostics over the first adapter's 18 stored episodes on the earlier panel, resampling 9 task clusters (10,000 resamples, seed 133). Pooled over 132 multi-option decisions, heap-head agreement is 0.174 against a chance rate of 0.269, and the last-label rate is 0.545. The agreement gap's interval includes zero while the last-label excess excludes it (Table~\ref{tab:appendix-m4}). M4 therefore keeps its diagnostic, non-confirmatory status.
% Evidence: #134 R2 outputs/choice-frontier/o4/metrics/m4-task-cluster-intervals.json subset (E134-20); bootstrap.* (E134-32);
% pooled_point_estimates.on_policy_decisions_k_ge_2 132, .observed_agreement 0.174, .chance 0.269, .selected_last_rate 0.545 (E134-21..25);
% stored classification "post-hoc diagnostics; M4 has diagnostic (non-confirmatory) status".

\begin{table}[t]
\caption{\textbf{Descriptive post-hoc M4 diagnostics for the first adapter.} Only the last-label excess has a task-clustered interval that excludes zero. The pooled column gives point differences from chance, and the task-clustered column adds 95\% intervals.}
\label{tab:appendix-m4}
\centering\footnotesize
\setlength{\tabcolsep}{4pt}
\begin{tabular}{lcc}
\toprule[0.4mm]
\rowcolor{mygray}\textbf{Quantity} & \textbf{Pooled} & \textbf{Task-clustered [95\% CI]} \\
\midrule
Heap-head agreement minus chance & $-0.095$ & $-0.069$ $[-0.161, +0.032]$ \\
Last-label rate minus chance & $+0.276$ & $+0.277$ $[+0.126, +0.402]$ \\
\bottomrule
\end{tabular}
\end{table}
% Evidence: statistics.observed_minus_chance_agreement.pooled_point -0.0950, .task_cluster_point -0.0694, .task_cluster_95pct_ci_10000_seed133
% [-0.1609, 0.0323] (E134-24, E134-26, E134-27); statistics.last_label_minus_chance.pooled_point 0.2762, .task_cluster_point 0.2775,
% .task_cluster_95pct_ci_10000_seed133 [0.1257, 0.4021] (E134-29..31).

\section{Results Detail: Choice-Frontier Validation Panel and Ladder}
\label{app:results-ladder}

% Paragraph purpose: State that the validation panel was frozen before evaluation and that every adjacent ladder gap on it is positive, pointing to Section sec:results-ladder.
This appendix shows how the validation panel was frozen and why the ladder test passes on it, with every arm's receipts behind Section~\ref{sec:results-ladder}.

% Paragraph purpose: Show that the validation panel was fixed by rules committed in advance, then give generation, screening, and freeze receipts.
The validation panel was fixed by rules committed in advance, the ladder protocol before task generation and the panel membership before evaluation. Of 240 generated candidates (12 snapshots, seeds 935000 to 935019), 127 met the cost filter ($8 \le C^* \le 20$), and 34 passed controls-only screening, which admits tasks on which random-valid is discriminative. The freeze rule sorts admitted tasks by domain and seed, keeps at most two per domain, and takes the first 12 (Table~\ref{tab:appendix-ladder-panel}).
% Evidence: #135 protocol docs/experiments/choice-frontier/issue-135-protocol.md (commit f823850 before generation); membership
% configs/experiments/choice-frontier-v2/membership.json (commit 47de39c before evaluation); 240 and 127 per issue-135-closeout.md
% Panel (E135-60, E135-61); freeze_rule and task_ids/domains (E135-63).
% Source conflict: issue-135-closeout.md says 38; membership.json per_domain_counts.*.admitted sums to 34 (E135-62, Main binding resolution).

\begin{table}[t]
\caption{\textbf{Frozen validation panel membership.} Random-valid reached the goal in some but not all screening episodes on every task. $R$ gives uncapped exact-reference expansions (greedy / w3), and Screen gives random-valid goals in the screening episodes.}
\label{tab:appendix-ladder-panel}
\centering\footnotesize
\setlength{\tabcolsep}{3pt}
\begin{tabular}{lccc}
\toprule[0.4mm]
\rowcolor{mygray}\textbf{Task} & \textbf{$C^*$} & \textbf{$R$} & \textbf{Screen} \\
\midrule
15puzzle-935006 & 8 & 46 / 12 & 1 \\
blocksworld-935001 & 10 & 26 / 26 & 4 \\
blocksworld-935004 & 10 & 13 / 13 & 2 \\
depot-935005 & 10 & 14 / 15 & 1 \\
depot-935011 & 9 & 13 / 13 & 1 \\
elevators-935000 & 8 & 11 / 11 & 2 \\
elevators-935013 & 8 & 11 / 11 & 2 \\
grid-935016 & 8 & 9 / 9 & 4 \\
storage-935000 & 10 & 20 / 17 & 1 \\
storage-935001 & 8 & 18 / 15 & 2 \\
towers\_of\_hanoi-935000 & 13 & 24 / 25 & 1 \\
towers\_of\_hanoi-935002 & 15 & 24 / 26 & 2 \\
\bottomrule
\end{tabular}
\end{table}
% Evidence: outputs/choice-frontier/v2/metrics/analysis.json optimal_costs.<task_id> (E135-65); membership.json
% per_task.<task_id>.R.best_first_add_greedy / .best_first_add_w3 (E135-66) and per_task.<task_id>.screen_random_goals (E135-67);
% screen seeds 101/202/303/404/505 x 2 algorithms at cap 2R per issue-135-closeout.md.

\begin{table}[t]
\caption{\textbf{Development-stage validation-panel arms.} M1 falls in order from the exact reference down the $\epsilon$ rungs to random-valid. M1 has task-clustered intervals, M2 and M3 are solved-only, and the rule selectors use privileged information unavailable to the policy.}
\label{tab:appendix-ladder-arms}
\centering\footnotesize
\setlength{\tabcolsep}{2.5pt}
\begin{tabular}{lcccc}
\toprule[0.4mm]
\rowcolor{mygray}\textbf{Arm} & \textbf{M1 [95\% CI]} & \textbf{2$\times$} & \textbf{M2 $\rho$} & \textbf{M3 $\kappa$} \\
\midrule
Exact reference & 0.875 [0.875, 0.875] & 1 & 1.0 & 1.021 \\
Exact-$\epsilon$ 0.25 & 0.793 [0.764, 0.824] & 0.992 & 1.089 & 1.024 \\
Exact-$\epsilon$ 0.50 & 0.603 [0.528, 0.670] & 0.95 & 1.246 & 1.032 \\
Exact-$\epsilon$ 0.75 & 0.347 [0.252, 0.445] & 0.717 & 1.374 & 1.023 \\
Random-valid & 0.021 [0.007, 0.038] & 0.083 & 1.692 & 1.042 \\
\midrule
\texttt{hadd-greedy} & 0.882 [0.774, 0.955] & 0.958 & 0.962 & 1.083 \\
\texttt{bfs-order} & 0.000 [0.000, 0.000] & 0 & undefined & undefined \\
Novelty-first selection & 0.021 [0.000, 0.052] & 0.167 & 1.909 & 1.000 \\
\texttt{worst-first} & 0.000 [0.000, 0.000] & 0 & undefined & undefined \\
\midrule
First adapter & 0.026 [0.010, 0.042] & 0.208 & 1.810 & 1.050 \\
Scaled adapter (seed 17) & 0.349 [0.172, 0.531] & 0.583 & 1.326 & 1.140 \\
Pretrained base & 0.000 [0.000, 0.000] & 0 & undefined & undefined \\
\bottomrule
\end{tabular}
\end{table}
% Evidence: outputs/choice-frontier/v2/metrics/analysis.json arms.<arm>.m1_auc, .m1_task_cluster_95pct_ci_10000_seed133,
% .m3_solved_coverage, .m2_solved_only_geometric_rho_task_weighted, .m3_solved_kappa_mean_task_weighted (E135-20..28; exact_reference
% M3 stored 1.0212962962962964, shown to 3 dp; bfs-order and worst-first M2/M3 stored null = undefined).
% Evidence: first adapter adapter_reevaluation.learned_adapter.* (E135-41..45); pretrained_base m1_auc 0.0, CI [0.0, 0.0], solved_fraction 0 (E135-53;
% M2/M3 undefined because nothing is solved); scaled adapter outputs/choice-frontier/v3/metrics/analysis.json per_seed.17.m1_auc, .ci95,
% .solved_at_2, .m2_solved_only_geometric_rho, .m3_solved_kappa (E136-20, E136-21, E136-28..30). hadd-greedy appears only here.
% Evidence: descriptive_only_arms list in outputs/choice-frontier/v2/metrics/analysis.json; base "emits no valid first choice under its frozen 1-call cap" (E135-53).

\begin{figure}[tbp]
\centering
\includegraphics[width=.88\linewidth]{figures/fig_budget_curves.pdf}
\caption{\textbf{Solved fraction against decision budget.} On all three panels the privileged rungs fall in order between the reference and random-valid, and the trained adapter's area exceeds random-valid's. Panels are validation, held-out (P2), and unscreened (P2u, descriptive), and the zero-shot base appears where measured. The adapter band is its three-seed range, not a confidence interval. Areas under these curves are M1, and the privileged rungs test only their own reference--random axis.}
\label{fig:budget-curves}
% Evidence: outputs/choice-frontier/v2/metrics/all-episode-metrics.json, v3/metrics/all-episode-metrics.json, v4/{seeds,panels}/metrics/all-episode-metrics.json, v6/metrics/zero-shot-episode-metrics.json, v4/panels/metrics/analysis.json; figures/fig_budget_curves.py asserts each area at three decimals.
\end{figure}

% Paragraph purpose: Show that the ladder test passes on the validation panel because every adjacent gap has a positive lower bound, then give the rule, the four gaps, and the replay receipt.
The ladder test passes on the validation panel because every adjacent gap has a positive lower bound. The pre-registered rule requires all four adjacent paired task-cluster lower bounds to be positive, with arms averaged over seeds within each task and algorithm, then over algorithms. The gaps widen toward the random end of the ladder: $+0.082$ $[+0.051, +0.111]$ for the exact reference minus the 0.25 rung, $+0.190$ $[+0.126, +0.252]$ for 0.25 minus 0.50, $+0.256$ $[+0.184, +0.339]$ for 0.50 minus 0.75, and $+0.326$ $[+0.233, +0.422]$ for 0.75 minus random-valid. All 696/696 CPU episodes and the 48 first-adapter model episodes were independently replayed.
% Evidence: episode_accounting.cpu_episodes 696 / .replay_passed 696 / .missing [] (E135-30..32); adapter_reevaluation.episodes 48 (E135-40);
% instrument_validity.rule and .aggregation (E135-01, E135-04); instrument_validity.pairs[0..3].difference and .ci95
% (E135-05, E135-06, E135-08, E135-09, E135-11, E135-12, E135-14, E135-15); verdict PASS (E135-01).

% Paragraph purpose: Show that the three documented deviations affect how the panel was generated and rendered, not which tasks it holds.
Three documented deviations affect how the panel was generated and rendered, not which tasks it holds. First, generation reuses the frozen candidate generator, which reproduces all 12 snapshots byte-for-byte although domain generators take the seed differently. Second, a first concurrent generation run corrupted the blocksworld outputs through a shared temporary file. That run was discarded, all candidates were regenerated serially and reproduced, and the deviation was disclosed before evaluation. Third, only the two additive native references seed the pre-rendered views, rather than four algorithms, because native BFS exceeds the 128-expansion reference ceiling on these tasks, and unseen states are rendered live.
% Evidence: issue-135-closeout.md Execution record and issue comment 5796987720 (sheet #135 deviations 1-3, quoted).

\section{Results Detail: Scaled Adapter and Seed Replication}
\label{app:results-adapter}

% Paragraph purpose: State that this appendix holds the adapter's frozen protocol, corpus, and recipe, and per-seed results showing the gain over random-valid at every seed; point to the subsections.
This appendix documents the adapter behind Section~\ref{sec:results-adapter}, from its frozen protocol to per-seed results showing a gain over random-valid at each of three training seeds, with the held-out panels, the DAgger ablation, and the zero-shot policy as subsections.

% Paragraph purpose: Show that each part of the adapter protocol was frozen before the step it governs, then define the estimand, rule, and margins.
Each part of the adapter protocol was frozen before the step it governs: the protocol before task generation, the corpus before training, and the schedule before launch. The endpoint $D$ is the mean over training seeds of the adapter's M1 minus random-valid's M1, with a task-cluster interval (10,000 draws, seed 133). The verdict POSITIVE requires a positive lower bound, $D \ge 0.05$, and a positive estimate at every seed, and the equivalence margin is $\pm 0.05$. Controls are reused from the ladder test.
% Evidence: #136 docs/experiments/choice-frontier/issue-136-protocol.md and configs/experiments/choice-frontier-v3/protocol.json (commit ecdda03
% before generation; corpus f16dfef before training; schedule f26875b before launch); outputs/choice-frontier/v3/metrics/analysis.json
% primary.estimand, primary.bootstrap.* (E136-02, E136-03), primary.rule, primary.materiality_margin 0.05, primary.equivalence_margin 0.05
% (E136-04, E136-05); reused_from_135 and recomputed_control_m1_matches_frozen_135 true (E136-06).

% Paragraph purpose: Show that the training corpus shares no task with any evaluation panel and does not place the target last more often than chance, then give the counts.
The training corpus shares no task with any evaluation panel, and its menu augmentation places the target last no more often than chance. Of 2,400 serially generated training tasks, 1,182 remained eligible after excluding overlap with every evaluation panel and deduplicating. The seed-major walk reached 2048 records per algorithm from 236 tasks in 12 domains. Two reseeded menu orders per record give 8,192 samples, with the target last in 17.58\% of samples against a chance rate of 17.58\%.
% Evidence: 2,400 and 1,182 per issue-136-closeout.md Corpus and issue comment 5798281103 (E136-41, E136-42); outputs/choice-frontier/v3/
% preparation/report.json counts.tasks 236 (E136-43), records_per_algorithm.* 2048 (E136-44), augmentation_check.samples 8192 (E136-45),
% augmentation_check.target_last_share 0.17578125 / .chance_last_rate_mean_inverse_menu_size 0.1757589 / .tolerance_percentage_points 5.0 /
% .pass true (E136-49); walk order and seeds 945000-945015 per closeout Corpus (E136-53). 12 domains per closeout ("236 tasks from all 12 domains").

% Paragraph purpose: Show that the adapter changes the first adapter's data, not its recipe, and that the seed-17 run passed every execution check; give its secondary endpoints.
The adapter changes the first adapter's data, not its recipe. The recipe matches the first adapter's except for corpus, augmentation, sample order, and seed (LoRA rank 64, alpha 128, 128 steps, seed 17). The smoke gate accepted 34/34 calls, and all 24/24 learned episodes, capped at $2R$, replayed with no mismatch and no invalid operation. For seed 17, the 2$\times$ solve fraction is 0.583, M2 $\rho$ 1.326, and M3 $\kappa$ 1.140. Over 691 multi-option decisions heap-head agreement is 0.210 against chance 0.125, the last-label rate is 0.205, and the first-label rate is 0.090.
% Evidence: protocol.json training.optimizer_updates 128 and recipe string (E136-61); closeout audit-train (r 64, alpha 128, 128 steps, seed 17);
% smoke.json model_calls 34, accepted_calls 34 (E136-65); cap 2 x R_t (E136-67); episode_accounting learned 24 / replayed 24 / mismatches [] (E136-31);
% 0 invalid (E136-66); per_seed.17.solved_at_2 / .m2_solved_only_geometric_rho / .m3_solved_kappa (E136-28..30); .on_policy_decisions_k_ge_2 691,
% .teacher_agreement_observed 0.210, .teacher_agreement_chance 0.125 (E136-25, E136-26); .last_label_rate 0.205, .first_label_rate 0.090
% (E136-23, E136-24); ladder_position.learned.17 (E136-33).

\begin{table}[t]
\caption{\textbf{Descriptive seed-17 per-task adapter gain over random-valid.} The gain is uneven across tasks, large on both storage and both elevators tasks and zero on both blocksworld tasks. The pre-registered three-seed endpoint is reported in Table~\ref{tab:appendix-seed-d}.}
\label{tab:appendix-adapter-tasks}
\centering\footnotesize
\setlength{\tabcolsep}{3pt}
\begin{tabular}{@{}lc|lc@{}}
\toprule[0.4mm]
\rowcolor{mygray}\textbf{Task} & \textbf{$D$} & \textbf{Task} & \textbf{$D$} \\
\midrule
15puzzle-935006 & $-0.0625$ & elevators-935013 & $0.6875$ \\
blocksworld-935001 & $0.0$ & grid-935016 & $0.3625$ \\
blocksworld-935004 & $0.0$ & storage-935000 & $0.7375$ \\
depot-935005 & $0.0625$ & storage-935001 & $0.7$ \\
depot-935011 & $0.7125$ & towers\_of\_hanoi-935000 & $0.2375$ \\
elevators-935000 & $0.5$ & towers\_of\_hanoi-935002 & $0.0$ \\
\bottomrule
\end{tabular}
\end{table}
% Evidence: outputs/choice-frontier/v3/metrics/analysis.json primary.per_task_difference.<task_id> (E136-09); primary.D and .ci95 (E136-02, E136-03);
% primary.verdict POSITIVE, primary.seeds [17] (E136-01, E136-07).
% Evidence: pooled three-seed D3 per #138 outputs/choice-frontier/v4/seeds/metrics/analysis.json primary.D3 0.2847 (0.285), primary.ci95 [0.1403, 0.4326] ([0.140, 0.433]).

% Paragraph purpose: Show that of the four documented choices only the single training seed bears on what the run can show, and that the later replication closed that gap.
Four documented choices shaped this run, and only the first bears on what it can show. First, a scope edit to one training seed left seed variance unmeasured in that run. A later pre-registered replication closed the gap by training seeds 29 and 71 with the byte-identical recipe and the seed-17 sample order, so the seeds change only the LoRA initialisation and dropout. All 48 of its new episodes were independently replayed with no mismatch. Second, the step budget matches the first adapter's recipe. Third, the seed-major corpus walk resolves an ambiguity in the protocol text. Fourth, both evaluation hooks first failed because the scheduler passes only the terminal path. After a fix that derives the attempt directory from that path, the audit re-ran on the same attempts and passed without affecting GPU work or episodes.
% Evidence: protocol.json training_seeds_note ("ticket edited to '1 seed (to speed up)'"), closeout Limitations (sheet #136 choice 1);
% protocol.json optimizer_updates 128 (choice 2); protocol.json corpus_selection.walk_order (choice 3); closeout Smoke gate, evaluation and
% replay execution note, jobs/cfv3-evaluate-*/1/hook-rerun.json ok true (choice 4). Ledger 4.84/30 GPU-h (LED-2).
% Evidence (replication): issue-138-protocol.md L11-12 ("the #136 recipe, byte-identical ... Only training_seed changes"), sample order
% pinned to order:17 (issue-138-closeout.md L84-86); A138 episode_accounting.v4_replayed 48, .v4_missing [], .v4_replay_mismatches [].

% Paragraph purpose: Show that the adapter beats random-valid at every seed on the validation panel while sitting near or below the 0.75 rung, with per-seed D, M1, and S.
On the validation panel the adapter beats random-valid at every seed, but it sits near or below the exact-$\epsilon$ 0.75 rung. Per-seed $D$ is $+0.328$ $[+0.155, +0.508]$ for seed 17, $+0.229$ $[+0.108, +0.354]$ for seed 29, and $+0.297$ $[+0.134, +0.466]$ for seed 71, so every lower bound is positive. Per-seed M1 is 0.349, 0.250, and 0.318, with a between-seed standard deviation of 0.051 (descriptive). The pre-registered difference from the 0.75 rung is $S = +0.002$, $-0.097$, and $-0.029$ per seed. Against the 0.50 rung the three-seed $S$ is $-0.298$ $[-0.468, -0.129]$ (descriptive), below the rung. Table~\ref{tab:appendix-seed-d} sets these values beside the held-out and unscreened panels.
% Evidence: A138 primary.per_seed.{17,29,71} 0.3281/0.2292/0.2969 and primary.per_seed_ci95.{17,29,71} [0.1552, 0.5083]/[0.1083, 0.3542]/[0.1344, 0.4656];
% seed_variance.m1_per_seed.{17,29,71}.m1_auc 0.3490/0.25/0.3177; seed_variance.m1_between_seed_sd 0.0506 (0.051, statistics.stdev n-1);
% separation.vs_eps_0.75.per_seed.{17,29,71} 0.0021/-0.0969/-0.0292 (0.002/-0.097/-0.029); separation.vs_eps_0.50.S -0.2976, .ci95
% [-0.4684, -0.1288] (role "descriptive only", P138 L51).

% Paragraph purpose: Show that the adapter picks the last menu label less often than the first adapter did (descriptive).
The adapter also picks the last menu label less often than the first adapter did. Its last-label rate per seed is 0.205, 0.192, and 0.171 (chance about 0.13, descriptive), against 0.300 (chance 0.144) for the first adapter on the same panel.
% Evidence: A138 per_seed.{17,29,71}.last_label_rate 0.2055 / 0.1917 / 0.1712, rounded half-up to 3 dp from the named keys.
% Evidence: A138 per_seed.{17,29,71}.teacher_agreement_chance 0.1247 / 0.1306 / 0.1267, written "about 0.13";
% first adapter outputs/choice-frontier/v2/metrics/analysis.json adapter_reevaluation.learned_adapter.m4_teacher_agreement.selected_last_rate 0.2997 (0.300) and .chance 0.1444 (0.144) (E135-50).

\subsection{Held-Out Panels}
\label{app:results-panels}

% Paragraph purpose: State that both fresh panels were frozen before any evaluation episode and that this subsection gives every arm's results on them.
This subsection shows that both fresh panels were frozen before any evaluation episode and gives every arm's results on them. The held-out panel appears as P2 in tables and the unscreened panel as P2u.

% Paragraph purpose: Show that the held-out panel repeats the validation panel's construction on fresh seeds and met its minimum size, so its tests are confirmatory.
The held-out panel repeats the validation panel's construction on fresh seeds. Both fresh panels come from 2,400 candidates generated from seeds 955000--955199. Removing overlap with earlier panels (1,168 candidates), duplicates (205), goals at the initial state (31), and $C^*$ out of range (536) left 460. The held-out panel then applies the validation panel's generation, $C^*$, and screening rules, keeping a candidate when $8 \le C^* \le 20$ and random-valid reaches the goal in 1 to 9 of its 10 screening episodes. The freeze walk, sorted by domain and seed, takes at most 2 tasks per domain up to 12 tasks, with a minimum of 8 tasks from 5 domains. Only 11 admitted tasks were available under the per-domain limit, so the panel froze with 11 tasks from 8 domains. That meets the minimum, so every held-out test is confirmatory.

% Paragraph purpose: Show that the unscreened panel drops the random-valid screen, so it shows descriptively whether the gain survives without it; give its composition and the zoo replay.
The unscreened panel drops the random-valid screen, so it shows descriptively whether the gain survives without it. It takes kept candidates outside the held-out panel on which the exact reference reaches the goal under both algorithms, at most 2 per domain, without consulting the screen. This gives 12 tasks from 7 domains, and all 12 happen to have 0/10 random-valid goals in screening. All 638/638 held-out-panel and 696/696 unscreened-panel zoo episodes were independently replayed, with none missing.
% Evidence: P2 rule issue-139-protocol.md L7, L36, L48-50; P2u rule L55-57; candidate pool issue-139-closeout.md L46-54 (2,400; 1,168 / 205 / 31 / 536; 460);
% P2 froze at 11 issue-139-closeout.md L28; A139 pooled.panels {p135 12, p2 11, p2u 12}; concentration.per_panel.{p2,p2u}.domains (8, 7 domains);
% "All 12 happen to have 0/10 random goals." issue-139-closeout.md L56; zoo replay issue-139-closeout.md L113-115.

% Paragraph purpose: Show that the seed range was extended once, before any evaluation episode and with every other rule unchanged; disclose the amendment in full.
The seed range was extended once, before any evaluation episode, with every other rule unchanged. The first freeze, from the original 40-seed pool 955000--955039, admitted 6 tasks from 5 domains and took the protocol's blocked exit. At that point only controls-only screening had run. The amendment extended the seeds to 955000--955199, kept the 480 earlier candidates unchanged, and generated the new seeds 955040--955199 serially after them. The generation mechanism, the exclusion sources, the duplicate check, the $C^*$ rule, $R_t$, the screen, the freeze walk, and the unscreened-panel rule all stayed fixed, and both panels were frozen from the combined pool.
% Evidence: issue-139-protocol.md Amendment A1 L118-127 (L120 trigger 6 tasks / 5 domains, no evaluation episode; L122 seeds extended,
% 480 candidates kept; L123 everything else unchanged); issue-139-closeout.md L58 and L5-6 (commit d567939 before any extension seed).

% Paragraph purpose: Show that the ladder test passes on the held-out panel and, descriptively, on the unscreened panel, with all eight adjacent gaps.
The ladder test passes on the held-out panel and, descriptively, on the unscreened panel. On the held-out panel the four pre-registered adjacent gaps, from the exact reference down to random-valid, are $+0.130$ $[+0.070, +0.190]$, $+0.239$ $[+0.183, +0.290]$, $+0.302$ $[+0.259, +0.340]$, and $+0.189$ $[+0.136, +0.248]$, so every lower bound is positive. On the unscreened panel the descriptive gaps are $+0.197$ $[+0.154, +0.243]$, $+0.281$ $[+0.191, +0.355]$, $+0.248$ $[+0.209, +0.290]$, and $+0.147$ $[+0.088, +0.226]$, and every lower bound is again positive. Table~\ref{tab:appendix-panel-arms} gives each arm's M1 on all three panels.
% Evidence: A139 ladder_p2.pairs[0..3].{difference,ci95} 0.1295 [0.0705, 0.1898] / 0.2386 [0.1830, 0.2898] / 0.3023 [0.2591, 0.3398] /
% 0.1886 [0.1364, 0.2477], verdict PASS, status confirmatory (P139 L93); unscreened_p2u.ladder.pairs[0..3] 0.196875 [0.1542, 0.2427] /
% 0.28125 [0.1906, 0.3552] / 0.2479 [0.2094, 0.2896] / 0.146875 [0.0875, 0.2260], verdict PASS, descriptive (P139 L96).
% Note: 0.196875 and 0.28125 have no 3-dp rendering on the sheet; printed at the table's 3 dp.

\begin{table}[t]
\caption{\textbf{Descriptive M1 by panel and arm.} On every panel the adapter's three-seed mean lies above random-valid and below the exact reference. Values have task-clustered 95\% intervals, and the exact-$\epsilon$ rungs and hadd-greedy use privileged information the policy never sees.}
\label{tab:appendix-panel-arms}
\centering\footnotesize
\setlength{\tabcolsep}{5pt}
\begin{tabular}{@{}lll@{}}
\toprule[0.4mm]
\rowcolor{mygray}\textbf{Panel} & \textbf{Arm} & \textbf{M1 [95\% CI]} \\
\midrule
\multirow{11}{*}{Validation (12)} & Exact reference & 0.875 [0.875, 0.875] \\
 & Exact-$\epsilon$ 0.25 & 0.793 [0.764, 0.824] \\
 & Exact-$\epsilon$ 0.50 & 0.603 [0.528, 0.670] \\
 & Exact-$\epsilon$ 0.75 & 0.347 [0.252, 0.445] \\
 & Random-valid & 0.021 [0.007, 0.038] \\
 & Pretrained base & 0.000 [0.000, 0.000] \\
 & Adapter, seed 17 & 0.349 [0.172, 0.531] \\
 & Adapter, seed 29 & 0.250 [0.125, 0.375] \\
 & Adapter, seed 71 & 0.318 [0.151, 0.490] \\
 & Adapter, three-seed mean & 0.306 [0.158, 0.455] \\
\midrule
 & hadd-greedy (privileged) & 0.882 [0.774, 0.955] \\
\midrule
\multirow{11}{*}{P2 (11)} & Exact reference & 0.875 [0.875, 0.875] \\
 & Exact-$\epsilon$ 0.25 & 0.745 [0.685, 0.805] \\
 & Exact-$\epsilon$ 0.50 & 0.507 [0.445, 0.566] \\
 & Exact-$\epsilon$ 0.75 & 0.205 [0.145, 0.269] \\
 & Random-valid & 0.016 [0.000, 0.048] \\
 & Pretrained base & 0.000 [0.000, 0.000] \\
 & Adapter, seed 17 & 0.506 [0.313, 0.699] \\
 & Adapter, seed 29 & 0.284 [0.119, 0.489] \\
 & Adapter, seed 71 & 0.506 [0.295, 0.705] \\
 & Adapter, three-seed mean & 0.432 [0.273, 0.602] \\
\midrule
 & hadd-greedy (privileged) & 0.928 [0.883, 0.966] \\
\midrule
\multirow{11}{*}{\legendbg{yellow!25}{P2u (12)}} & Exact reference & 0.875 [0.875, 0.875] \\
 & Exact-$\epsilon$ 0.25 & 0.678 [0.632, 0.721] \\
 & Exact-$\epsilon$ 0.50 & 0.397 [0.331, 0.478] \\
 & Exact-$\epsilon$ 0.75 & 0.149 [0.091, 0.227] \\
 & Random-valid & 0.002 [0.000, 0.006] \\
 & Pretrained base & 0.000 [0.000, 0.000] \\
 & Adapter, seed 17 & 0.198 [0.031, 0.396] \\
 & Adapter, seed 29 & 0.120 [0.026, 0.224] \\
 & Adapter, seed 71 & 0.135 [0.031, 0.250] \\
 & Adapter, three-seed mean & 0.151 [0.030, 0.288] \\
\midrule
 & hadd-greedy (privileged) & 0.906 [0.883, 0.933] \\
\bottomrule
\end{tabular}
% Evidence: outputs/choice-frontier/v4/panels/metrics/analysis.json arms.{p135,p2,p2u}.{exact_reference,exact-eps-0.25,exact-eps-0.50,exact-eps-0.75,
% random_valid,pretrained_base,learned_adapter_s17,learned_adapter_s29,learned_adapter_s71,learned_adapter_seed_mean,hadd-greedy}.{m1,ci95},
% rounded half-up to 3 dp from the named keys (descriptive, P139 L108). Note: arms.p2u.exact-eps-0.25.m1 is 0.678125, printed at 3 dp.
\end{table}

% Paragraph purpose: Show that every seed's gain over random-valid has a positive lower bound on every panel, and state the limits on reading the rung separations.
Every seed's gain over random-valid has a positive lower bound on every panel, but the rung separations need care. Table~\ref{tab:appendix-seed-d} gives $D$ per seed and $D_3$ with intervals on each panel and pooled over 35 tasks, together with the separation $S$ from the exact-$\epsilon$ rungs. Seed 29 has the lowest per-seed $D$ on every panel, and on the held-out panel its separation point from the 0.75 rung is $+0.080$ (descriptive). No equivalence margin was pre-registered for $S$, so a non-separated $S$ is not equivalence to the rung. The held-out separation rests on 11 task clusters, with 8 positive and 3 negative per-task differences, and percentile intervals can under-cover at this size.
% Evidence: A139 heldout_p2.separation.vs_eps_0.75.per_task_difference signs 8 positive / 3 negative / 0 zero
% of 11 (negative: 15puzzle-955075, blocksworld-955022, blocksworld-955127; listing, descriptive).
% Evidence: "Seed 29 is weaker. On P2 its separation point is +0.080, and it has the lowest D3 on every panel." issue-139-closeout.md L171;
% A139 heldout_p2.separation.vs_eps_0.75.per_seed.29 0.0795 (0.080).

\begin{table}[t]
\caption{\textbf{Seed and rung contrasts on development panels.} The adapter separates above random-valid on every panel, and above the 0.75 rung only on the held-out panel. Estimates are paired task-clustered 95\% intervals for adapter minus random-valid ($D$) and minus a privileged $\epsilon$ rung ($S$), and yellow rows are descriptive. \cmark\ means separated above, $=$ not separated, and $\blacktriangledown$ separated below.}
\label{tab:appendix-seed-d}
\centering\footnotesize
\setlength{\tabcolsep}{4pt}
\begin{tabular}{@{}llcc@{}}
\toprule[0.4mm]
\rowcolor{mygray}\textbf{Panel} & \textbf{Contrast} & \textbf{Estimate [95\% CI]} & \textbf{Verdict} \\
\midrule
\rowcolor{yellow!25}Validation (12) & $D$, seed 17 & $+0.328$ $[+0.155, +0.508]$ &  \\
\rowcolor{yellow!25} & $D$, seed 29 & $+0.229$ $[+0.108, +0.354]$ &  \\
\rowcolor{yellow!25} & $D$, seed 71 & $+0.297$ $[+0.134, +0.466]$ &  \\
 & $D_3$ & $+0.285$ $[+0.140, +0.433]$ & \cmark \\
 & $S$ vs. $\epsilon{=}0.75$ & $-0.041$ $[-0.187, +0.102]$ & $=$ \\
\rowcolor{yellow!25} & $S$ vs. $\epsilon{=}0.50$ & $-0.298$ $[-0.468, -0.129]$ & $\blacktriangledown$ \\
\midrule
\rowcolor{yellow!25}P2 (11) & $D$, seed 17 & $+0.490$ $[+0.294, +0.688]$ &  \\
\rowcolor{yellow!25} & $D$, seed 29 & $+0.268$ $[+0.093, +0.477]$ &  \\
\rowcolor{yellow!25} & $D$, seed 71 & $+0.490$ $[+0.276, +0.694]$ &  \\
 & $D_3$ & $+0.416$ $[+0.250, +0.591]$ & \cmark \\
 & $S$ vs. $\epsilon{=}0.75$ & $+0.227$ $[+0.053, +0.403]$ & \cmark \\
\rowcolor{yellow!25} & $S$ vs. $\epsilon{=}0.50$ & $-0.075$ $[-0.247, +0.089]$ & $=$ \\
\midrule
\rowcolor{yellow!25}P2u (12) & $D$, seed 17 & $+0.196$ $[+0.031, +0.396]$ &  \\
\rowcolor{yellow!25} & $D$, seed 29 & $+0.118$ $[+0.026, +0.220]$ &  \\
\rowcolor{yellow!25} & $D$, seed 71 & $+0.133$ $[+0.031, +0.248]$ &  \\
\rowcolor{yellow!25} & $D_3$ & $+0.149$ $[+0.030, +0.283]$ & \cmark \\
\rowcolor{yellow!25} & $S$ vs. $\epsilon{=}0.75$ & $+0.002$ $[-0.145, +0.157]$ & $=$ \\
\rowcolor{yellow!25} & $S$ vs. $\epsilon{=}0.50$ & $-0.246$ $[-0.400, -0.082]$ & $\blacktriangledown$ \\
\midrule
\rowcolor{yellow!25}Pooled (35) & $D$, seed 17 & $+0.334$ $[+0.223, +0.450]$ &  \\
\rowcolor{yellow!25} & $D$, seed 29 & $+0.203$ $[+0.124, +0.290]$ &  \\
\rowcolor{yellow!25} & $D$, seed 71 & $+0.301$ $[+0.200, +0.411]$ &  \\
\rowcolor{yellow!25} & $D_3$ & $+0.279$ $[+0.191, +0.375]$ & \cmark \\
\rowcolor{yellow!25} & $S$ vs. $\epsilon{=}0.75$ & $+0.058$ $[-0.040, +0.156]$ & $=$ \\
\bottomrule
\end{tabular}
% Evidence: validation panel A138 primary.{per_seed,per_seed_ci95,D3,ci95,verdict}; separation.vs_eps_0.75.{S,ci95,verdict} (co-primary, P138 L44-52);
% separation.vs_eps_0.50.{S,ci95,verdict} (descriptive only, P138 L51). P2 A139 heldout_p2.primary.* (P139 L94); heldout_p2.separation.vs_eps_0.75.*
% (P139 L95); heldout_p2.separation.vs_eps_0.50.* (descriptive, issue-139-closeout.md L30). P2u A139 unscreened_p2u.primary.* (P139 L96);
% unscreened_p2u.separation.* (descriptive, P139 L96). Pooled A139 pooled.primary.*, pooled.separation_vs_eps_0.75.* (descriptive, P139 L97-98);
% no pooled separation against the 0.50 rung is stored. Values rounded half-up to 3 dp from the named keys.
\end{table}

% Paragraph purpose: Show that no single domain carries the gain over the validation and held-out panels, while on the unscreened panel it comes from four of twelve tasks (all descriptive).
Over the validation and held-out panels no single domain carries the gain, while on the unscreened panel it comes from four of twelve tasks. Table~\ref{tab:appendix-concentration} counts tasks by the sign of their three-seed difference and leaves out one domain at a time over the validation and held-out panels. Over the 35-task union 22 tasks are positive, 12 are zero, and 1 is negative. On the unscreened panel the eight zero tasks are both 15-puzzle, both blocksworld, both depot, and both driverlog tasks, on which neither the adapter nor random-valid solves within $2R_t$. Every leave-one-domain-out $D_3$ lies from $+0.305$ to $+0.395$, and every lower bound is at least $+0.185$. All of these results are descriptive.
% Evidence: A139 concentration.all_36_counts {positive 22, zero 12, negative 1}; concentration.per_panel.p2u.counts.zero 8 and
% .per_task_difference (issue-139-closeout.md L73); concentration.leave_one_domain_out_p2_union_p135.* (P139 L101-104, descriptive).
% Note: key concentration.all_36_counts holds the 35-task union (P2 froze at 11; issue-139-closeout.md L28).

\begin{table}[t]
\caption{\textbf{Descriptive concentration of three-seed gains.} Leaving out any one domain keeps $D_3$ positive with a positive lower bound. The upper rows count per-task signs, and the lower rows leave one domain out of the validation--P2 union.}
\label{tab:appendix-concentration}
\centering\footnotesize
\setlength{\tabcolsep}{4pt}
\begin{tabular}{lccc}
\toprule[0.4mm]
\rowcolor{mygray}\textbf{Panel} & \textbf{Positive} & \textbf{Zero} & \textbf{Negative} \\
\midrule
Validation (12) & 9 & 2 & 1 \\
P2 (11) & 9 & 2 & 0 \\
P2u (12) & 4 & 8 & 0 \\
Union (35) & 22 & 12 & 1 \\
\midrule
\rowcolor{mygray}\textbf{Domain left out} & \textbf{$D_3$} & \textbf{95\% interval} & \textbf{Tasks left} \\
\midrule
15-puzzle & $+0.384$ & $[+0.269, +0.501]$ & 21 \\
blocksworld & $+0.395$ & $[+0.269, +0.520]$ & 19 \\
depot & $+0.353$ & $[+0.224, +0.484]$ & 20 \\
elevators & $+0.305$ & $[+0.185, +0.430]$ & 20 \\
ferry & $+0.331$ & $[+0.212, +0.455]$ & 21 \\
grid & $+0.331$ & $[+0.204, +0.464]$ & 20 \\
storage & $+0.332$ & $[+0.209, +0.459]$ & 21 \\
Tower of Hanoi & $+0.384$ & $[+0.262, +0.507]$ & 20 \\
visitall & $+0.318$ & $[+0.211, +0.422]$ & 22 \\
\bottomrule
\end{tabular}
% Evidence: A139 concentration.per_panel.{p135,p2,p2u}.counts (9/2/1, 9/2/0, 4/8/0) and concentration.all_36_counts (22/12/1);
% concentration.leave_one_domain_out_p2_union_p135.<domain>.{D3,ci95,tasks_left}, all verdict_rule_applied POSITIVE (P139 L104).
% depot D3 is stored as 0.3525; 0.353 is the issue-139-closeout.md L79 rendering.
% Note: key concentration.all_36_counts holds the 35-task union (P2 froze at 11; issue-139-closeout.md L28).
\end{table}

\begin{figure}[tbp]
\centering
\includegraphics[width=.78\linewidth]{figures/fig_task_gains.pdf}
\caption{\textbf{Descriptive task-level choice-frontier gains.} Most tasks gain, and the unscreened panel's improvement is concentrated on four tasks. Each point is the three-seed adapter mean minus random-valid M1, averaged across the two additive algorithms for a task, and colours identify domains. These points are not task-level significance tests.}
\label{fig:task-gains}
% Evidence: outputs/choice-frontier/v4/panels/metrics/analysis.json concentration.{per_panel,all_36_counts}; figures/fig_task_gains.py checks the plotted 35-task union (22 positive, 12 zero, one negative) and the unscreened panel (4 positive, 8 zero).
\end{figure}


\subsection{Illustrative Search Episodes}
\label{app:qualitative}
These rule-selected traces illustrate what the adapter saw and selected on one solved validation task and one unsolved held-out task. They do not estimate a task-level effect.
% Selection rules fixed before inspecting candidate episodes; figures are illustrative descriptive development-stage examples, not new efficacy claims.
% Evidence: outputs/choice-frontier/v4/seeds/metrics/all-episode-metrics.json records.{task,algorithm,arm,training_seed,auc,decision_cap,termination}; outputs/choice-frontier/v2/metrics/all-episode-metrics.json records.{task,algorithm,arm,auc}; outputs/choice-frontier/v4/panels/metrics/all-episode-metrics.json.p2[].{task,algorithm,arm,decision_cap,termination}.
% Evidence: outputs/choice-frontier/v4/seeds/metrics/all-episode-metrics.json records.{task,algorithm,training_seed,auc,teacher_agreements,teacher_decisions}; outputs/choice-frontier/v2/metrics/all-episode-metrics.json records.{task,algorithm,arm,auc}; outputs/choice-frontier/v3/evaluation/episodes/choice-frontier-v2__grid-expanded-935016/best_first_add_w3-learned_adapter-s17-17.json.gz events and result; outputs/choice-frontier/v2/zoo/episodes/choice-frontier-v2__grid-expanded-935016/best_first_add_w3-{exact_reference,random_valid}-m2-17.json.gz events and result. Script asserts three-seed M1 0.792, five-seed random-valid M1 0.075, difference 0.717, goal at decision 10 for exact and adapter, random-valid cap 18, displayed decision indices, and adapter reference-head agreement 7/9.
% Evidence: outputs/choice-frontier/v2/views/grid-expanded-935016/scenes/catalog.json.gz stored_scene_size=[128,128], path_bindings, states[].scene_path; v3/evaluation/views/.../views.json.gz states[].scene_path. Script re-renders each selected frame at 128 px, asserts pixel equality with its cached policy PNG, and then uses identical renderer arguments at 256 px without labels. Catalog frames use their original empty object-name set; episode frames use task_context.objects_by_type.
\begin{figure}[tbp]
\centering
\includegraphics[width=.94\linewidth]{figures/fig_qualitative_success.pdf}
\caption{\textbf{On a solved validation task, the adapter follows the reference to the goal.} In this rule-selected weighted-A* episode, the seed-17 adapter agrees with the reference on 7/9 comparable decisions, and both reach the goal at decision 10, while random-valid exhausts its decision cap at 18. Rows show expanded scenes, and the plot traces their $h_{\mathrm{add}}$ after the first expansion. The example is illustrative, not an efficacy estimate.}
\label{fig:qualitative-success}
\end{figure}

% Evidence: outputs/choice-frontier/v3/evaluation/episodes/choice-frontier-v2__grid-expanded-935016/best_first_add_w3-learned_adapter-s17-17.json.gz events[9].{menu,raw_output,trusted_runtime_result}, previous event frontier_after.head and all prior admissions[].trusted_runtime_result.h; v3/evaluation/views/.../views.json.gz states[].vfg; script asserts the twelve displayed h_add values [6,10,2,4,2,10,2,5,2,3,8,0] and menu decision 10.
% Evidence: outputs/choice-frontier/v4/seeds/metrics/all-episode-metrics.json records.{teacher_agreements,teacher_decisions,teacher_chance_sum,task,algorithm,training_seed}; outputs/choice-frontier/v3/evaluation/episodes/ and outputs/choice-frontier/v4/seeds/evaluation/episodes/ adapter events.{menu,raw_output,trusted_runtime_result.admissions}; script asserts 429/2117 agreement, uniform teacher-chance sum 269.711, 2231 total rank-histogram decisions, observed bins [1041,392,295,279,224], uniform expected bins [650.023,420.860,431.063,420.860,308.195].
% Evidence: fig_qualitative.py verifies every displayed menu thumbnail against its source 128 px PNG before rendering the same VFG at 256 px.
\begin{figure}[tbp]
\centering
\includegraphics[width=.94\linewidth]{figures/fig_qualitative_decision.pdf}
\caption{\textbf{The adapter picks the reference's state at the episode's largest menu and prefers low-$h_{\mathrm{add}}$ states overall.} At decision 10 of Figure~\ref{fig:qualitative-success}, the adapter and the privileged reference both select $c_{11}$ from 12 choices. The policy sees only opaque labels and scenes, not the heuristic values or the reference pick. The histogram describes the adapter's heuristic ranks across validation decisions against uniform choice on the same menus and is not an efficacy estimate.}
\label{fig:qualitative-decision}
\end{figure}

% Evidence: outputs/choice-frontier/v4/panels/metrics/all-episode-metrics.json.p2[].{task,algorithm,arm,training_seed,decision_cap,decisions,termination}; outputs/choice-frontier/v4/panels/p2/evaluation/episodes/choice-frontier-v4__15puzzle-expanded-955075/best_first_add_greedy-learned_adapter-s17-17.json.gz events[].trusted_runtime_result.{expanded_state_id,admissions} and result; outputs/choice-frontier/v4/panels/p2/zoo/episodes/choice-frontier-v4__15puzzle-expanded-955075/best_first_add_greedy-exact_reference-m2-17.json.gz events and result; script asserts cap 92, exact goal decision 47, adapter minimum h_add=9 at decision 40 and final h_add=35 at decision 92.
% Evidence: outputs/choice-frontier/v4/panels/views/15puzzle-expanded-955075/scenes/frames/state-{000000,000060}.png and p2/evaluation/views/.../state-000126.png are original cached 128x128 policy scenes. The initial scene omits embedded tile glyphs present in later cached frames; a direct VFG re-render changes its content, so the displayed pixels come from the cached PNGs.
\begin{figure}[tbp]
\centering
\includegraphics[width=.94\linewidth]{figures/fig_qualitative_failure.pdf}
\caption{\textbf{On a held-out puzzle, the adapter stalls where the reference solves.} On this rule-selected puzzle, the reference reaches the goal at decision 47, while every adapter seed exhausts the cap at 92. The shown seed-17 adapter's best-so-far $h_{\mathrm{add}}$ stops improving at 9 after decision 40. The first cached scene lacks tile glyphs that later scenes show. The trace is illustrative and does not establish why the adapter fails.}
\label{fig:qualitative-failure}
\end{figure}

\clearpage

\subsection{DAgger Ablation}
\label{app:results-cf-dagger}

% Paragraph purpose: State that the pre-registered DAgger ablation finds no detectable change in choice quality after one on-policy round at one seed.
This subsection reports the pre-registered DAgger ablation, which finds no detectable change in the adapter's choice quality after one on-policy round at one seed.

% Paragraph purpose: Show that the ablation asks whether one on-policy round lifts the adapter above the 0.75 rung, and define the procedure and endpoints before any result.
The ablation asks whether one round of on-policy data lifts the adapter above the exact-$\epsilon$ 0.75 rung. It collects on-policy records with the seed-17 adapter under each additive algorithm, aggregates them with replay records from the frozen corpus, retrains, and evaluates on the validation and held-out panels. The primary endpoint is the paired separation $S$ of the DAgger adapter from the 0.75 rung pooled over the 23 tasks, with SEPARATED\_ABOVE requiring a positive lower bound. The co-primary is the change $\Delta$ in M1 from the pre-DAgger adapter, with equivalence margin $\pm 0.05$.
% Evidence: #140 outputs/choice-frontier/v5/metrics/analysis.json primary.rule, primary.tasks 23, primary.panels {p2 11, v2 12} (P140 L66-72);
% co_primary.equivalence_margin 0.05, co_primary.rule.order IMPROVED/EQUIVALENT/WORSE/INCONCLUSIVE (P140 L74-77); collection cells
% best_first_add_greedy-s17 and best_first_add_w3-s17 (collection.cells.*).

% Paragraph purpose: Show that a deadline cut the ablation to one seed and half the planned corpus before any outcome; disclose the amendment.
A deadline cut the ablation to one seed and half the planned corpus, and the cut was made before any DAgger training or evaluation outcome. The amendment, at 11:21 UTC on 2026-09-24, kept only seed 17 because the three-seed plan could not fit a hard wall-clock deadline. The running seed-29 collection was cancelled with 1.79 GPU-h still charged, and seed 71 was not run. Each algorithm's training set became the first 512 on-policy seed-17 records plus 512 replay records, which the menu augmentation turns into 2048 samples over 64 steps, with the recipe otherwise unchanged.
% Evidence ("half the planned corpus"): wording copied from the Seed Replication and Follow-Up Limits paragraph ("trains one seed on half the planned corpus").
% Evidence: issue-140-protocol.md Amendment A1 heading L118 (11:21 UTC, before any outcome), L120 (deadline), L122 (s29 cancelled, s71 not run),
% L123 (first 512 + 512, 2048 samples, 64 steps); issue-140-closeout.md L80 (1.79 GPU-h charged); A140 collection.cells.<cell>.{on_policy_records 512,
% replay_records 512, samples 2048}; B5 attempts[4..5].total 64; A140 primary.seeds [17].

% Paragraph purpose: Show that one round at one seed produces no detectable change, so the ablation is inconclusive; give endpoints, descriptive secondaries, and replay.
One DAgger round at one seed produces no detectable change in measured choice quality, so the ablation is inconclusive. The primary endpoint is NOT\_SEPARATED, with $S = +0.115$ $[-0.026, +0.258]$ over 23 tasks. The co-primary is INCONCLUSIVE, with $\Delta = -0.030$ $[-0.114, +0.057]$ against the $\pm 0.05$ margin. As descriptive secondary results, $S$ is $+0.278$ $[+0.082, +0.473]$ on the held-out panel and $-0.034$ $[-0.197, +0.135]$ on the validation panel, and the last-label rate on the validation panel falls from 0.205 to 0.137. All 24/24 validation-panel and 22/22 held-out-panel DAgger episodes were independently replayed with no mismatch.
% Evidence: A140 episode_accounting.{v2,p2}.episodes_replayed 24 / 22, replay_mismatches [] (issue-140-closeout.md L53);
% primary.verdict NOT_SEPARATED, primary.S_pool 0.1152 (0.115), primary.ci95 [-0.0255, 0.2576] ([-0.026, 0.258]), primary.tasks 23;
% co_primary.verdict INCONCLUSIVE, co_primary.delta_pool -0.0299 (-0.030), co_primary.ci95 [-0.1141, 0.0571] ([-0.114, 0.057]);
% per_panel.p2.S 0.2784 (0.278), .S_ci95 [0.0818, 0.4727]; per_panel.v2.S -0.0344 (-0.034), .S_ci95 [-0.1969, 0.1354] (descriptive secondary, P140 L79-88);
% per_seed.v2.17.pre_dagger.last_label_rate 0.2055 (0.205), per_seed.v2.17.dagger.last_label_rate 0.1371 (0.137) (descriptive, P140 L84).

\subsection{Zero-Shot Policy}
\label{app:results-cf-zeroshot}

% Paragraph purpose: State that the untrained base, prompted under the same contract, scores at the random-valid level.
This subsection shows that the untrained base model, prompted under the same contract, observation, and cap, is not separated from random-valid.

% Paragraph purpose: Define the zero-shot second policy as the prompted base with no adapter.
The zero-shot policy is the pretrained Qwen3-VL-8B-Instruct base with no adapter, prompted to choose by the goal. It receives the frozen scene-only choice-frontier observation and one appended system message, which states the JSON format and asks for the frontier state whose scene appears closest to satisfying the goal pages. The policy uses a fixed three-rule output extractor, greedy decoding at inference seed 17, and the adapter's cap of $2R_t$ calls. It ran on both additive algorithms over the validation and held-out panels, 46 episodes in all.

% Paragraph purpose: Show that the trained second backbone was ruled out before any outcome because it did not fit the compute cap.
A trained second backbone was ruled out before any outcome because it did not fit the compute cap. The preferred arm, InternVL3.5-8B trained with the frozen three-seed recipe, has a mean observation of 11,748 tokens against 3,111 for Qwen3-VL, which puts training alone at 43 to 55 GPU-h against the 30 GPU-h cap. A reduced recipe would break the byte-identical condition, so the arm was not run in any form. The protocol was committed before any GPU launch, and no amendment followed.
% Evidence: configs/experiments/choice-frontier-v6/protocol.json arm.system_message_suffix (JSON-format sentence plus "Choose the frontier state whose scene appears closest to satisfying the goal pages.").
% Evidence: #141 configs/experiments/choice-frontier-v6/protocol.json arm.condition zero_shot_base, model_id/model_revision 0c351dd0, arm.adapter null,
% arm.system_message_suffix, arm.extraction (3 rules), inference.do_sample false, inference_seed 17; outputs/choice-frontier/v6/metrics/analysis.json
% arm.decision_cap "2 x R_t"; evaluation.model_episodes 46 (P141 L16-41). InternVL: protocol.json arm_choice.feasibility_evidence.internvl_mean_tokens
% 11747.9985 (11,748) and .qwen_mean_tokens 3110.5199 (3,111), rounded to the nearest token; arm_choice.reason "43-55 GPU-h training alone";
% cap 30 GPU-h (P141 L13-14); "not run in any form" (P141 L14, issue-141-closeout.md L7). Commit 0fa497c before launch; "No amendments." (issue-141-closeout.md L3).

% Paragraph purpose: Show that the zero-shot base scores at the random-valid level and below both the 0.75 rung and the trained adapter, as pre-registered, with stage labels and the held-out margin detail.
The zero-shot base is not separated from random-valid and falls below both the 0.75 rung and the trained adapter, as pre-registered (Table~\ref{tab:appendix-zeroshot}). Here $D_3$ names the pre-registered learned-minus-random-valid endpoint for both arms, a single realisation for the zero-shot arm. Validation-panel verdicts are development-stage and held-out-panel verdicts are held-out confirmatory. On the held-out panel the upper bound of $D_3$ exceeds the $\pm 0.05$ margin by 0.001, so the frozen rule returns INCONCLUSIVE rather than EQUIVALENT. The arm's M1 is 0.026 [0.000, 0.063] on the validation panel, 0.017 [0.000, 0.051] on the held-out panel, and 0.022 [0.003, 0.046] pooled, below the exact-$\epsilon$ 0.75 rung on every panel. The pre-registered prediction, $D_3$ within the margin with $S$ and the adapter contrast $A$ separated below on both panels, held.
% Evidence: analysis.json primary.v2.stage "development", primary.p2.stage "held-out confirmatory"; arm.realisations "one per (task, algorithm)".
% Evidence: primary.p2.ci95[1] 0.05114 vs margin 0.05 (issue-141-closeout.md L32); per_panel.v2.m1_arm 0.0260, .m1_arm_ci95 [0.0, 0.0625]; per_panel.p2.m1_arm 0.0170,
% .m1_arm_ci95 [0.0, 0.0511]; pooled.m1_arm 0.0217, .m1_arm_ci95 [0.0027, 0.0462]; ladder_position.*.next_higher_rung exact-eps-0.75 (descriptive, P141 L66);
% Evidence (prediction, copied from the block below): Prediction P141 L72, held per issue-141-closeout.md L54-59.

% Paragraph purpose: Show that the base always emits a valid choice but agrees with the reference's heap head at about the chance rate (descriptive), with replay.
The base always emits a valid choice, but its choices agree with the reference's heap head at about the chance rate. All 1,592 model calls produced a valid choice through the first extractor rule. Heap-head agreement minus chance is $+0.003$ on the validation panel and $+0.001$ on the held-out panel, the last-label rate is 0.328 and 0.285 against a chance rate of about 0.13 to 0.14, and 4 of 46 episodes solve at $2R_t$ (descriptive). All 46/46 episodes were independently replayed with no mismatch.
% Evidence (continued from the endpoint block above, behaviour and replay):
% per_panel.{v2,p2}.behaviour.model_calls 836 + 756 (stored-row sum 1,592), valid_output_rate 1.0, extraction_rules.json_object on every call;
% .teacher_agreement_minus_chance 0.0034 / 0.0012; .last_label_rate 0.3283 / 0.2847; last-label chance about 0.13-0.14 per issue-141-closeout.md L46;
% terminations.goal_reached 3 + 1 of 24 + 22 (issue-141-closeout.md L48-49). Prediction P141 L72, held per issue-141-closeout.md L54-59.
% episode_accounting.{v2,p2}.episodes_replayed 24 / 22, replay_mismatches [] (issue-141-closeout.md L66). Rounded half-up to 3 dp from the named keys.

\begin{table}[t]
\caption{\textbf{Zero-shot base contrasts by panel.} The base is not separated from random-valid and falls below both the 0.75 rung and the trained adapter on every panel. Estimates are paired task-clustered 95\% intervals with descriptive exact sign-flip $p$ values, and pooled rows are descriptive. $\equiv$ applies the pre-registered $\pm0.05$ margin, $?$ is inconclusive, $\blacktriangledown$ separated below, and $|D_3|<0.05$ is a descriptive point-only comparison.}
\label{tab:appendix-zeroshot}
\centering\footnotesize
\setlength{\tabcolsep}{3pt}
\begin{tabular}{@{}llccc@{}}
\toprule[0.4mm]
\rowcolor{mygray}\textbf{Panel} & \textbf{Contrast} & \textbf{Estimate [95\% CI]} & \textbf{Verdict} & \textbf{$p$} \\
\midrule
Validation (12) & $D_3$ vs. random-valid & $+0.005$ $[-0.025, +0.044]$ & $\equiv$ & 0.906 \\
 & $S$ vs. $\epsilon{=}0.75$ & $-0.321$ $[-0.422, -0.226]$ & $\blacktriangledown$ & $4.88\times10^{-4}$ \\
 & $A$ vs. adapter & $-0.280$ $[-0.425, -0.142]$ & $\blacktriangledown$ & 0.004 \\
\midrule
P2 (11) & $D_3$ vs. random-valid & $+0.001$ $[-0.048, +0.051]$ & $?$ & 1.000 \\
 & $S$ vs. $\epsilon{=}0.75$ & $-0.188$ $[-0.239, -0.140]$ & $\blacktriangledown$ & $9.77\times10^{-4}$ \\
 & $A$ vs. adapter & $-0.415$ $[-0.589, -0.248]$ & $\blacktriangledown$ & 0.004 \\
\midrule
\rowcolor{yellow!25}Pooled (23) & $D_3$ vs. random-valid & $+0.003$ $[-0.024, +0.033]$ & $|D_3|<0.05$ & 0.896 \\
\rowcolor{yellow!25} & $S$ vs. $\epsilon{=}0.75$ & $-0.257$ $[-0.324, -0.197]$ & $\blacktriangledown$ & $2.38\times10^{-7}$ \\
\rowcolor{yellow!25} & $A$ vs. adapter & $-0.344$ $[-0.459, -0.232]$ & $\blacktriangledown$ & $7.63\times10^{-6}$ \\
\bottomrule
\end{tabular}
% Evidence: outputs/choice-frontier/v6/metrics/analysis.json primary.{v2,p2}.{D3,ci95,verdict,sign_flip_p}; co_primary.{v2,p2}.{S,A}.{value,ci95,verdict,sign_flip_p};
% pooled.{D3,S,A}.{D3|S|value,ci95,verdict,sign_flip.p_two_sided} (P141 L57-59, L65-67). Values rounded half-up to 3 dp from the named keys; p values below 0.001
% rounded to 3 significant figures (raw 0.00048828125, 0.0009765625, 2.384e-07, 7.629e-06).
\end{table}

\section{Results Detail: External Identity Audit}
\label{app:results-external}

% Paragraph purpose: State that ScienceWorld was chosen from a survey of published decision interfaces as a positive control for the identity audit, with LLM-First Search as a subsection.
This appendix shows how ScienceWorld was chosen from a survey of published decision interfaces, and why its audit serves as a positive control for the identity audit behind Section~\ref{sec:results-external}. The LLM-First Search audit follows as a subsection.

% Paragraph purpose: Show that the survey found no published interface of our enumeration type and that ScienceWorld was the most cited CPU-feasible candidate; give the filter, types, feasibility, and selection.
The survey found no published evaluation that uses our enumeration contract, so the external audits run on interfaces where the model picks or scores candidates. It admitted published evaluations from 2022 on in which the model chooses from a runtime-supplied candidate set inside a planning or search loop, with public tasks and code. Seven candidates passed, five of type B (the model picks one valid action per step), two of type C (the model scores candidates that a search orders), and none of type A, our enumeration contract (Table~\ref{tab:appendix-survey}). The closest published analogue, LFG \citep{shah2023lfg}, has the model score candidates in the enumeration style, but it was screened out because no evaluation harness was released. Four candidates are not CPU-feasible under their papers' budgets. For WebShop (free-text search, no step limit) and LLM-First Search (an unreported token cap) this was a judgement call. LLM-First Search was the only candidate predicted to saturate, a prediction the later audit tested (Appendix~\ref{app:results-external-lfs}). ScienceWorld is the most cited feasible candidate, and the runner-up shares its prediction.
% Evidence (LFG): docs/experiments/external-audit/issue-137-survey.md Screened out, row "Shah et al., LFG" ("the closest published case to our type A"; "the repo ships only the LLM frontier-scoring notebooks; no Habitat/HM3D harness or task list").
% Evidence: docs/experiments/external-audit/issue-137-survey.md inclusion filter, Screened out (E137-41), Candidate table (E137-40), interface
% types per row (5 B, 2 C, 0 A; E137-42); CPU feasibility and selection rule (E137-43, E137-44); judgement call quoted from
% issue-137-closeout.md Limitations and issue comment 5801181178. Type definitions paraphrase the survey rows.

% Paragraph purpose: Show that only LLM-First Search counts its budget in tokens, which a zero-token chooser never exhausts.
The surveyed interfaces differ in how they count budget, and only one counts a resource that random-valid does not spend. Six of the seven count budgets in steps, plan-length multiples, or iterations. Only LLM-First Search counts tokens, which a chooser that makes no model call never exhausts (Appendix~\ref{app:results-external-lfs}).
% Evidence: budget units per the tab:appendix-survey evidence comments (issue-137-survey.md Candidate table, budget column);
% zero-token random-valid per issue-142-protocol.md L57-62 (LLM-First Search Audit subsection).

\begin{table}[t]
\caption{\textbf{Surveyed decision interfaces and identity audits.} Both audited external interfaces register divergent choices, and random-valid matches the reference only on our runtime's enumeration contract, whose serial tie-break rule causes the identity. B selects an action, C scores search candidates, and A enumerates candidates. ``div.'' means divergent and ``sat.'' predicts saturation. $=$ means identical choices, while \cmark\ marks an observed divergence.}
\label{tab:appendix-survey}
\centering\footnotesize
\setlength{\tabcolsep}{2pt}
\begin{tabular}{@{}p{0.23\linewidth}p{0.15\linewidth}cccp{0.19\linewidth}p{0.16\linewidth}@{}}
\toprule[0.4mm]
\rowcolor{mygray}\textbf{Interface} & \textbf{Venue} & \textbf{Type} & \textbf{CPU} & \textbf{Pred.} & \textbf{Audit} & \textbf{Budget} \\
\midrule
AutoPlanBench \citep{stein2025autoplanbench} & ICAPS 2025 & B & \cmark & div. & --- & plan length $\times$ \\
ScienceWorld \citep{wang2022scienceworld} & EMNLP 2022 & B & \cmark & div. & \cmark\ 1,000/1,000 & 100 steps \\
GLAM (BabyAI-Text) & ICML 2023 & B & \cmark & div. & --- & steps (code) \\
WebShop & NeurIPS 2022 & B & \xmark & div. & --- & no step limit \\
RAP & EMNLP 2023 & C & \xmark & div. & --- & iterations \\
LLM-First Search \citep{herr2025llmfirst} & arXiv 2025 & C & \xmark$^{\dagger}$ & sat.$^{\dagger}$ & \cmark\ 400/400 & tokens \\
EmbodiedBench & ICML 2025 & B & \xmark & div. & --- & steps \\
\midrule
Our enumeration (additive) & This paper & A & --- & --- & $=$ 48/48 & calls \\
Our choice frontier & This paper & --- & --- & --- & \cmark\ 18/18 & 2$R$ decisions \\
\bottomrule
\end{tabular}
\par\vspace{2pt}\raggedright\footnotesize $^{\dagger}$Audited later with a GPU reference model; paper-budget saturation follows by construction (Appendix~\ref{app:results-external-lfs}). \cmark/\xmark\ in CPU denotes feasibility; ``---'' means no audit or no survey designation, not a negative result.
\end{table}
% Evidence: issue-137-survey.md Candidate table (venue, interface type, CPU feasible?, prediction columns) and Selection by the fixed rule
% (E137-40, E137-42, E137-43). Only LLM-First Search predicts saturated per the survey's closing statement.
% Dagger note: audited later on GPU (#142); the prediction held under the paper budget (issue-142-closeout.md L3-5).
% Evidence: ScienceWorld audit and budget: #137 outputs/external-audit/v1/scienceworld/identity-audit.json verdict CHOICE_REGISTERED, tasks 200, pairs_checked 1000,
% pairs_divergent 1000, first_divergence_index max 0, success.reference 0.915, success.random_valid_mean 0.0; envStepLimit 100 per issue-137-closeout.md.
% LLM-First Search: outputs/external-audit/v2/lfs/identity-audit.json pairs_identical 0 of pairs_checked 400; paper_budget.verdict SATURATED, paper_budget.success.random_valid_mean 1.0; supplementary.tokens_median.random_valid 0.0; matched_decision_budget.success.random_valid_mean 0.0775 (#142). Token-cap mechanism per issue-142-closeout.md L11-15.
% WebShop budget: "free-text search, no step limit" per the survey paragraph above (issue-137-survey.md, issue-137-closeout.md Limitations).
% Our enumeration contract (48/48): outputs/native-arms/v1/identity-audit.json pairs_identical 48 of pairs_checked 48; outputs/native-arms/v1/identity-audit-submission-order.json rules.sorted_candidate_order.pairs_identical 48 (#134 E134-03). Budget in decision calls per Section 4 (docs/experiments/expanded-study/synthesis-v1/README.md).
% Our choice-frontier contract (18/18): outputs/choice-frontier/v1/evaluation/identity-audit.json (pairs_divergent 18). Budget per docs/experiments/choice-frontier/issue-132-design.md sections 2 and 6
% (decision_cap = expansion_cap = 2 x reference_expansions). Type A is the survey's label for our enumeration contract; the choice-frontier contract was not typed by the survey (n/a).
% Evidence (budget units): docs/experiments/external-audit/issue-137-survey.md Candidate table, budget column.
% AutoPlanBench: "twice the length of the plan generated by ... GBFS with the FF heuristic" (Sec. 5). GLAM: max_steps 64 or 128 in the released code.
% RAP: "10 or 20 MCTS iterations, depth limit = task length". EmbodiedBench: "The maximum number of environment steps is 30 for high-level tasks,
% 20 for EB-Navigation, and 15 for EB-Manipulation" (Sec. 5.1).

% Paragraph purpose: Show that success on ScienceWorld measures choice because the reference usually succeeds and random-valid never does; give the frozen protocol and results.
On ScienceWorld success measures choice, because the reference usually succeeds and random-valid never does. The frozen protocol takes the first 200 test variations under the paper's 100-step limit with five random-valid seeds (17, 5077, 6131, 7409, and 8527), and ten re-run episodes are byte-identical. The reference completes 183 of 200 (0.915), including 3/9 boil, 6/9 melt, 4/9 freeze, 6/9 change-state, 135/135 use-thermometer, and 29/29 measure-melting-point variations. It falls short only where 17 of the 36 Matter variations need more than 100 moves. Random-valid completes none, and 997 of its 1000 episodes end in a failure score.
% Evidence: docs/experiments/external-audit/issue-137-protocol.md (frozen at 914339d); outputs/external-audit/v1/scienceworld/identity-audit.json
% tasks 200, pairs_checked 1000, success.random_valid_per_seed seeds (E137-02, E137-03, E137-07); envStepLimit 100 per issue-137-closeout.md;
% outputs/external-audit/v1/scienceworld/determinism.json all_identical true, episodes 10 (E137-18). CPU only (LED-4).
% Evidence: identity-audit.json pairs_divergent 1000, first_divergence_index max 0, verdict CHOICE_REGISTERED, success.reference 0.915,
% success.random_valid_mean 0.0 and per seed 0.0 (E137-01, E137-03..07); 183 of 200 per issue-137-closeout.md; per-task counts (E137-16);
% episode fates 997/1000 and 3 step limit (E137-17); 17 of 36 Matter tasks exceed 100 moves (E137-20).

% Paragraph purpose: Show that the audit covers a narrow slice of ScienceWorld on a newer simulator version; state the scope, score, and version caveats.
The audit covers a narrow slice of ScienceWorld on a newer simulator version. The slice holds 6 of 30 task types (200 of 1,819 test variations), on all of which the paper's own Random-Valid score is 0.00. The paper reports normalized scores, so its DRRN score of 0.084 on these tasks is not a success rate. The audit uses ScienceWorld 1.3.0, whose notes warn that gold paths can change relative to 1.2.x, and the paper's simulator version was not re-run.
% Evidence: issue-137-closeout.md Limitations (6 of 30, 200 of 1,819, Random-Valid 0.00, task 4-2 outside the slice; E137-19, sheet #137
% judgement calls 2 and 3); identity-audit.json published_model_success.value 0.084, random_fraction_of_published 0.0 (E137-09, E137-10);
% supplementary.published_random_valid_score_same_tasks 0.0 (E137-13); upstream pin ScienceWorld 1.3.0 (E137-15).

\subsection{LLM-First Search Audit}
\label{app:results-external-lfs}

% Paragraph purpose: State that this subsection shows why a token budget cannot bind a chooser that spends no tokens, using the LLM-First Search audit behind the budget check.
A token budget cannot bind a chooser that spends no tokens. The LLM-First Search audit shows the consequence and why step~2 of Box~\ref{box:procedure} counts decisions instead.

% Paragraph purpose: Show that LLM-First Search keeps every unexpanded alternative, which makes a zero-token random chooser a complete search; define random-valid and the decision unit.
LLM-First Search keeps every unexpanded alternative, and this makes a zero-token random chooser a complete search. It is a type-C interface. The model scores candidates, and the search expands the argmax child or pops the highest-valued frontier node, so every unexpanded alternative stays in the frontier. Random-valid answers every query with independent uniform draws, spends zero tokens, and makes no model call. It uses the five seeds of the first audit, and a decision is one node expansion after the root.
% Evidence: issue-142-protocol.md L11-13 (type C), L57-62 (random-valid, zero tokens), L66-68 (decision); seeds 17/5077/6131/7409/8527 (protocol L61) (copied from the block below).

% Paragraph purpose: Show that the audit runs the upstream code unchanged under the paper's budgets, and declare the tasks and the four deviations.
The audit runs the upstream code unchanged under the paper's own budgets and declares four deviations. The upstream code (commit \texttt{3025bdaa}, Apache-2.0) is imported unchanged and not vendored. The fixed rule audits every configuration whose published GPT-4o win rate exceeds 10\%, which excludes Sudoku 6$\times$6 and keeps Countdown 3, 5, and 7 and Sudoku 4$\times$4, with games 0 to 19 of each, 80 tasks in all. Budgets are the paper's per-game token caps (1,000,000 for Countdown and 100,000 for Sudoku 4$\times$4). The deviations are these. The paper's GPT-4o is replaced by the local Qwen3-30B-A3B-Instruct-2507, so reference success is that model's. The reference runs once per task. Sudoku 6$\times$6 is excluded. Random-valid uses a faster frontier implementation that produced byte-identical logs on 57/57 probe episodes.
% Evidence: issue-142-protocol.md L11-13 (type C), L57-62 (random-valid, zero tokens), L66-68 (decision), L14-16 (upstream, not vendored), L24-37 (rule, games 0-19),
% L31-34 (caps), L191-200 (deviations 1-4); identity-audit.json interface_type C, upstream.commit 3025bdaa3add..., upstream.license Apache-2.0, tasks 80,
% per_config.{cd3,cd5,cd7,su4}.tasks 20; lfs-probe/fast-equivalence.json checked 57, identical 57; seeds 17/5077/6131/7409/8527 (protocol L61).

% Paragraph purpose: Show that random-valid diverges from the reference early in every pair, with replay receipts.
Random-valid diverges from the reference in every pair, and early. No pair of the 400 is identical, and the first divergence has median index 0 and maximum 2. All 80/80 reference and 400/400 random-valid episodes were independently replayed with no mismatch.
% Evidence: identity-audit.json replay.reference {80, 80, 0, 0}, replay.random_valid {400, 400, 0, 0} (lfs/replay.json); pairs_identical 0 of pairs_checked 400;
% first_divergence_index.median 0, .max 2;

% Paragraph purpose: Show that under the paper's token budget random-valid wins every episode, and that this saturation follows by construction.
Under the paper's token budget random-valid nevertheless wins every episode, and this saturation follows by construction. Random-valid spends no tokens and the frontier keeps every alternative, so it is a complete randomized search, and a pre-audit probe had won 100/100 non-audited episodes. Its only way to lose was a 500,000-expansion safety guard, which never bound. Random-valid used a median of 201 expansions and a maximum of 102,606, against a median of 29.5 and a maximum of 1,300 for the reference. Random-valid wins at every seed, so the verdict is SATURATED, and every reference loss ends at the token cap (Table~\ref{tab:appendix-lfs}).
% Evidence: issue-142-protocol.md L45 (500,000 guard), issue-142-closeout.md L30, L74 (never bound); probe 100/100 per issue-142-protocol.md L166-167;
% identity-audit.json supplementary.expansions_median {random_valid 201.0, reference 29.5}, supplementary.expansions_max {random_valid 102606, reference 1300};
% paper_budget.verdict SATURATED, .success.random_valid_per_seed all 1.0; supplementary.terminated.reference {success 66,
% token_limit 14};

% Paragraph purpose: Show that a budget counted in decisions lets success register choice, with the matched-budget result, the predictions, published success, and the ledger.
A budget counted in decisions lets success register choice again. Within the reference's own number of expansions per task, random-valid wins 0.0775 (per seed 0.0625, 0.075, 0.075, 0.0875, and 0.0875), so the verdict is CHOICE\_REGISTERED. Doubling that budget raises random-valid's success only to 0.19. The reference's matched-budget success equals its token-budget success by construction, so the gap between the two arms in Table~\ref{tab:appendix-lfs} is the evidence that choices matter. The matched-budget prediction held, and the paper-budget saturation follows by construction. Published GPT-4o success on these configurations is 0.768. One amendment, made before any audit episode, only reordered the smoke and audit start, and the ledger spent 2.89 of 12 GPU-h with no paid API use.
% Evidence (0.19 at twice the reference's expansions, copied from the caption-curve comment below): matched_decision_budget.random_valid_decision_budget_curve.random_valid_success [0.0775, 0.0975, 0.14, 0.17, 0.19] at multipliers [1.0, 1.25, 1.5, 1.75, 2.0].
% Evidence: matched_decision_budget.random_valid_decision_budget_curve.reference_success_at_own_budget 0.825.
% Evidence: matched_decision_budget.success.random_valid_mean 0.0775, .random_valid_per_seed {0.0625, 0.075, 0.075, 0.0875, 0.0875}, .verdict CHOICE_REGISTERED;
% prediction {paper_budget SATURATED, matched CHOICE_REGISTERED}, held per issue-142-closeout.md L3-5; published_model_success.value 0.7684 (0.768);
% Amendment A1 issue-142-protocol.md L212-223 (commit e89a793); budget.json total_gpu_hours 2.8910, gpu_hours_cap 12.0, paid_api_usd 0.0.

\begin{table}[t]
\caption{\textbf{Descriptive LLM-First Search budget audit.} Random-valid saturates under the paper's token budget but falls far below the reference under a matched decision budget. Differences are substituted local reference minus random-valid success with task-cluster 95\% intervals. The matched-budget reference success equals its paper-budget success by construction, and \cmark\ denotes registered choice and $=$ saturation.}
\label{tab:appendix-lfs}
\centering\footnotesize
\setlength{\tabcolsep}{4pt}
\begin{tabular}{@{}lccc@{}}
\toprule[0.4mm]
\rowcolor{mygray}\textbf{Configuration} & \textbf{Reference} & \textbf{Random-valid} & \textbf{Difference [95\% CI]} \\
\midrule
\multicolumn{4}{@{}l}{\textit{Paper token budget ($=$, saturated)}} \\
Countdown 3 (20) & 1.000 & 1.000 & $0.000$ $[0.000, 0.000]$ \\
Countdown 5 (20) & 0.750 & 1.000 & $-0.250$ $[-0.450, -0.100]$ \\
Countdown 7 (20) & 0.650 & 1.000 & $-0.350$ $[-0.550, -0.150]$ \\
Sudoku 4$\times$4 (20) & 0.900 & 1.000 & $-0.100$ $[-0.250, 0.000]$ \\
All (80) & 0.825 & 1.000 & $-0.175$ $[-0.263, -0.100]$ \\
\midrule
\multicolumn{4}{@{}l}{\textit{Matched decision budget (\cmark, choice registered)}} \\
Countdown 3 (20) & 1.000 & 0.120 & $+0.880$ $[+0.770, +0.970]$ \\
Countdown 5 (20) & 0.750 & 0.110 & $+0.640$ $[+0.370, +0.870]$ \\
Countdown 7 (20) & 0.650 & 0.020 & $+0.630$ $[+0.410, +0.840]$ \\
Sudoku 4$\times$4 (20) & 0.900 & 0.060 & $+0.840$ $[+0.700, +0.950]$ \\
All (80) & 0.825 & 0.0775 & $+0.747$ $[+0.648, +0.838]$ \\
\bottomrule
\end{tabular}
\end{table}
% Evidence: outputs/external-audit/v2/lfs/identity-audit.json per_config.{cd3,cd5,cd7,su4}.{paper_budget,matched_decision_budget}.{success.reference,success.random_valid_mean,
% contrast.mean_difference_reference_minus_random,contrast.bootstrap_95ci,verdict}; top-level paper_budget.* and matched_decision_budget.* (-0.175 [-0.2625, -0.1];
% raw 0.7474999999999999 [0.6475, 0.8375]; rendered 0.747 from the stored raw value). Rounded half-up from the stored repr; 0.0775 kept verbatim. Descriptive statistics per issue-142-protocol.md L144-156.
% Caption curve: matched_decision_budget.random_valid_decision_budget_curve.{multipliers [1.0,1.25,1.5,1.75,2.0], random_valid_success [0.0775,0.0975,0.14,0.17,0.19], auc_normalized 0.1353125 (0.135)}.

\section{Results Detail: Seed Replication and Follow-Up Limits}
\label{app:results-boundary}

% Paragraph purpose: State that the enumeration-contract seed replication reproduces the seed-17 verdicts, and that this appendix lists every follow-up window's development-stage limits.
This appendix shows that the enumeration-contract seed replication reproduces the seed-17 verdicts, and it lists the development-stage limits of every follow-up window.

% Paragraph purpose: Show that both replicated enumeration-contract cells give the same verdict at all three training seeds.
Both replicated enumeration-contract cells give the same verdict at all three training seeds. The prospective replication froze seeds 17, 29, and 71 before launch, independently replayed 534 new episodes, and executed on two cells (Table~\ref{tab:appendix-seeds}). In the greedy $\times$ multimodal replication cell (greedy best-first $\times$ multimodal-state, 24-problem unseen panel), per-seed success is 21/24, 21/24, and 22/24 (mean 0.889, range $[0.875, 0.917]$). The contrast against the pretrained base is positive at every seed ($+0.875$/$+0.875$/$+0.917$, no sign flip), and the contrast against random-valid stays negative ($-0.125$/$-0.125$/$-0.083$). On the DAgger iteration-1 cell set (BFS $\times$ three modalities, DAgger and continued-SFT arms), every arm and seed scores 0/72 unseen, with paired contrast $+0.000$ $[+0.000, +0.000]$ at each seed, and the published iteration-2 verdict (1/72 vs 1/72, seed 17) is unchanged.
% Evidence: see the comment below (seed-replication analysis.json).

% Paragraph purpose: Show that the replication covers only part of the enumeration matrix, so cross-seed aggregation stays descriptive.
The replication covers only part of the enumeration matrix. Cell A covers one of the three modalities of the greedy additive cell. The remaining two modalities and all other enumeration-contract cells use seed 17 only, and cross-seed aggregation is descriptive ($n{=}3$). The image-only-arm proposal has since executed under the native-arms v2 amendment (Appendix~\ref{app:results-native}), and its v1 gate remains as history.
% Evidence: outputs/expanded-study/v1/seed-replication/analysis.json baseline_process_sft.aggregated and dagger_iteration_1.aggregated.
% #127 closed 2026-09-21. #128 TODO retired: superseded by #129 (executed) and #130 (stress tests); see issue-129-closeout.md.


% Paragraph purpose: Show that every follow-up window is development-stage and none touched the frozen 45-task final evaluation; list each window's limits.
Every follow-up window is development-stage, and none accessed the frozen 45-task final evaluation. The native-arms window evaluates an 18-unit clean panel with one training seed per cell. The R1-R4 follow-up window reuses the frozen native-arms adapters for R1, R2, and R4, and trains one seed-17 run per cell for R3. The choice-frontier window trains one seed per algorithm on a 9-task panel whose discriminative subset holds one task. The comparator zoo is a CPU-only exploratory analysis. The adapter has three training seeds on the validation panel, the held-out panel, and the unscreened panel, and the DAgger ablation trains one seed on half the planned corpus. The shared nine-task panel draws from 8 domains with 7 compact and 2 expanded variants.
% Evidence: A139 training_seeds {p135, p2, p2u} [17, 29, 71]; A140 primary.seeds [17]; held-out manifest never read (issue-139-closeout.md L11).
% Evidence: issue-129-closeout.md, issue-130-closeout.md, issue-132-closeout.md (Erratum item 3), #133 closing comment; held-out manifest
% status; panel composition per outputs/choice-frontier/o4/metrics/optimal-costs.json.

% Paragraph purpose: State the undecided prerequisites that must be recorded before the held-out final primary evaluation can be frozen.
\paragraph{Final primary evaluation prerequisites.}
No authorized protocol yet exists for the frozen 45-task final evaluation. Freezing one requires recorded decisions on the following prerequisites. (1) The condition-level data source and Search-Trace Segment construction with per-algorithm allocation. (2) Compute parity and checkpoint-selection rules. (3) The episode-budget unit, the whole-instance split procedure, and the fixed Search Memory capacity. (4) The estimand, the uncertainty interval or test, the success threshold, the invariant-compliance rule, and the validity and invalid-operation denominators for every cell. Until these choices are recorded, the final evaluation protocol cannot be frozen and no final comparative efficacy finding exists.
% Held-out final primary evaluation: no open ticket; stated in prose with no visible marker per the round-2 TODO convention.
% The preregistration prerequisites in writing_design_tree.md are undecided.

\section{Decision Contracts}
\label{app:design-contracts}

% Paragraph purpose: State that the two decision contracts differ in what the policy emits and what the runtime owns, so their results are not comparable; point to the table.
The two decision contracts differ in what the policy emits and what the runtime owns, so their results are not comparable. Table~\ref{tab:design-contracts} sets them side by side, and Section~\ref{sec:policy} defines both.

\begin{table}[t]
\caption{\textbf{Two non-comparable decision contracts.} The policy emits and the runtime checks different units under the enumeration and choice-frontier interfaces. The identity-audit row records the executed scope, not a cross-contract performance comparison.}
\label{tab:design-contracts}
\centering\footnotesize
\setlength{\tabcolsep}{3pt}
\begin{tabular}{@{}p{0.17\linewidth}p{0.39\linewidth}p{0.39\linewidth}@{}}
\toprule[0.4mm]
\rowcolor{mygray}\textbf{Attribute} & \textbf{Enumeration} & \textbf{Choice frontier} \\
\midrule
Policy output & Grounded, unscored Typed Search Operation & Opaque frontier label \\
Frontier owner & Search Memory; visited / score & Runtime; seeded per-decision labels \\
Observation & Modality adapter & Visual: task / initial / 128px frontier / goal \\
Learned budget & $2\times$ reference decisions; base: 1 call & $2\times$ reference expansions and decisions; goal free \\
Identity audit & Diagnostic, additive cells & Pre-registered, 18 pairs \\
\bottomrule
\end{tabular}
\end{table}
% Evidence: docs/experiments/choice-frontier/issue-132-design.md sections 2 and 6 (menu permutation master seed 51131,
% visual-only payload, 18 (task, algorithm) control pairs); outputs/native-arms/v1/identity-audit.json (enumeration-contract audit).

\section{Use of LLM Statement}
\label{app:llm-usage}

% Paragraph purpose: Disclose how large language models were used in preparing this paper.
Large language model agents drafted and revised the manuscript text and wrote the figure scripts under the authors' direction. An agent pass searched for citations, and a human verified every cited source. Agent reviewers ran repeated review rounds on the drafts. LLM agents also wrote and ran experiment and analysis code, drafted protocols and two pre-evaluation protocol amendments on the authors' behalf, and compiled evidence sheets, and the authors approved each protocol before its outcome. The authors verified every reported number against its pinned artifact and take full responsibility for the content.
